% PROVENANCE: This manuscript was written entirely by AI.
% Models that drafted or edited it:
%   Claude Opus 5.5 (Anthropic), in Claude Code, outline, proof notes and
%     drafting, October 2026
%   GPT-6 Astra (OpenAI), in Codex, verification of proofs, computations and
%     review, October 2026
%   GPT-6 Astra (OpenAI), in Codex (effort ultra), block-by-block review (W5)
%     with minor edits, 2026-10-08
%   GPT-6 Astra (OpenAI), in Codex (effort high), application of the prose edits
%     decided by Claude (P-apply), 2026-10-08
%   Claude Sonnet 5.5 (Anthropic), in Claude Code, round 2: wording changes,
%     drafts of Appendix B paragraphs, 2026-10-09 (reviewed by Claude Opus 5.5)
%   Claude Opus 5.5 (Anthropic), in Claude Code, round 2: second checks,
%     corrections and edits of Sections 3 and 6 to 10 and Appendix B, 2026-10-09
% Detailed history: git log; log/CONVERSATION.md and PROGRESS.md of the
% repository; Appendix D.
%
% Research report on the OpenAI preprints on the finitistic dimension conjecture
% and the Auslander--Reiten conjecture (September 2026). Sources: this file
% (preamble and outline) and sections/*.tex, included with \input.
% Build with latexmk (see latexmkrc): pdfLaTeX and biber, bibliography in
% library.bib.

\documentclass[a4paper]{amsart}

% Allow line breaking at long run-in headings and repository URLs.
\setlength{\emergencystretch}{1em}
\hyphenation{OpenAI}

% ============================================================
% Core mathematics
% ============================================================

\usepackage{mathtools}
\usepackage{amssymb}
\usepackage{eucal}

% ============================================================
% Color, graphics, and diagrams
% ============================================================

\usepackage{xcolor}
\usepackage{eso-pic}
\usepackage{graphicx}
\usepackage{tikz-cd}
\usepackage{array}

% ============================================================
% Theorem infrastructure
% ============================================================

\usepackage{thmtools}

% ============================================================
% Bibliography
% ============================================================

\usepackage[
  backend=biber,
  style=alphabetic,
  isbn=false,
  doi=false,
  url=false,
  eprint=true,
  giveninits=true,
  backref=false,
  maxbibnames=99
]{biblatex}

\renewbibmacro{in:}{}
% The long label [StacksProj] widens the hanging indent; allow looser lines.
\appto{\bibsetup}{\sloppy}
\addbibresource{library.bib}

% ============================================================
% Hyperlinks and cross-references
% ============================================================

\usepackage{hyperref}

\hypersetup{
  colorlinks,
  linkcolor={red!80!black},
  citecolor={blue!80!black},
  urlcolor={blue!80!black}
}

% Last package, and before the theorem declarations (see below).
\usepackage{cleveref}

% ============================================================
% Core notation
% ============================================================

\newcommand{\kk}{\mathbf{k}}
\newcommand{\ZZ}{\mathbb{Z}}
\newcommand{\op}{\mathrm{op}}

% \id       -> id
% \id[A]    -> id_A
\NewDocumentCommand{\id}{s o}{%
  \IfBooleanTF{#1}{1}{\operatorname{id}}%
  \IfValueT{#2}{_{#2}}%
}

% \set{x}, \set{x}[condition]; starred form scales delimiters.
\NewDocumentCommand{\set}{s m o}{%
  \IfBooleanTF{#1}
    {\left\{#2\IfValueT{#3}{\,\middle|\,#3}\right\}}
    {\{#2\IfValueT{#3}{\,\mid\,#3}\}}%
}

% Hom and derived Hom; optional square argument is the base/category.
\NewDocumentCommand{\Hom}{o m m}{%
  \operatorname{Hom}\IfValueT{#1}{_{#1}}(#2,#3)%
}
\NewDocumentCommand{\RHom}{o m m}{%
  \mathbb{R}\!\operatorname{Hom}\IfValueT{#1}{_{#1}}(#2,#3)%
}

% Derived tensor product. Use ordinary \otimes for the underived tensor.
\NewDocumentCommand{\Lotimes}{o}{%
  \otimes^{\mathbb{L}}\IfValueT{#1}{_{#1}}%
}

% \DerCat{\A}, \DerCat[b]{\A}
\NewDocumentCommand{\DerCat}{o d<> m}{%
  \operatorname{D}%
  \IfValueT{#1}{^{\mathrm{#1}}}%
  \IfValueT{#2}{_{#2}}%
  (#3)%
}

\DeclareMathOperator{\img}{im}

% Define only category letters needed by the article.
\newcommand{\category}[1]{\mathcal{#1}}
\newcommand{\A}{\category{A}}

% ============================================================
% Paper-specific notation
% ============================================================

% Notice printed in the top margin of every page (requested by Gustavo Jasso).
\AddToShipoutPictureFG{%
  \AtPageUpperLeft{%
    \raisebox{-1.1cm}{\makebox[\paperwidth]{%
      \textcolor{red}{\normalfont\bfseries NOT VERIFIED -- PROVIDED AS IS}}}}}

% ============================================================
% Theorem environments
% ============================================================

\numberwithin{equation}{section}
% The shared counter needs section-qualified PDF destinations.
\renewcommand{\theHequation}{\theHsection.\arabic{equation}}

\declaretheorem[style=plain,sibling=equation]{theorem}
\declaretheorem[style=plain,sibling=equation]{proposition}
\declaretheorem[style=plain,sibling=equation]{lemma}
\declaretheorem[style=plain,sibling=equation]{corollary}
\declaretheorem[style=plain,numbered=no,name=Theorem]{theorem*}

\declaretheorem[style=definition,sibling=equation]{definition}
\declaretheorem[
  style=definition,
  sibling=equation,
  Refname={Notation,Notations}
]{notation}
\declaretheorem[style=definition,sibling=equation]{construction}

\declaretheorem[style=remark,sibling=equation]{remark}
\declaretheorem[style=remark,sibling=equation]{example}

\declaretheorem[
  style=definition,
  numbered=no
]{conventions}

% ============================================================
% Metadata
% ============================================================

\title[OpenAI's counterexamples to the finitistic dimension conjecture]{Report on OpenAI's counterexamples to the little finitistic
  dimension conjecture for finite-dimensional algebras}

\author[AI agents (Claude Opus 5.5 and GPT-6 Astra)]{Prepared autonomously by AI
  agents (Claude Opus~5.5 and GPT-6 Astra)}
\address[G.~Jasso]{Mathematisches Institut, Universität zu Köln,
  Weyertal 86--90, 50931 Köln, Germany}
\email{gjasso@math.uni-koeln.de}
\urladdr{https://gustavo.jasso.info}

\begin{document}

\begin{abstract}
We reconstruct, with proofs written out in full (not verified by a human),
the two constructions of
finite-dimensional algebras of infinite little finitistic dimension released by
OpenAI in September 2026: one realises a selection functor on modules over a
finitely presented algebra by a bimodule complex and a trivial extension; the
other obtains a counterexample to the Auslander--Reiten conjecture from a
bimodule and a morphism in the stable module category of a symmetric algebra. We
also prove criteria for infinite finitistic dimension and
obstructions to simpler constructions.
\end{abstract}

\maketitle

\setcounter{tocdepth}{1}
\tableofcontents

\section*{Caveat lector}

This report was produced as an aid for verifying and understanding the
counterexamples to the little finitistic dimension conjecture in the preprints
of OpenAI~\cite{OAI26findim,OAI26ar}. It was orchestrated by Gustavo Jasso and
prepared by AI agents, as described in \Cref{app:ai-declaration}. It has not
been verified by a human, and it is provided as is.

\section{Introduction}
\label{sec:introduction}

The little finitistic dimension of a finite-dimensional algebra $A$ is the
supremum of the projective dimensions of the finitely generated $A$-modules of
finite projective dimension. Bass asked whether it is always
finite~\cite{Bas60}; see~\cite{ZHui95} for the history of the resulting
finitistic dimension conjecture. In September 2026, OpenAI released preprints
that claim counterexamples. The first~\cite{OAI26findim} constructs a
finite-dimensional complex algebra of infinite little finitistic dimension. The
second~\cite{OAI26ar} constructs a counterexample to the Auslander--Reiten
conjecture over the field $\mathbb{F}_2(q,H_1,H_2)$; by a theorem of Auslander
and Reiten~\cite[Theorem~1.1(b)]{AR75a}, together with
\Cref{thm:strong-nakayama}, such a counterexample yields a further
algebra of infinite little finitistic dimension.

This article reconstructs both constructions with proofs written out in full
(not verified by a human), organised
around general statements, some more general than the applications require. It
also gives criteria and obstructions that explain the choices in the
constructions. All proofs were written and checked by AI models, as
described in \Cref{app:verification,app:ai-declaration}; no human has checked
them, and the article has not been peer reviewed. The repository that
accompanies the preprints poses the existence statements of both constructions
as challenges in the proof assistant Lean and provides formal proofs of them;
\Cref{app:verification} compares these statements with the present article and
records what has been checked.

\subsection*{The main construction}

A \emph{selection functor}
${H=\Psi\otimes_R-}$ on finite-dimensional modules over an algebra $R$, possibly
infinite-dimensional, annihilates modules at arbitrarily late times
(\Cref{def:selection-data}). A bounded complex $P$ of bimodules over a
finite-dimensional algebra $B$ of global dimension at most two realises $H$, in
the
sense that ${P\Lotimes[B]M(Y)\cong M(HY)\oplus M(HY)[3]}$ for a fully faithful
embedding $M$ (\Cref{thm:realisation}). The complex $P$ is simulated by an
ordinary bimodule $X$ over an algebra $\Delta$ of finite global dimension, and
the trivial extension ${\Delta\ltimes X}$ turns the extinction times of $H$ into
projective dimensions (\Cref{prop:simulation,coro:extinction}). The result is
the following theorem, proved in \Cref{sec:simulation}; the main preprint
treats the case of a central idempotent and an endomorphism of~$R$.

\begin{theorem*}[\Cref{thm:main}]
  Let $(R,\Psi)$ be selection data over a field, with $R$ finitely presented
  and with unbounded extinction. Then there is a finite-dimensional algebra
  ${A=\Delta\ltimes X}$, with $\Delta$ of finite global dimension, such that
  ${\operatorname{findim}A=\infty}$.
\end{theorem*}

Selection data with unbounded extinction exist: the complex group algebra of a
finitely presented group with a family of independent central involutions,
shifted by an automorphism, provides them (\Cref{thm:selection-unbounded}),
which gives the counterexample of the main preprint (\Cref{coro:main}). The
selection data must be infinite-dimensional in an essential way: their rank
functions on finite-dimensional modules must span an infinite-dimensional space
(\Cref{prop:rank-obstruction}), which excludes finite-dimensional algebras,
commutative noetherian algebras and universal localisations of
finite-dimensional hereditary algebras
(\Cref{coro:rank-obstruction}). The realisation is the only step that is not
explicit: it lifts a diagram from a Verdier quotient by the calculus of
fractions, and the argument gives no bound on the length $l$ of the complex $P$,
hence on the number ${3(l+2)}$ of simple modules of the final algebra
(\Cref{rem:uncontrolled-length,rem:what-is-explicit}). The argument works with
an object ${U\oplus U[3]}$ instead of the summand $U$ that encodes $\Psi$: the
summand $U$ need not lie in the Verdier quotient, while ${U\oplus U[3]}$ does
(\Cref{prop:odd-double}), and with the shift $3$ the obstruction to splitting
the realisation lies in an $\operatorname{Ext}^4$ over an algebra of global
dimension at most two (\Cref{rem:why-three}).

\subsection*{The counterexample to the Auslander--Reiten conjecture}

The second construction uses the following result over an arbitrary field
(\Cref{thm:conversion}): for a symmetric algebra $E$, a
bimodule $F$ projective on both sides, a module $S$ and a stable morphism
${v\colon S\to FS}$, if the maps ${(g,j)\mapsto F(g)v-v[a]j}$ between Tate
cohomology groups are bijective in positive degrees and surjective with nonzero
kernel in degree zero, then a module over the triangular algebra
$\left(\begin{smallmatrix}E&0\\F&E\end{smallmatrix}\right)$ is a counterexample
to the Auslander--Reiten conjecture. The construction uses a ten-dimensional
algebra $C$, whose simple module has an aperiodic resolution, the symmetric
algebra ${T=C\ltimes DC}$, over which this simple module has polynomial Ext
algebra on a generator of degree three, and the tensor square ${E=T\otimes T}$,
over which two cones and two twists produce the bimodule $F$
(\Cref{sec:ar-route}). With \Cref{thm:ar-to-findim} this gives an algebra of
infinite little finitistic dimension with nine simple modules
(\Cref{coro:ar-findim}); the triangular algebra used to obtain it has eight
simple modules and dimension of the order of $10^{35}$ (\Cref{rem:ar-sizes}).

\subsection*{Criteria and obstructions}

The work that led to this article began with a search for a counterexample that
could be verified by hand. The search did not succeed, and its results are the
general statements of \Cref{sec:criteria,sec:extinction,sec:obstructions}.
Failure of the strong Nakayama conjecture produces explicit modules of every
finite projective dimension (\Cref{thm:strong-nakayama}); simple modules
witnessing this failure always come from a corner of the algebra and a
counterexample to the Auslander--Reiten conjecture over that corner
(\Cref{prop:simple-witness});
and finite projective dimensions, as well as finite extinction times, are
bounded on modules of bounded dimension
(\Cref{prop:bounded-dimension,prop:bounded-extinction}), where the assertion
about extinction times is for trivial extensions of algebras of finite global
dimension. For the construction of the counterexample to the Auslander--Reiten
conjecture, the analogous construction over $T$
itself, with one cone instead of two, fails in degree zero
whenever the simple module has polynomial Ext algebra on a generator of degree
three (\Cref{prop:one-factor,coro:one-factor}); the proof uses Tate duality,
which also forces the vanishing of the corresponding Toda brackets
(\Cref{thm:tate-obstruction}). Whether a counterexample with a small explicit
presentation exists remains open.

\subsection*{Structure}

\Cref{sec:preliminaries} fixes conventions and recalls Tate duality and Toda
brackets. \Cref{sec:criteria} contains the general criteria. The main
construction occupies \Cref{sec:extinction,sec:selection,sec:realisation,%
sec:simulation}: trivial extensions and extinction, selection data and the
group, the realisation by a bimodule complex, and the simulation with the main
theorem. The counterexample to the Auslander--Reiten conjecture is constructed
in \Cref{sec:conversion,sec:ar-route},
and \Cref{sec:obstructions} discusses obstructions to smaller constructions.
\Cref{app:computations} describes the computations used,
\Cref{app:verification} the verification and formalisation record,
\Cref{app:companion} the third preprint of the release, and
\Cref{app:ai-declaration} the use of AI in the preparation of this article.

\section{Conventions and preliminaries}
\label{sec:preliminaries}

\begin{conventions}
  Throughout, $\kk$ is a field and ${D=\Hom[\kk]{-}{\kk}}$ denotes the
  $\kk$-linear duality. Algebras are associative and unital; unless stated
  otherwise, they are finite-dimensional over $\kk$, although the algebras of
  selection data in \Cref{sec:selection} need not be. Modules are finitely
  generated \emph{left} modules unless they are called right modules, and right
  $A$-modules are identified with left $A^{\op}$-modules. Composition of maps
  and products of paths in a quiver are written from right to left, so that an
  arrow ${a\colon i\to j}$ lies in $e_jAe_i$. Complexes are graded
  cohomologically, with ${K[s]^n=K^{n+s}}$ and ${d_{K[s]}=(-1)^sd_K}$. For a module
  $M$ we write ${M^*=\Hom[A]{M}{A}}$, which is a right module, $\Omega M$ for the
  kernel of a projective cover of $M$, and $\operatorname{Tr}M$ for the
  transpose computed from a minimal projective presentation: if
  ${P_1\to P_0\to M\to 0}$ is minimal, then
  ${\operatorname{Tr}M=\operatorname{coker}(P_0^*\to P_1^*)}$.
\end{conventions}

The \emph{little finitistic dimension} of an algebra $A$ is
\[
  \operatorname{findim}A\coloneqq\sup\set*{\operatorname{pd}_AM}[\substack{
    M\text{ a finitely generated left }A\text{-module}\\
    \operatorname{pd}_AM<\infty}].
\]
The \emph{little finitistic dimension conjecture} asserts that
${\operatorname{findim}A<\infty}$ for every finite-dimensional
algebra~$A$~\cite{Bas60}; see~\cite{ZHui95} for its history. The finitistic
dimension of $A^{\op}$ is the \emph{right} finitistic dimension of $A$; the two
need not agree.

\subsection{Stable categories of symmetric algebras}
\label{subsec:stable}

Let $A$ be a symmetric algebra. Its stable module category
$\operatorname{\underline{mod}}A$ is triangulated, with suspension $[1]$ given
by the cosyzygy $\Omega^{-1}$, so that ${[-1]=\Omega}$, and with triangles
induced by short exact sequences. For modules $X$ and $Y$ and an integer $a$, we
write
\[
  \widehat{\operatorname{Ext}}{}^a_A(X,Y)\coloneqq
  \underline{\operatorname{Hom}}_A(X,Y[a]),
\]
so that ${\widehat{\operatorname{Ext}}{}^a_A(X,Y)=\operatorname{Ext}^a_A(X,Y)}$
for ${a\geq1}$. Composition makes $\widehat{\operatorname{Ext}}{}^*_A(X,X)$ a
graded algebra.

\begin{theorem}[{Tate duality; \cite[Section~2, (2.1)]{Lin13}}]
  \label{thm:tate-duality}
  Let $A$ be a symmetric algebra. For all modules $X$, $Y$ and every integer
  $a$, there is an isomorphism
  \[
    D\widehat{\operatorname{Ext}}{}^a_A(X,Y)\cong
    \widehat{\operatorname{Ext}}{}^{-1-a}_A(Y,X),
  \]
  natural in $X$ and $Y$.
\end{theorem}

Equivalently~\cite[Section~2, (2.2)]{Lin13}, there is a natural nondegenerate
bilinear form
\[
  \langle-,-\rangle\colon\widehat{\operatorname{Ext}}{}^{a}_A(Y,X)\times
  \widehat{\operatorname{Ext}}{}^{-1-a}_A(X,Y)\longrightarrow\kk,
\]
and it is compatible with composition: ${\langle\zeta\eta,\tau\rangle=
\langle\zeta,\eta\tau\rangle}$ whenever the composites are
defined~\cite[Section~2, (2.8)]{Lin13}.

\subsection{Toda brackets}
\label{subsec:toda}

Let $\mathcal{T}$ be a triangulated category and let
\[
  X\xrightarrow{f}Y\xrightarrow{g}Z\xrightarrow{h}W
\]
be morphisms with ${gf=0}$ and ${hg=0}$. Choose a triangle
\[
  Y\xrightarrow{g}Z\xrightarrow{i}C_g\xrightarrow{q}Y[1].
\]
A \emph{defining
system} is a pair of morphisms ${a\colon X[1]\to C_g}$ and ${b\colon C_g\to W}$ with
${qa=f[1]}$ and ${bi=h}$; such pairs exist since ${gf=0}$ and ${hg=0}$. The
\emph{Toda bracket} ${\langle h,g,f\rangle\subseteq\Hom[\mathcal{T}]{X[1]}{W}}$
is the set of the composites $ba$ over all defining systems. It is a coset of
the \emph{indeterminacy}
\[
  h\circ\Hom[\mathcal{T}]{X[1]}{Z}+\Hom[\mathcal{T}]{Y[1]}{W}\circ f[1].
\]
In a graded endomorphism ring $\widehat{\operatorname{Ext}}{}^*_A(X,X)$, the
bracket of homogeneous classes of degrees $p$, $q$ and $r$ lies in degree
${p+q+r-1}$.

\begin{lemma}
  \label{lemma:toda-juggling}
  In the situation above, let ${u\colon X'\to X}$ be a morphism. Then
  \[
    \langle h,g,f\rangle\circ u[1]\subseteq\langle h,g,fu\rangle.
  \]
\end{lemma}

\begin{proof}
  If $(a,b)$ is a defining system for $\langle h,g,f\rangle$, then
  ${q(a\circ u[1])=(fu)[1]}$ and ${bi=h}$, so that ${(a\circ u[1],b)}$ is a defining
  system for $\langle h,g,fu\rangle$, with value ${(ba)\circ u[1]}$.
\end{proof}

\subsection{Trivial extensions}
\label{subsec:trivial-extensions}

For an algebra $\Delta$ and a $\Delta$-bimodule $X$, the \emph{trivial
extension} is the algebra ${\Delta\ltimes X=\Delta\oplus X}$ with multiplication
\[
  (d,x)(d',x')=(dd',dx'+xd').
\]
Every $\Delta$-module is a ${(\Delta\ltimes X)}$-module through the projection
${\Delta\ltimes X\to\Delta}$; we call such modules \emph{inflated}.

For a finite-dimensional algebra $C$, the trivial extension ${C\ltimes DC}$ by
the bimodule $DC$, with actions ${(a\phi)(c)=\phi(ca)}$ and ${(\phi a)(c)=
\phi(ac)}$, is symmetric: the linear form ${\operatorname{tr}(c,\phi)=\phi(1)}$
satisfies ${\operatorname{tr}((a,\phi)(b,\psi))=\psi(a)+\phi(b)}$, which is
symmetric in the two factors, and the associated bilinear form is
nondegenerate, since for ${a\neq0}$ there is $\psi$ with ${\psi(a)\neq0}$, and for
${a=0\neq\phi}$ there is $b$ with ${\phi(b)\neq0}$. Hence,
${t\mapsto\operatorname{tr}(-\,t)}$ is an isomorphism of bimodules from
${C\ltimes DC}$ to its dual.

\section{Criteria for infinite finitistic dimension}
\label{sec:criteria}

Both constructions use a module with prescribed homological behaviour. The
following dimension-shifting argument produces modules of every finite
projective dimension.

\begin{proposition}
  \label{prop:coresolution}
  Let $\A$ be an abelian category with enough projective objects, and suppose
  given short exact sequences
  \[
    0\to c_0\to p_0\to c_1\to0\qquad\text{and}\qquad
    0\to c_n\to p_n\to c_{n+1}\to0\quad(n\geq1)
  \]
  in which every $p_n$ is projective, $c_0$ is projective and $c_1$ is not.
  Then ${\operatorname{pd}c_n=n}$ for every ${n\geq1}$.
\end{proposition}

\begin{proof}
  A short exact sequence ${0\to c_n\to p_n\to c_{n+1}\to0}$ with $p_n$
  projective gives ${\operatorname{pd}c_{n+1}\leq\operatorname{pd}c_n+1}$, so
  ${\operatorname{pd}c_n\leq n}$ by induction. Since $c_1$ is not projective,
  ${\operatorname{pd}c_1=1}$. Let ${n\geq1}$ and suppose that
  ${\operatorname{pd}c_n=n}$. For every object $V$ and every ${j\geq1}$, the long
  exact sequence of ${\operatorname{Ext}(-,V)}$ and the vanishing of
  $\operatorname{Ext}^{\geq1}(p_n,V)$ give an isomorphism
  ${\operatorname{Ext}^{j+1}(c_{n+1},V)\cong\operatorname{Ext}^j(c_n,V)}$.
  Choosing $V$ with ${\operatorname{Ext}^n(c_n,V)\neq0}$ and ${j=n}$ yields
  ${\operatorname{pd}c_{n+1}=n+1}$.
\end{proof}

\begin{remark}
  \Cref{prop:coresolution} is formally verified in Lean, against Mathlib, in an
  arbitrary abelian category; see \Cref{app:verification}.
\end{remark}

\subsection{The strong Nakayama conjecture}
\label{subsec:strong-nakayama}

The following argument shows that finiteness of the finitistic dimension of
$A^{\op}$ implies the strong Nakayama conjecture for $A$~\cite[Section~3.2,
Proposition~5]{CB19}; it also gives the exact projective dimensions of the
modules constructed.

\begin{theorem}
  \label{thm:strong-nakayama}
  Let $A$ be a finite-dimensional algebra and ${E\neq0}$ a right $A$-module
  such that
  \[
    \operatorname{Ext}^i_{A^{\op}}(E,A)=0\qquad\text{for all }i\geq0.
  \]
  Let ${\cdots\to P_1\to P_0\to E\to0}$ be a minimal projective resolution, and
  put ${C_n}={\operatorname{coker}(P_{n-1}^*}\to{ P_n^*)}$ for ${n\geq1}$. Then
  ${C_n\cong\operatorname{Tr}\Omega^{n-1}E}$ and ${\operatorname{pd}_AC_n=n}$ for
  every ${n\geq1}$. In particular, ${\operatorname{findim}A=\infty}$.
\end{theorem}

\begin{proof}
  The complex ${0\to P_0^*\to P_1^*\to P_2^*\to\cdots}$ of projective left
  $A$-modules has cohomology ${\operatorname{Ext}^*_{A^{\op}}(E,A)=0}$, so it is
  exact, and it splits into short exact sequences
  ${0\to C_0\to P_1^*\to C_1\to0}$ and ${0\to C_n\to P_{n+1}^*\to C_{n+1}\to0}$,
  where ${C_0=P_0^*}$. If $C_1$ were projective, the map ${P_0^*\to P_1^*}$ would be
  a split monomorphism; since finitely generated projective modules are
  reflexive, its dual ${P_1\to P_0}$ would then be a split epimorphism, and
  ${E=\operatorname{coker}(P_1\to P_0)}$ would vanish. Hence,
  \Cref{prop:coresolution} applies. Finally,
  ${P_n\to P_{n-1}\to\Omega^{n-1}E\to0}$ is a minimal presentation, so that
  \[
    \operatorname{Tr}\Omega^{n-1}E=\operatorname{coker}(P_{n-1}^*\to P_n^*)
    =C_n.\qedhere
  \]
\end{proof}

Since ${\operatorname{Ext}^i_{A^{\op}}(E,A)\cong D\operatorname{Tor}^A_i(E,DA)}$,
the hypothesis says that the derived Nakayama functor ${-\Lotimes[A]DA}$
annihilates $E$. It forces ${\operatorname{pd}E=\infty}$: if
${\operatorname{pd}E=n<\infty}$, then ${\operatorname{Ext}^n(E,A)\neq0}$.

\begin{remark}
  \Cref{thm:strong-nakayama} is formally verified in Lean over an arbitrary
  ring, including the equivalence between the vanishing of the groups
  $\operatorname{Ext}^i_{A^{\op}}(E,A)$ and the exactness of the dual complex;
  the identification of $C_n$ with a transpose is proved for the presentation
  given by the resolution, without its minimality. See
  \Cref{app:verification}.
\end{remark}

The next proposition reformulates the hypothesis of \Cref{thm:strong-nakayama}
as a property of a module of projective dimension one, such as $C_1$.

\begin{proposition}
  \label{prop:infinitely-torsionfree}
  Let $C$ be an $A$-module without nonzero projective summands and with
  ${\operatorname{pd}C=1}$. The following are equivalent:
  \begin{enumerate}
    \item the right module ${E=\operatorname{Tr}C}$ satisfies
      ${\operatorname{Ext}^i_{A^{\op}}(E,A)=0}$ for all ${i\geq0}$;
    \item $C$ is not projective and
      ${\operatorname{Ext}^i_{A^{\op}}(\operatorname{Tr}C,A)=0}$ for all ${i\geq1}$;
    \item $C$ is not projective and admits an exact sequence
      ${0\to C\to Q_2\to Q_3\to\cdots}$ with each $Q_m$ projective, in which every
      map from $C$ or from a cokernel to the next $Q_m$ is a left
      $\operatorname{add}A$-approximation.
  \end{enumerate}
  Moreover, ${E\cong\operatorname{Ext}^1_A(C,A)}$.
\end{proposition}

\begin{proof}
  Choose a minimal projective resolution ${0\to Q_0\to Q_1\to C\to0}$. Then
  ${E=\operatorname{coker}(Q_1^*\to Q_0^*)\cong\operatorname{Ext}^1_A(C,A)}$ and
  ${E^*=\ker(Q_0\to Q_1)=0}$; moreover ${E\neq0}$, since otherwise ${Q_1^*\to Q_0^*}$
  would split and $C$ would be projective. This proves that (1) and (2) are
  equivalent. Assume (1). Dualising a minimal projective resolution of $E$ gives
  an exact complex whose first cokernel is ${\operatorname{Tr}E=C}$, since
  ${\operatorname{Tr}\operatorname{Tr}C\cong C}$ for a module $C$ without nonzero
  projective summands, and whose tail
  is a projective coresolution of $C$; the dual of each of its short exact
  sequences is a segment of the resolution of $E$, hence exact, which is the
  approximation property in (3). Conversely, dualising the sequence in (3) gives
  an exact sequence ${\cdots\to Q_3^*\to Q_2^*\to C^*\to0}$; splicing it with
  ${0\to C^*\to Q_1^*\to Q_0^*\to E\to0}$ gives a projective resolution of $E$
  whose dual ${0\to Q_0\to Q_1\to Q_2\to\cdots}$ is exact, and (1) follows.
\end{proof}

\begin{corollary}
  \label{coro:infinitely-many-torsionless}
  If a finite-dimensional algebra $A$ admits a module $E$ as in
  \Cref{thm:strong-nakayama}, then $A$ has infinitely many isomorphism classes
  of indecomposable torsionless modules of finite projective dimension.
\end{corollary}

\begin{proof}
  The modules $C_n$ of \Cref{thm:strong-nakayama} embed into the projective
  modules $P_{n+1}^*$, by the short exact sequences in its proof, so their
  indecomposable summands are torsionless and of finite projective
  dimension. If there were only finitely many of them, their projective
  dimensions would bound ${\operatorname{pd}C_n=n}$ for all $n$.
\end{proof}

\subsection{Counterexamples to the Auslander--Reiten conjecture}
\label{subsec:auslander-reiten}

Let $\Lambda$ be an Artin algebra. The \emph{Auslander--Reiten conjecture}
asserts that a $\Lambda$-module $M$ with
${\operatorname{Ext}^i_\Lambda(M,M\oplus\Lambda)=0}$ for all ${i\geq1}$ is
projective. Auslander and Reiten showed that its failure
produces a failure of the generalised Nakayama conjecture for an endomorphism
algebra~\cite[Theorem~1.1(b)]{AR75a}. Combined with \Cref{thm:strong-nakayama},
this gives a counterexample to the finitistic dimension conjecture with one
more simple module.

\begin{theorem}
  \label{thm:ar-to-findim}
  Let $\Lambda$ be a finite-dimensional algebra and $M$ a nonprojective
  $\Lambda$-module with ${\operatorname{Ext}^i_\Lambda(M,M\oplus\Lambda)}={0}$ for
  all ${i\geq1}$. Put ${G=\Lambda\oplus M}$ and
  ${\Gamma}={\operatorname{End}_\Lambda(G)^{\op}}$,
  and let ${\pi\colon P\to M}$ be a projective cover. Then the $\Gamma$-module
  \[
    S\coloneqq\operatorname{coker}\bigl(\Hom[\Lambda]{G}{\pi}\bigr)
  \]
  is nonzero and ${\operatorname{Ext}^i_\Gamma(S,\Gamma)=0}$ for all ${i\geq0}$.
  Consequently, the algebra $\operatorname{End}_\Lambda(G)$ has infinite
  little finitistic dimension.
\end{theorem}

\begin{proof}
  Write ${\mathbb{F}=\Hom[\Lambda]{G}{-}}$, which restricts to an equivalence
  ${\operatorname{add}G\simeq\operatorname{proj}\Gamma}$ with
  ${\mathbb{F}G=\Gamma}$.
  Put ${K_0=\ker\pi}$ and choose right $\operatorname{add}G$-approximations
  ${g_j\colon G_j\to K_j}$ with ${K_{j+1}=\ker g_j}$; they are surjective since
  ${\Lambda\in\operatorname{add}G}$. The functor $\mathbb{F}$ is left exact and
  the
  maps $\mathbb{F}g_j$ are surjective, so
  \[
    \cdots\to\mathbb{F}G_1\to\mathbb{F}G_0\to\mathbb{F}P\to\mathbb{F}M\to S\to0
  \]
  is a projective resolution of $S$. If $\mathbb{F}\pi$ were surjective, the
  identity of the summand $M$ of $G$ would lift through $\pi$, and $M$ would be
  projective; hence ${S\neq0}$. For ${Y\in\operatorname{add}G}$ there is a natural
  isomorphism ${\Hom[\Gamma]{\mathbb{F}Y}{\Gamma}\cong\Hom[\Lambda]{Y}{G}}$, so
  that
  $\operatorname{Ext}^*_\Gamma(S,\Gamma)$ is the cohomology of
  \[
    0\to\Hom[\Lambda]{M}{G}\to\Hom[\Lambda]{P}{G}\to\Hom[\Lambda]{G_0}{G}\to
    \Hom[\Lambda]{G_1}{G}\to\cdots,
  \]
  with $\Hom[\Lambda]{M}{G}$ in degree zero. Since
  ${\operatorname{Ext}^{\geq1}_\Lambda(G,G)}={0}$, the resolution
  ${\cdots}\to{ G_0}\to{ P}\to{ M}$ consists of $\Hom[\Lambda]{-}{G}$-acyclic modules
  and
  computes $\operatorname{Ext}^*_\Lambda(M,G)$; as
  ${\operatorname{Ext}^{\geq1}_\Lambda(M,G)}={0}$ and the first map is injective,
  the complex is exact. The last assertion is \Cref{thm:strong-nakayama} applied
  to ${\Gamma^{\op}}={\operatorname{End}_\Lambda(G)}$.
\end{proof}

Replacing $M$ by an indecomposable nonprojective summand preserves the
hypotheses of \Cref{thm:ar-to-findim}, since $\operatorname{Ext}$ is additive.
If $\Lambda$ is basic, the algebra $\operatorname{End}_\Lambda(G)$ then has
exactly one simple module more than $\Lambda$.

Conversely, simple modules witnessing failure of the strong Nakayama conjecture
always arise in this
way, on a corner of the algebra. In the next proposition, $Af$, $M$ and their
endomorphisms are taken as \emph{right} $B$-modules.

\begin{proposition}
  \label{prop:simple-witness}
  Let $S$ be a simple right $A$-module with
  ${\Hom[A^{\op}]{S}{A}}={0}={\operatorname{Ext}^1_{A^{\op}}(S,A)}$. Let $e$ be the
  sum
  of the primitive idempotents of a complete decomposition of $1$ corresponding
  to $S$, ${f=1-e}$, ${B=fAf}$ and ${M=eAf}$. Then left multiplication induces an
  algebra isomorphism
  \[
    A\stackrel{\sim}{\longrightarrow}\operatorname{End}_{B^{\op}}(Af)
    =\operatorname{End}_{B^{\op}}(B\oplus M),
  \]
  $B$ has exactly one isomorphism class of simple modules fewer than $A$, and
  $M$ is not a projective right $B$-module. If moreover
  ${\operatorname{Ext}^i_{A^{\op}}(S,A)=0}$ for all ${i\geq0}$, then
  ${\operatorname{Ext}^i_{B^{\op}}(M,B\oplus M)=0}$ for all ${i\geq1}$.
\end{proposition}

\begin{proof}
  For a right $A$-module $V$, let ${\eta_V\colon V\to\Hom[B^{\op}]{Af}{Vf}}$ send
  $v$ to ${x\mapsto vx}$. Multiplication by $f$ identifies $\eta_Vf$ with the
  identity of $Vf$, so the kernel and cokernel of $\eta_V$ are modules over
  $A/AfA$, whose composition factors are copies of $S$. If
  ${\Hom[A^{\op}]{S}{V}=0}$, then $\eta_V$ is injective. If moreover
  ${\operatorname{Ext}^1_{A^{\op}}(S,V)=0}$, then
  ${\operatorname{Ext}^1_{A^{\op}}(T,V)=0}$ for every $A/AfA$-module $T$ by
  induction on the length, so $\eta_V$ splits; its cokernel is a summand of
  $\Hom[B^{\op}]{Af}{Vf}$ annihilated by $f$, hence zero, since
  ${\Hom[A^{\op}]{T}{\Hom[B^{\op}]{Af}{Vf}}\cong\Hom[B^{\op}]{Tf}{Vf}=0}$. For
  ${V=A}$ this gives the isomorphism. Passing from $A$ to $B$ removes exactly the
  isomorphism class of $S$. If $M$ were projective, ${B\oplus M}$ would be a
  progenerator of $B$ and $A$ would be Morita equivalent to $B$, which has fewer
  simple modules.

  Assume now the vanishing in all degrees. A minimal injective resolution
  ${A\to I^\bullet}$ of the right regular module has no summand with socle $S$,
  since the number of such summands of $I^j$ is the dimension of
  ${\operatorname{Ext}^j_{A^{\op}}(S,A)}$ over ${\operatorname{End}_{A^{\op}}(S)}$,
  so each $I^j$ is a sum of modules $D(Ae_t)$ with ${e_t\leq f}$. Then
  $I^jf$ is a direct sum of modules $D(Be_t)$, hence an injective right
  $B$-module, and, by the first part of the proof applied to ${V=I^j}$, the
  complex
  ${\Hom[B^{\op}]{Af}{I^\bullet f}\cong\Hom[A^{\op}]{A}{I^\bullet}}$ computes
  $\operatorname{Ext}^*_{B^{\op}}(Af,Af)$, which therefore vanishes in positive
  degrees.
\end{proof}

\begin{example}
  The two vanishing conditions in the first part of
  \Cref{prop:simple-witness} do not suffice for the last assertion. Let $S$ be
  the right simple module at the vertex $3$ of the algebra
  ${A}={\kk(1\xrightarrow{a}2\xrightarrow{b}3)/(ba)}$. Its minimal projective resolution is
  ${0\to e_1A\to e_2A\to e_3A\to S\to0}$, and the dual complex shows that
  ${\Hom[A^{\op}]{S}{A}}={0}={\operatorname{Ext}^1_{A^{\op}}(S,A)}$, while
  ${\operatorname{Ext}^2_{A^{\op}}(S,A)}\cong{ Ae_1/\kk a}\neq{0}$. The
  conclusion fails as well: here ${B=\kk(1\xrightarrow{a}2)}$ and $M$ is the simple
  right $B$-module at the vertex $2$, with projective resolution
  ${0\to e_1B\to e_2B\to M\to0}$, so that
  ${\operatorname{Ext}^1_{B^{\op}}(M,B)\cong Be_1/\kk a\neq0}$.
\end{example}

\subsection{Triangular matrix algebras and bounds on projective dimension}
\label{subsec:gluing}

Triangular matrix algebras cannot produce a counterexample to the strong
Nakayama conjecture when both
diagonal algebras have finite little finitistic dimension.

\begin{proposition}
  \label{prop:triangular}
  Let ${A=\begin{psmallmatrix}B&0\\M&C\end{psmallmatrix}}$ for a $C$-$B$-bimodule
  $M$, so that right $A$-modules are triples $(X,Y,\varphi)$ with $X$ a right
  $B$-module, $Y$ a right $C$-module and ${\varphi\colon Y\otimes_CM\to X}$. Let
  ${E=(X,Y,\varphi)\neq0}$ satisfy ${\RHom[A^{\op}]{E}{A}=0}$, and put
  ${Z=\operatorname{cone}(Y\Lotimes[C]M\to X)}$ in
  ${\DerCat[-]{\operatorname{mod}B^{\op}}}$.
  Then ${\RHom[B^{\op}]{Z}{B}=0}$. If ${Z\neq0}$, the little finitistic dimension of
  $B$ is infinite; if ${Z=0}$, then ${Y\neq0}$, ${\RHom[C^{\op}]{Y}{C}=0}$, and the
  little finitistic dimension of $C$ is infinite.
\end{proposition}

\begin{proof}
  Let ${e=\operatorname{diag}(1,0)}$ and ${f=\operatorname{diag}(0,1)}$, let $i_*$
  inflate right $B$-modules and put ${j_!Y=Y\Lotimes[C]fA}$. The derived counit
  gives
  a triangle ${j_!Y\to E\to i_*Z\to(j_!Y)[1]}$. Since ${eA=i_*B}$, adjunction gives
  ${\RHom[B^{\op}]{Z}{B}\cong\RHom[A^{\op}]{E}{eA}=0}$. If ${Z\neq0}$, represent it
  by a nonzero minimal bounded-above complex ${P^\bullet}$ of finitely generated
  projective right $B$-modules, with largest nonzero degree $m$. Its dual
  ${0\to(P^m)^*\to(P^{m-1})^*\to\cdots}$ is exact, and its first cokernel is not
  projective, since otherwise the first map would split, against the
  minimality of ${P^\bullet}$; \Cref{prop:coresolution} applied to the successive
  cokernels gives unbounded finite projective dimensions. If ${Z=0}$, then
  ${E\simeq j_!Y}$, so ${Y\neq0}$, and
  ${\RHom[C^{\op}]{Y}{C}\cong\RHom[A^{\op}]{j_!Y}{fA}=0}$ since ${(fA)f=C}$; now
  \Cref{thm:strong-nakayama} applies to $C$.
\end{proof}

Finite projective dimensions are bounded on modules of bounded
dimension. The argument is due to Schofield, as explained by Happel in his
survey~\cite[Section~2.3]{Hap90}; the openness also follows from the
semicontinuity of $\dim\operatorname{Ext}$ on module
varieties~\cite[Corollary~2.6]{GLS24}. Both sources work over an algebraically
closed field; the argument below works over any field.

\begin{proposition}
  \label{prop:bounded-dimension}
  Let $A$ be a finite-dimensional algebra and ${d\geq0}$ an integer. The
  subset of the variety $\operatorname{Rep}_d(A)$ of $d$-dimensional modules of
  projective dimension at most $n$ is open, and projective dimension is bounded
  on the modules of finite projective dimension in $\operatorname{Rep}_d(A)$.
  In particular, modules of unbounded finite projective dimension have unbounded
  dimension.
\end{proposition}

\begin{proof}
  Let $J$ be the radical of $A$. A module $M$ satisfies
  ${\operatorname{pd}M\leq n}$ if and only if
  ${\operatorname{Tor}^A_{n+1}(A/J,M)=0}$. Choose a resolution ${Q_\bullet}$ of the
  right module $A/J$ by finitely generated free modules. Over
  $\operatorname{Rep}_d(A)$, the complex
  \[
    Q_{n+2}\otimes_AM\longrightarrow Q_{n+1}\otimes_AM\longrightarrow
    Q_n\otimes_AM
  \]
  has terms of fixed dimension and differentials given by matrices whose entries
  are polynomial in the point $M$; the vanishing of its middle homology is a
  lower bound on the sum of the ranks of the two differentials, which is an open
  condition. The open subsets ${\set{M}[\operatorname{pd}M\leq n]}$ form an
  ascending chain in a noetherian topological space and are therefore
  stationary.
\end{proof}

In \Cref{thm:strong-nakayama}, the dimension of $\operatorname{Tr}N$ is at
most $(\dim A)^2\dim N$. Hence, by \Cref{prop:bounded-dimension}, the syzygies
$\Omega^nE$ of a module $E$ as in \Cref{thm:strong-nakayama} have unbounded
dimension:
families of modules of constant dimension whose syzygies move along a
parameter, such as those of Schulz~\cite[Example~7]{Schu95}, cannot themselves
satisfy the hypothesis of that theorem.

\section{Trivial extensions and extinction}
\label{sec:extinction}

Let $\Delta$ be a finite-dimensional algebra and $X$ a finite-dimensional
$\Delta$-bimodule, and put
\[
  A=\Delta\ltimes X,\qquad\Phi=X\Lotimes[\Delta]-.
\]
For an inflated $A$-module $N$, ${\Phi^rN\not\simeq0}$ forces
${\operatorname{pd}_AN\geq r}$, and when $\Delta$ has finite global dimension, an
inflated module $N$ has finite projective dimension exactly when
${\Phi^tN\simeq0}$ for some $t$. Hence, for such $\Delta$,
${\operatorname{findim}A=\infty}$ as soon as finite extinction times of $\Phi$
are unbounded on inflated modules. Both statements are due to Minamoto and
Yamaura~\cite{MY20}; we give a proof in the present conventions, without
assuming that $X$ is flat on either side.

\subsection{The bar decomposition}
\label{subsec:bar-decomposition}

Complexes are graded cohomologically, with the conventions of
\Cref{sec:preliminaries}; the differential of a tensor product of complexes is
${d(u\otimes v)=du\otimes v+(-1)^{|u|}u\otimes dv}$. We write
${\DerCat[-]{\operatorname{Mod}\Delta}}$ for the derived category of
bounded-above complexes of left $\Delta$-modules.

\begin{theorem}[{Minamoto--Yamaura~\cite[Lemma~4.13(4) and proof of
    Theorem~4.17]{MY20}}]
  \label{thm:bar-decomposition}
  Let $N$ be a left $\Delta$-module, regarded as an $A$-module through the
  projection ${A\to\Delta}$. There is an isomorphism, natural in~$N$,
  \[
    \Delta\Lotimes[A]N\cong\bigoplus_{r\geq0}\Phi^rN[r]
    \quad\text{in }\DerCat[-]{\operatorname{Mod}\Delta},
  \]
  where ${\Phi^0N=N}$.
\end{theorem}

The source states the decomposition in terms of the internal grading of the
graded algebra $A$, for right modules; the statement above is its translation
to left modules through the opposite algebras. The proof below avoids the
grading, and it does not assume that $X$ is flat on either side. That the
derived powers of $X$ appear, and not the ordinary ones, matters already in
small examples; see \Cref{ex:derived-powers}.

\begin{proof}
  Choose a resolution ${\widetilde{X}\to X}$ by finitely generated projective
  $\Delta$-bimodules with ${\widetilde{X}^j=0}$ for ${j>0}$, and a projective
  resolution ${\widetilde{N}\to N}$ over $\Delta$ with ${\widetilde{N}^j=0}$ for
  ${j>0}$. The complex ${\widetilde{A}=\Delta\oplus\widetilde{X}}$, with
  multiplication ${(d,x)(d',x')=(dd',dx'+xd')}$, is a dg algebra in which
  $\widetilde{X}$ is an ideal of square zero, and the augmentation
  ${\widetilde{A}\to A}$ is a quasi-isomorphism of dg algebras.

  \emph{The bar module.} For ${r\geq0}$ put
  ${F_r=\widetilde{X}^{\otimes_\Delta r}\otimes_\Delta\widetilde{N}}$, with
  ${F_0=\widetilde{N}}$; each $F_r$ is a bounded-above complex of projective left
  $\Delta$-modules, since for a projective bimodule $V$ and any left module $L$
  the module ${V\otimes_\Delta L}$ is a direct summand of a direct sum of modules
  ${\Delta\otimes_\kk L}$. Let
  \[
    \mathcal{B}=\bigoplus_{r\geq0}(\widetilde{A}\otimes_\Delta F_r)[r],
  \]
  and write $s^rw$ for the element of the summand of index $r$ corresponding
  to ${w\in\widetilde{A}\otimes_\Delta F_r}$. Let $\widetilde{A}$ act by
  ${a\cdot s^rw=(-1)^{r|a|}s^r(aw)}$, and define the differential by
  \[
    \begin{gathered}
      D(s^rw)=(-1)^rs^r\,dw+b(s^rw),\\
      b\bigl(s^r(a\otimes q_1\otimes\cdots\otimes q_r\otimes v)\bigr)=
      s^{r-1}(aq_1\otimes q_2\otimes\cdots\otimes q_r\otimes v)
    \end{gathered}
  \]
  for ${r>0}$, and ${b=0}$ on the summand of index zero; here $d$ is the
  differential of the tensor product ${\widetilde{A}\otimes_\Delta F_r}$. The
  multiplication ${\widetilde{A}\otimes_\Delta F_r\to\widetilde{A}\otimes_\Delta
  F_{r-1}}$ is a chain map, so the mixed terms of $D^2$ cancel, with
  coefficients ${(-1)^r+(-1)^{r-1}}$; and ${b^2=0}$ because ${q_1q_2=0}$ in
  $\widetilde{A}$. Moreover, ${b(a\cdot z)=(-1)^{|a|}a\cdot b(z)}$, since
  ${r|a|\equiv|a|+(r-1)|a|}$ modulo two. Thus, $\mathcal{B}$ is a dg
  $\widetilde{A}$-module. The other faces of the usual bar construction
  vanish: they multiply two elements of $\widetilde{X}$, or let
  $\widetilde{X}$ act on $\widetilde{N}$.

  \emph{The bar module resolves $N$.} Let ${\epsilon\colon\mathcal{B}\to
  \widetilde{N}}$ be the map induced by ${\widetilde{A}\to\Delta}$ on the summand
  of index zero and zero elsewhere, where $\widetilde{N}$ is regarded as an
  $\widetilde{A}$-module through ${\widetilde{A}\to\Delta}$. As a complex of
  left $\Delta$-modules, ${\widetilde{A}\otimes_\Delta F_r=F_r\oplus F_{r+1}}$,
  according to the decomposition ${\widetilde{A}=\Delta\oplus\widetilde{X}}$.
  Hence, $\mathcal{B}$ is the direct sum of $\widetilde{N}$, in index zero, and
  of the complexes
  \[
    F_r[r]\oplus F_r[r-1]\qquad(r\geq1),
  \]
  in which the first summand is the part of index $r$ with leading factor in
  $\Delta$, the second is the part of index ${r-1}$ with leading factor in
  $\widetilde{X}$, and the component of $D$ from the first to the second is the
  identity of the underlying graded module, given by $b$. Each of these
  complexes is the mapping cone of an identity map, so the map of degree ${-1}$
  sending the second summand identically to the first is a contraction. These
  contractions give $h$ with ${Dh+hD=\id-\iota\epsilon}$, where $\iota$ is the
  inclusion of $\widetilde{N}$. Hence, $\epsilon$ is a quasi-isomorphism.

  \emph{Base change to $A$.} Filter $\mathcal{B}$ by the subcomplexes of the
  summands of index at most $r$. The subquotients are
  ${(\widetilde{A}\otimes_\Delta F_r)[r]}$, and the filtration is split as a
  filtration of graded $\widetilde{A}$-modules. For an acyclic right dg
  $\widetilde{A}$-module $Z$, the complex
  \[
    Z\otimes_{\widetilde{A}}(\widetilde{A}\otimes_\Delta F_r)[r]\cong
    (Z\otimes_\Delta F_r)[r],
  \]
  with the sign ${(-1)^{r|z|}}$ on ${z\otimes s^rw}$, is acyclic, because a
  bounded-above complex of projective modules is K-flat: its stupid truncations
  from below are finite extensions of projective modules, and it is their
  filtered colimit. By induction on $r$ and passage to the colimit,
  ${Z\otimes_{\widetilde{A}}\mathcal{B}}$ is acyclic, so $\mathcal{B}$ is K-flat
  over $\widetilde{A}$. Since ${\widetilde{A}\to A}$ is a quasi-isomorphism, so is
  ${\mathcal{B}\to L\coloneqq A\otimes_{\widetilde{A}} \mathcal{B}}$, and the
  augmentation ${L\to N}$ induced by $\epsilon$ and ${\widetilde{N}\to N}$ is a
  quasi-isomorphism. The terms of $L$ are direct sums of modules
  ${(A\otimes_\Delta F_r^j)}$, which are projective over $A$, and they vanish in
  positive degrees; in each degree only finitely many indices $r$ contribute.
  Hence, $L$ is a bounded-above projective resolution of~$N$.

  \emph{Tensoring with $\Delta$.} In ${\Delta\otimes_AL}$ the bar differential
  vanishes, since its image has a leading factor in the ideal $X$, which acts
  as zero on~$\Delta$. Hence,
  \[
    \Delta\otimes_AL\cong\bigoplus_{r\geq0}F_r[r],
  \]
  with the differential ${(-1)^rd_{F_r}}$ on the summand of index $r$, which is
  the differential of $F_r[r]$. Finally, $F_r$ represents $\Phi^rN$: by
  induction on $r$, $F_{r-1}$ is a bounded-above complex of projective modules
  representing $\Phi^{r-1}N$, hence K-flat, and
  ${F_r=\widetilde{X}\otimes_\Delta F_{r-1}\to X\otimes_\Delta F_{r-1}}$ is a
  quasi-isomorphism. Naturality follows by choosing $\widetilde{N}$
  functorially, for instance as the bar resolution of $N$ over~$\kk$.
\end{proof}

\begin{example}
  \label{ex:derived-powers}
  Let ${\Delta=\kk(0\xrightarrow{a}1\xrightarrow{b}2)}$, a path algebra of
  dimension six, and let ${X=S_0\otimes_\kk S_2'}$, where $S_i$ and $S_i'$ denote
  the simple left and right modules at the vertex $i$. Then $A$ is the quotient
  of the path algebra of the cyclic quiver ${0\to1\to2\to0}$, with the new arrow
  ${x\colon2\to0}$, by the relations ${ax=0=xb}$. For ${N=S_1}$ the ordinary tensor
  product ${X\otimes_\Delta N}$ is zero, whereas the projective resolution
  \[
    0\longrightarrow e_1\Delta\xrightarrow{\,b\cdot\,}e_2\Delta\longrightarrow
    S_2'\longrightarrow0
  \]
  gives ${\Phi N\cong S_0[1]}$ and ${\Phi^2N=0}$. The minimal projective resolution
  \[
    0\longrightarrow Ae_1\xrightarrow{\,\cdot a\,}Ae_0\xrightarrow{\,\cdot x\,}
    Ae_2\xrightarrow{\,\cdot b\,}Ae_1\longrightarrow S_1\longrightarrow0,
  \]
  by right multiplications, shows that ${\operatorname{pd}_AS_1=3}$, in agreement
  with \Cref{prop:pd-formula}: the two finite terms of the formula are
  ${\operatorname{pd}_\Delta S_1=1}$ and ${\operatorname{pd}_\Delta(S_0[1])+1=3}$.
\end{example}

\subsection{Projective dimension and extinction}
\label{subsec:pd-formula}

Let ${\DerCat[-]{\operatorname{mod}\Delta}}$ be the derived category of
bounded-above complexes with finitely generated cohomology. For a nonzero
object $C$ of it, let $\operatorname{pd}_\Delta C$ be the least integer $p$ such
that $C$ is represented by a bounded-above complex of projective modules
vanishing in degrees below ${-p}$, or $\infty$ if there is none; for a module
this is its projective dimension, and
${\operatorname{pd}(C[s])=\operatorname{pd}C+s}$. We put
${\operatorname{pd}0=-\infty}$.

\begin{lemma}
  \label{lemma:pd-detection}
  Let $\Lambda$ be a finite-dimensional algebra and $C$ a nonzero object of
  ${\DerCat[-]{\operatorname{mod}\Lambda}}$. Then
  \[
    \operatorname{pd}_\Lambda C=\sup\set{n\in\ZZ}[\Hom{C}{S[n]}\neq0
    \text{ for a simple module }S].
  \]
\end{lemma}

\begin{proof}
  Represent $C$ by a bounded-above complex $J$ of finitely generated
  projective modules. If some differential ${J^{n}\to J^{n+1}}$ is nonzero on the
  tops of its source and target, it has an invertible component between
  isomorphic indecomposable summands, and a change of bases splits off a
  contractible complex ${P\xrightarrow{\id}P}$. Performing these cancellations
  from the highest degree downwards, finitely many in each degree, we may assume
  that $J$ is minimal: every differential has image in the radical of its
  target. Then $\Hom[\Lambda]{J}{S}$ has zero differential for every simple
  $S$, so ${\Hom{C}{S[n]}=\Hom[\Lambda]{J^{-n}}{S}}$, which is nonzero for some
  $S$ exactly when ${J^{-n}\neq0}$. A representative vanishing below ${-p}$ forces
  these groups to vanish for ${n>p}$, and the minimal one realises the bound.
\end{proof}

\begin{proposition}[{Minamoto--Yamaura~\cite[Corollary~4.11]{MY20}}]
  \label{prop:pd-formula}
  Let $N$ be a nonzero finitely generated left $\Delta$-module, regarded as an
  $A$-module. Then
  \[
    \operatorname{pd}_AN=\sup_{r\geq0}\bigl(\operatorname{pd}_\Delta(\Phi^rN)
    +r\bigr).
  \]
  In particular, ${\Phi^rN\not\simeq0}$ implies ${\operatorname{pd}_AN\geq r}$.
\end{proposition}

\begin{proof}
  Every simple $A$-module $S$ is annihilated by $X$: the submodule $XS$ is $0$
  or $S$, and ${XS=S}$ would give ${S=X^2S=0}$. Hence, the simple $A$-modules are
  the
  simple $\Delta$-modules, inflated. For such $S$, adjunction and
  \Cref{thm:bar-decomposition} give
  \[
    \operatorname{Ext}^n_A(N,S)\cong\Hom[\DerCat{\Delta}]{\Delta\Lotimes[A]N}
    {S[n]}\cong\prod_{r\geq0}\Hom[\DerCat{\Delta}]{\Phi^rN}{S[n-r]},
  \]
  where the factors with ${r>n}$ vanish since $\Phi^rN$ is represented by a
  complex concentrated in nonpositive degrees. Taking the supremum over $n$
  and $S$ and applying \Cref{lemma:pd-detection} over $A$ and over $\Delta$
  gives the formula. If ${\Phi^rN\not\simeq0}$, then
  ${\operatorname{pd}_\Delta\Phi^rN\geq0}$, again by \Cref{lemma:pd-detection}.
\end{proof}

\begin{corollary}
  \label{coro:extinction}
  Suppose that the left and the right global dimensions of $\Delta$ are at most
  $d_L$ and $d_R$, respectively, and let $N$ be a nonzero finitely generated
  left $\Delta$-module. Then ${\operatorname{pd}_AN<\infty}$ if and only if
  ${\Phi^tN\simeq0}$ for some ${t\geq1}$, and in that case
  \[
    \operatorname{pd}_AN=\max_{0\leq r<t}\bigl(\operatorname{pd}_\Delta
    (\Phi^rN)+r\bigr)\leq d_L+(t-1)(d_R+1).
  \]
  In particular, if the finite extinction times ${\min\set{t}[\Phi^tN\simeq0]}$
  of finitely generated $\Delta$-modules $N$ are unbounded, then
  ${\operatorname{findim}A=\infty}$.
\end{corollary}

\begin{proof}
  Truncating a projective bimodule resolution of $X$ at its $d_R$-th syzygy,
  which is projective as a right module, gives a resolution ${X'\to X}$ in
  degrees ${[-d_R,0]}$ by bimodules that are projective as right modules. Its
  tensor powers compute the iterates of $\Phi$, so $\Phi^rN$ has cohomology in
  degrees ${[-rd_R,0]}$. A complex $C$ with finitely generated cohomology in
  degrees $[u,0]$ satisfies ${\operatorname{pd}C\leq d_L-u}$, by
  \Cref{lemma:pd-detection} applied to the triangles of its cohomology
  truncations: in a triangle ${C_1\to C\to C_2\to C_1[1]}$, if
  ${\Hom{C_1}{S[n]}=0=\Hom{C_2}{S[n]}}$, then ${\Hom{C}{S[n]}=0}$, and
  ${H^i(C)[-i]}$ has projective dimension at most ${d_L-i}$. Hence,
  ${\operatorname{pd}_\Delta\Phi^rN\leq d_L+rd_R}$. If ${\Phi^tN\simeq0}$, then
  ${\Phi^rN\simeq0}$ for all ${r\geq t}$, and \Cref{prop:pd-formula} gives the
  displayed equality and bound. Conversely, if ${\operatorname{pd}_AN=p<\infty}$,
  then ${\Phi^{p+1}N\simeq0}$ by \Cref{prop:pd-formula}. For the last assertion,
  a module of extinction time ${t\geq1}$ has ${\Phi^{t-1}N\not\simeq0}$, so its
  projective dimension is finite and at least ${t-1}$.
\end{proof}

\subsection{Two constraints on extinction}
\label{subsec:extinction-constraints}

\begin{remark}
  \label{rem:k0-invisibility}
  Suppose that $\Delta$ has finite global dimension, with $n$ simple modules.
  The functor $\Phi$ preserves the bounded derived category of finitely
  generated modules, by the proof of \Cref{coro:extinction}, and induces an
  endomorphism $[\Phi]$ of ${K_0(\Delta)\cong\ZZ^n}$. If ${\Phi^tN\simeq0}$, then
  ${[\Phi]^t[N]=0}$, so $[N]$ lies in the generalised kernel of $[\Phi]$ on
  $\mathbb{Q}^n$, which is annihilated by $[\Phi]^n$. Hence,
  ${[\Phi^rN]=[\Phi]^r[N]=0}$ for ${r\geq n}$, although $\Phi^rN$ may be nonzero. In
  the main construction,
  \Cref{prop:simulation} and \Cref{thm:realisation} give
  ${\Phi^2M(Y)\simeq M(HY)[b]\oplus M(HY)[b+3]}$, whose class is zero because
  the two shifts differ by an odd integer, as in \Cref{prop:odd-double}.
\end{remark}

For an integer ${d\geq0}$, let $\operatorname{Rep}_d(\Delta)$ be the set of
$\kk$-algebra homomorphisms ${\Delta\to M_d(\kk)}$, that is, of left
$\Delta$-module structures on $\kk^d$. It is a Zariski-closed subset of the
affine space $\Hom[\kk]{\Delta}{M_d(\kk)}$, and therefore a noetherian
topological space.

\begin{proposition}
  \label{prop:bounded-extinction}
  Suppose that $\Delta$ has finite right global dimension, and let ${d\geq0}$.
  Then the finite extinction times ${\min\set{t}[\Phi^tN\simeq0]}$ of the modules
  ${N\in\operatorname{Rep}_d(\Delta)}$ are bounded.
\end{proposition}

\begin{proof}
  Let ${X'\to X}$ be a bounded resolution by finite-dimensional bimodules that
  are projective as right modules, as in the proof of \Cref{coro:extinction}.
  For each $t$, the complex $X'^{\otimes_\Delta t}$ is bounded, with
  finite-dimensional terms projective as right modules, and computes $\Phi^t$.
  Let $f_1,\dots,f_s$ be a complete set of orthogonal primitive idempotents of
  $\Delta$. Every finitely generated projective right $\Delta$-module is
  isomorphic to $\bigoplus_if_i\Delta^{m_i}$ for some integers $m_i$, so that
  ${\dim_\kk(V\otimes_\Delta N)}={\sum_im_i\dim_\kk f_iN}$ for such a module $V$.
  The function ${N}\mapsto{\dim_\kk f_iN}$ is the rank of an idempotent matrix
  depending polynomially on $N$, and it is locally constant: the conditions
  ${\operatorname{rank}f_i\geq c}$ and ${\operatorname{rank}(1-f_i)\geq d-c}$ are
  open, and the two ranks add up to $d$. Hence, $\operatorname{Rep}_d(\Delta)$
  is the disjoint union of finitely many open and closed subsets on each of
  which all the numbers $\dim_\kk f_iN$ are constant.

  Fix such a subset $Z$ and an integer ${t\geq0}$, and write the terms of
  $X'^{\otimes_\Delta t}$ as direct summands $\varepsilon_j\Delta^{n_j}$ of free
  right modules, with idempotent matrices $\varepsilon_j$ over~$\Delta$. For
  ${N\in Z}$, the complex ${C=X'^{\otimes_\Delta t}\otimes_\Delta N}$ is then
  identified with a complex of subspaces ${C^j=\varepsilon_j(N^{n_j})}$, whose
  dimensions do not depend on $N$, and whose differentials are the restrictions
  of linear maps ${D^j\colon N^{n_j}\to N^{n_{j+1}}}$ vanishing on the
  complements ${(1-\varepsilon_j)(N^{n_j})}$ and given by matrices whose entries
  are polynomial functions of~$N$. The rank of the differential ${C^j\to
  C^{j+1}}$ is the rank of $D^j$. The complex is exact at $C^j$ if and only if
  ${\operatorname{rank}D^{j-1}+\operatorname{rank}D^j\geq\dim C^j}$, which is an
  open condition. Hence, the sets ${U_t=\set{N\in Z}[\Phi^tN\simeq0]}$ are
  open. They form an ascending chain, since $\Phi$ preserves zero objects, so
  the chain is stationary, say from $T_Z$ on. Every ${N\in Z}$ of finite
  extinction time therefore has extinction time at most $T_Z$, and the maximum
  of the finitely many $T_Z$ bounds all finite extinction times in
  $\operatorname{Rep}_d(\Delta)$.
\end{proof}

Compare \Cref{prop:bounded-dimension}, which bounds finite projective dimensions
on modules of bounded dimension.

\section{Selection data}
\label{sec:selection}

The main construction uses an exact tensor endofunctor on finite-dimensional
modules whose finite extinction times are unbounded. The example
in~\cite[Section~2]{OAI26findim} is defined over the complex group algebra of a
finitely presented group.

\subsection{Selection data and extinction}
\label{subsec:selection-data}

\begin{definition}
  \label{def:selection-data}
  Let $\kk$ be a field. \emph{Selection data} over $\kk$ consist of a
  $\kk$-algebra $R$ and an $R$-bimodule $\Psi$ on which the two actions of
  $\kk$ agree and which is finitely generated and projective as a right
  $R$-module. The \emph{selection functor} is
  \[
    H=\Psi\otimes_R-\colon\operatorname{mod}_{\mathrm{fd}}R\longrightarrow
    \operatorname{mod}_{\mathrm{fd}}R.
  \]
  The \emph{extinction time} of a
  finite-dimensional left $R$-module $Y$ is the least integer ${t\geq0}$ with
  ${H^tY=0}$, or $\infty$ if there is none. The selection data have
  \emph{unbounded extinction} if the finite extinction times of
  finite-dimensional modules are unbounded.
\end{definition}

The functor $H$ is well defined and exact: $\Psi_R$ is a direct summand of some
$R^n$, so that, as a vector space, $HY$ is a direct summand of $Y^n$, naturally
in $Y$, and exactness is detected on underlying vector spaces; see
\Cref{lemma:selection-matrix} for an explicit description. Iterates of $H$ are
written $H^t$, with $H^0$ the identity.

\begin{example}
  \label{ex:twisted-corner}
  Let ${e\in R}$ be a central idempotent and ${\alpha\colon R\to R}$ a homomorphism
  of unital $\kk$-algebras, and let ${\Psi={}_\alpha(eR)}$ be the bimodule $eR$
  with its right regular action and the left action ${r\cdot x=\alpha(r)x}$.
  Then $\Psi$ is projective as a right module, and ${HY\cong{}_\alpha(eY)}$ is
  the $R$-submodule $eY$ of $Y$ with $r$ acting as $\alpha(r)$: the map
  ${\Psi\otimes_RY\to eY}$, ${x\otimes y\mapsto xy}$, has the inverse
  ${v\mapsto e\otimes v}$, since ${e\otimes xy=ex\otimes y=x\otimes y}$. This is
  the case treated in~\cite[Section~2]{OAI26findim}.
\end{example}

Selection data with unbounded extinction must detect infinitely many
independent numerical invariants of finite-dimensional modules. To make this
precise, let $K_0(R^{\op})$ be the Grothendieck group of finitely generated
projective right $R$-modules. For a finite-dimensional left $R$-module $Y$ the
\emph{rank function} of $Y$ is the additive map
\[
  \chi_Y\colon K_0(R^{\op})\longrightarrow\ZZ,\qquad
  [P]\longmapsto\dim_\kk(P\otimes_RY),
\]
and the bimodule $\Psi$ induces the endomorphism
${\mathsf{T}[P]=[P\otimes_R\Psi]}$
of $K_0(R^{\op})$; here ${P\otimes_R\Psi}$ is finitely generated and projective
since $P$ is a direct summand of some $R^n$. Associativity of the tensor
product gives
\begin{equation}
  \label{eq:rank-shift}
  \chi_Y(\mathsf{T}x)=\chi_{HY}(x)\qquad\text{and}\qquad
  \dim_\kk H^tY=\chi_Y(\mathsf{T}^t[R])
\end{equation}
for every ${x\in K_0(R^{\op})}$ and every ${t\geq0}$. Let $\mathcal{F}$ be the
$\mathbb{Q}$-vector space of functions from the set of isomorphism classes of
finite-dimensional left $R$-modules to $\mathbb{Q}$, and let
\[
  \chi\colon K_0(R^{\op})\otimes_\ZZ\mathbb{Q}\longrightarrow\mathcal{F},
  \qquad x\longmapsto(Y\mapsto\chi_Y(x)).
\]
For the selection data, put
\[
  \mathcal{V}=\chi\bigl(\operatorname{span}_{\mathbb{Q}}
  \set{\mathsf{T}^t[R]}[t\geq0]\bigr)\subseteq\mathcal{F}.
\]

\begin{proposition}
  \label{prop:rank-obstruction}
  Let $(R,\Psi)$ be selection data, and suppose that ${r=\dim_{\mathbb Q}\mathcal
  V<\infty}$. Then every
  finite-dimensional left $R$-module of finite extinction time has extinction
  time at most~$r$.
\end{proposition}

\begin{proof}
  By~\eqref{eq:rank-shift}, if ${\chi(x)=0}$ then ${\chi(\mathsf{T}x)=0}$, since
  $HY$ is again finite-dimensional. Hence, $\mathsf{T}$ induces an endomorphism
  $\overline{\mathsf{T}}$ of $\mathcal{V}$, and $\mathcal{V}$ is spanned by the
  vectors $\overline{\mathsf{T}}^tv$ for ${t\geq0}$, where ${v=\chi([R])}$. The
  Fitting decomposition of $\overline{\mathsf{T}}$ gives
  ${\mathcal{V}=\mathcal{V}_0\oplus\mathcal{V}_1}$, where $\overline{\mathsf{T}}$
  is nilpotent on $\mathcal{V}_0$, hence
  ${\overline{\mathsf{T}}^r\mathcal{V}_0=0}$,
  and invertible on $\mathcal{V}_1$. Write ${v=v_0+v_1}$ accordingly, and let
  ${\mathcal{V}_2\subseteq\mathcal{V}_1}$ be the span of the vectors
  $\overline{\mathsf{T}}^tv_1$ for ${t\geq0}$. The restriction of
  $\overline{\mathsf{T}}$ to $\mathcal{V}_2$ is injective, hence bijective, so
  for every $t_0$ the vectors $\overline{\mathsf{T}}^tv_1$ with ${t\geq t_0}$
  span ${\overline{\mathsf{T}}^{t_0}\mathcal{V}_2=\mathcal{V}_2}$.

  Let $Y$ be a finite-dimensional module of finite extinction time $t_0$, and
  let ${\phi_Y\colon\mathcal{V}\to\mathbb{Q}}$ be evaluation at $Y$. Then
  ${\phi_Y(\overline{\mathsf{T}}^tv)=\dim_\kk H^tY=0}$ for ${t\geq t_0}$
  by~\eqref{eq:rank-shift}. For ${t\geq\max(t_0,r)}$ we have
  ${\overline{\mathsf{T}}^tv=\overline{\mathsf{T}}^tv_1}$, so $\phi_Y$ vanishes on
  $\mathcal{V}_2$. Hence,
  \[
    \dim_\kk H^rY=\phi_Y(\overline{\mathsf{T}}^rv_0)+
    \phi_Y(\overline{\mathsf{T}}^rv_1)=0.\qedhere
  \]
\end{proof}

\begin{corollary}
  \label{coro:rank-obstruction}
  Selection data $(R,\Psi)$ over $\kk$ have bounded extinction in each of the
  following cases.
  \begin{enumerate}
    \item\label{it:rank-finite-list} There are finitely many finitely
      generated projective right $R$-modules such that every finitely
      generated projective right $R$-module is a finite direct sum of copies
      of them; for example, $R$ is finite-dimensional.
    \item\label{it:rank-noetherian} The algebra $R$ is commutative and
      noetherian.
    \item\label{it:rank-localisation} The algebra $R$ is a universal
      localisation of a finite-dimensional hereditary algebra.
  \end{enumerate}
\end{corollary}

\begin{proof}
  By \Cref{prop:rank-obstruction}, it suffices to show that $\chi$ has image
  of finite dimension. In case~\eqref{it:rank-finite-list} the classes of the
  finitely many modules span ${K_0(R^{\op})\otimes_\ZZ\mathbb{Q}}$; if $R$ is
  finite-dimensional, the Krull--Schmidt theorem shows that every finitely
  generated projective module is a finite direct sum of indecomposable direct
  summands of $R$, of which there are finitely many up to isomorphism. In
  case~\eqref{it:rank-localisation}, let ${R=B_\Sigma}$ be a universal
  localisation of a finite-dimensional hereditary algebra $B$. By a theorem of
  Schofield~\cite[Theorem~2.3]{Sch07a}, every universal localisation of $B$ is
  of the form studied in~\cite{Sch07a}, and for these the map from
  $K_0(B^{\op})$ to $K_0(B_\Sigma^{\op})$, for right modules as
  in~\cite{Sch07a}, is surjective~\cite[Lemma~4.1]{Sch07a}. Since
  $K_0(B^{\op})$ is finitely generated by case~\eqref{it:rank-finite-list}, so
  is $K_0(B_\Sigma^{\op})$.

  In case~\eqref{it:rank-noetherian}, for a finitely generated projective
  $R$-module $P$, the function
  \[
    \rho_P\colon\operatorname{Spec}R\longrightarrow\ZZ,\qquad
    \mathfrak{p}\longmapsto
    \dim_{\kappa(\mathfrak{p})}P\otimes_R\kappa(\mathfrak{p})
  \]
  is locally constant~\cite[Tag~00NX, Lemma~10.78.2]{Stacks}. The space
  $\operatorname{Spec}R$ has finitely many irreducible
  components~\cite[Tag~00FR, Lemma~10.31.6]{Stacks}; each of them is connected
  and contained in a connected component, so $\operatorname{Spec}R$ has
  finitely many connected components $Z_1,\dots,Z_c$, and $\rho_P$ takes a
  constant value $\rho_P(Z_i)$ on each of them. A finite-dimensional module $Y$
  has finite length, its composition factors are of the form $R/\mathfrak{m}$
  with $\mathfrak{m}$ maximal and $R/\mathfrak{m}$ finite-dimensional, and
  \[
    \dim_\kk(P\otimes_RR/\mathfrak{m})=\rho_P(\mathfrak{m})\dim_\kk
    R/\mathfrak{m}.
  \]
  Since ${P\otimes_R-}$ is exact,
  \[
    \chi_Y([P])=\sum_{i=1}^c\rho_P(Z_i)\,d_i(Y),
  \]
  where $d_i(Y)$ is the sum of the dimensions of the composition factors
  $R/\mathfrak{m}$ of $Y$ with ${\mathfrak{m}\in Z_i}$. Hence, the image of
  $\chi$ lies in the span of the $c$ functions $d_i$.
\end{proof}

In case~\eqref{it:rank-noetherian}, the argument bounds the dimension of the
image of $\chi$ by the number of connected components, without requiring
$K_0(R)$ to have finite rank. In the example below, $\mathcal V$ is
infinite-dimensional; see \Cref{rem:selection-visible-rank}.

\subsection{The group}
\label{subsec:group}

We write ${[x,y]=xyx^{-1}y^{-1}}$ for elements of a group, and put
${I=\set{1,2,3}}$. Let $\varepsilon_1,\varepsilon_2,\varepsilon_3$ be the
standard basis of $\ZZ^3$, and for ${n=(n_1,n_2,n_3)\in\ZZ^3}$ and commuting
elements $T_1,T_2,T_3$ write ${T^n=T_1^{n_1}T_2^{n_2}T_3^{n_3}}$.

\begin{definition}
  \label{def:selection-group}
  Let $G$ be the group with generators $T_1$, $T_2$, $T_3$ and $U_i(r)$,
  $V_i(r)$, $W_{ij}(r)$ for ${i,j\in I}$ with ${i\neq j}$ and ${r\in\ZZ}$, subject
  to the following relations, for all integers $r$ and $s$ and all indices
  satisfying the stated conditions: the $T_\ell$ commute pairwise;
  \begin{equation}
    \label{eq:group-conjugation}
    T_\ell S(r)T_\ell^{-1}=S(r+\lambda_S(\varepsilon_\ell))
  \end{equation}
  for every generator type ${S\in\set{U_i,V_i,W_{ij}}}$, where the
  \emph{weights} are the homomorphisms ${\lambda_S\colon\ZZ^3\to\ZZ}$ given by
  ${\lambda_{U_i}(n)=-n_i}$, ${\lambda_{V_i}(n)=n_i}$ and
  ${\lambda_{W_{ij}}(n)=n_i-n_j}$;
  \begin{equation}
    \label{eq:group-relations}
    \begin{gathered}
      U_i(r)^2=1,\qquad
      [U_i(r),U_j(s)]=[V_i(r),V_j(s)]=[U_i(r),V_j(s)]=1\quad(i\neq j),\\
      [W_{ij}(s),U_h(r)]=1\quad(h\neq i),\qquad
      [W_{ij}(s),V_h(r)]=1\quad(h\neq j),\\
      [U_i(r),W_{ij}(s)]=U_j(r+s),\qquad[W_{ij}(s),V_j(r)]=V_i(r+s).
    \end{gathered}
  \end{equation}
\end{definition}

No relation between two generators of the same type with different parameters
and no relation between two generators of types $W_{ij}$ is imposed. The main
preprint describes $G$ as motivated by finitely presented matrix groups of
Abels type, and its commutator arguments parallel the proof
of~\cite[Proposition~4.9, relations (4.9) and (4.10)]{San22}; the proofs below
are self-contained.

\begin{proposition}
  \label{prop:group-presentation}
  The group $G$ is finitely presented. More precisely, put ${u_i=U_i(0)}$,
  ${v_i=V_i(0)}$ and ${w_{ij}=W_{ij}(0)}$. Then $G$ is generated by the fifteen
  elements $T_\ell$, $u_i$, $v_i$ and $w_{ij}$, subject to the following
  finitely many relations: the $T_\ell$ commute pairwise; the relations
  \eqref{eq:group-relations} with ${r=s=0}$; and ${[T_h,u_i]=[T_h,v_i]=1}$ for
  ${h\neq i}$, together with ${[T_iT_j,w_{ij}]=[T_k,w_{ij}]=1}$ for
  ${\set{i,j,k}=I}$.
\end{proposition}

\begin{proof}
  Let $P$ be the group with the stated presentation. For a type $S$ write $s_S$
  for the generator $u_i$, $v_i$ or $w_{ij}$, and choose
  ${q_{U_i}=-\varepsilon_i}$, ${q_{V_i}=\varepsilon_i}$ and
  ${q_{W_{ij}}=\varepsilon_i}$, so that ${\lambda_S(q_S)=1}$. The vectors $b$ with
  ${[T^b,s_S]=1}$ among the relations of $P$ form a basis of the kernel $K_S$ of
  $\lambda_S$: for $U_i$ and $V_i$ the kernel is ${\set{n}[n_i=0]}$, with basis
  $\varepsilon_h$ for ${h\neq i}$, and for $W_{ij}$ it is ${\set{n}[n_i=n_j]}$,
  whose elements are ${n_i(\varepsilon_i+\varepsilon_j)+n_k\varepsilon_k}$.
  Hence, $T^b$ commutes with $s_S$ for every ${b\in K_S}$. In $P$ put
  \[
    \widetilde{S}(r)=T^{rq_S}s_ST^{-rq_S}\qquad(r\in\ZZ).
  \]
  For ${n\in\ZZ^3}$ we have ${n-\lambda_S(n)q_S\in K_S}$, and the elements
  $T_1,T_2,T_3$ commute, so
  \begin{equation}
    \label{eq:torus-conjugation}
    T^n\widetilde{S}(r)T^{-n}=\widetilde{S}(r+\lambda_S(n)).
  \end{equation}
  This gives the relations~\eqref{eq:group-conjugation} for the elements
  $\widetilde{S}(r)$, and conjugating ${u_i^2=1}$ by $T^{rq_{U_i}}$ gives
  ${\widetilde{U}_i(r)^2=1}$. Each of the remaining relations
  in~\eqref{eq:group-relations} involves two generators of types $S$ and $S'$
  with prescribed parameters, and it is obtained by conjugating the
  corresponding relation with parameters zero by $T^n$, for an ${n\in\ZZ^3}$
  with the prescribed values of $\lambda_S(n)$ and $\lambda_{S'}(n)$, and using
  \eqref{eq:torus-conjugation}. Such vectors $n$ are, for distinct $i,j,k$:
  \[
    \begin{array}{ll@{\qquad}ll}
      [U_i(r),U_j(s)]\colon&-r\varepsilon_i-s\varepsilon_j, &
      [W_{ij}(s),V_i(r)]\colon&r\varepsilon_i+(r-s)\varepsilon_j,\\{}
      [V_i(r),V_j(s)]\colon&r\varepsilon_i+s\varepsilon_j, &
      [W_{ij}(s),V_k(r)]\colon&s\varepsilon_i+r\varepsilon_k,\\{}
      [U_i(r),V_j(s)]\colon&-r\varepsilon_i+s\varepsilon_j, &
      [U_i(r),W_{ij}(s)]\colon&-r\varepsilon_i-(r+s)\varepsilon_j,\\{}
      [W_{ij}(s),U_j(r)]\colon&(s-r)\varepsilon_i-r\varepsilon_j, &
      [W_{ij}(s),V_j(r)]\colon&(r+s)\varepsilon_i+r\varepsilon_j.\\{}
      [W_{ij}(s),U_k(r)]\colon&s\varepsilon_i-r\varepsilon_k, &&
    \end{array}
  \]
  In the last two
  cases the weight of the third generator is ${\lambda_{U_j}(n)=r+s}$ and
  ${\lambda_{V_i}(n)=r+s}$, as required. Hence, all defining relations of $G$
  hold for the elements $T_\ell$ and $\widetilde{S}(r)$ of $P$.

  Conversely, the relations~\eqref{eq:group-conjugation} give
  ${T^nS(r)T^{-n}=S(r+\lambda_S(n))}$ in $G$ for all ${n\in\ZZ^3}$, which implies
  the relations ${[T^b,S(0)]=1}$ for ${b\in K_S}$; the other relations of $P$ are
  instances of relations of $G$. Thus, there are homomorphisms ${P\to G}$ sending
  $s_S$ to $S(0)$ and ${G\to P}$ sending $S(r)$ to $\widetilde{S}(r)$, both
  fixing the $T_\ell$. Their composites fix all generators, since
  ${\widetilde{S}(0)=s_S}$ and ${T^{rq_S}S(0)T^{-rq_S}=S(r)}$.
\end{proof}

\begin{proposition}
  \label{prop:central-involutions}
  For every ${N\in\ZZ}$ the element
  \[
    z_N=[U_i(a),V_i(b)]\qquad(i\in I,\ a,b\in\ZZ,\ a+b=N)
  \]
  of $G$ does not depend on the choice of $i$, $a$ and $b$. It is central and
  satisfies ${z_N^2=1}$.
\end{proposition}

\begin{proof}
  Let ${i\neq j}$ and ${a,s,b\in\ZZ}$, and put ${x=U_i(a)}$, ${w=W_{ij}(s)}$,
  ${v=V_j(b)}$, ${p}={[x,w]}={U_j(a+s)}$ and ${q}={[w,v]}={V_i(s+b)}$. By the relations
  \eqref{eq:group-relations}, the element $x$ commutes with $v$, and $q$
  commutes with $p$ and with $v$. Using ${xwx^{-1}=pw}$ and ${wvw^{-1}=qv}$, we
  compute
  \[
    xqx^{-1}=x(wvw^{-1}v^{-1})x^{-1}=(pw)v(pw)^{-1}v^{-1}
    =p(qv)p^{-1}v^{-1}=q[p,v].
  \]
  Since $q$ commutes with $p$ and $v$, it commutes with $[p,v]$, and
  multiplying by $q^{-1}$ gives ${[x,q]=[p,v]}$, that is,
  \begin{equation}
    \label{eq:commutator-transfer}
    [U_i(a),V_i(s+b)]=[U_j(a+s),V_j(b)].
  \end{equation}
  If ${a+b=c+d}$ and ${i\neq j}$, then~\eqref{eq:commutator-transfer} with
  ${s=c-a}$ and $b$ replaced by $d$ gives ${[U_i(a),V_i(b)]=[U_j(c),V_j(d)]}$.
  Two splittings at the same index $i$ are both equal to $[U_j(N),V_j(0)]$ for
  some ${j\neq i}$. This proves the independence.

  Let $g$ be a generator of $G$. If $g$ is $U_h(r)$ or $V_h(r)$, write $z_N$ at
  an index ${i\neq h}$; then $g$ commutes with $U_i(a)$ and $V_i(b)$
  by~\eqref{eq:group-relations}. If ${g=W_{ij}(r)}$, write $z_N$ at the index
  ${k\notin\set{i,j}}$, and use the relations ${[W_{ij},U_k]=[W_{ij},V_k]=1}$. If
  ${g=T_\ell}$, then~\eqref{eq:group-conjugation} gives
  \[
    T_\ell z_NT_\ell^{-1}=[U_i(a-\delta_{\ell i}),V_i(b+\delta_{\ell i})]=z_N.
  \]
  Hence, $z_N$ is central. Finally, with ${x=U_i(a)}$ and ${y=V_i(b)}$, centrality
  of ${z_N=[x,y]}$ gives ${xyx^{-1}=z_Ny}$ and then ${x^2yx^{-2}=z_N^2y}$; since
  ${x^2=1}$, we obtain ${z_N^2=1}$.
\end{proof}

\begin{proposition}
  \label{prop:shift-automorphism}
  For every ${c\in\ZZ}$ there is an automorphism $\beta_c$ of $G$ fixing the
  $T_\ell$, the $V_i(r)$ and the $W_{ij}(r)$, with ${\beta_c(U_i(r))=U_i(r+c)}$.
  Moreover, ${\beta_c\beta_d=\beta_{c+d}}$, and ${\alpha_G\coloneqq\beta_1}$
  satisfies ${\alpha_G(z_N)=z_{N+1}}$ for every $N$.
\end{proposition}

\begin{proof}
  The assignment sends each defining relation to a defining relation:
  \begin{itemize}
    \item the relations~\eqref{eq:group-conjugation} for $U_i$ go to
      \[
        T_\ell U_i(r+c)T_\ell^{-1}=U_i(r+c-\delta_{\ell i});
      \]
    \item the relations
      \begin{gather*}
        U_i(r)^2=1,\qquad[U_i(r),U_j(s)]=1,\\
        [U_i(r),V_j(s)]=1,\qquad[W_{ij}(s),U_h(r)]=1
      \end{gather*}
      go to the instances with ${r+c}$ and ${s+c}$ in place of the parameters of
      $U$;
    \item the relation ${[U_i(r),W_{ij}(s)]=U_j(r+s)}$ goes to its instance with
      first parameter ${r+c}$;
    \item all other relations are fixed.
  \end{itemize}
  Hence,
  $\beta_c$ is an endomorphism, and the identities ${\beta_c\beta_d=\beta_{c+d}}$
  and ${\beta_0=\id}$ hold on generators, so $\beta_c$ is an automorphism with
  inverse $\beta_{-c}$. Applying $\beta_1$ to ${z_N=[U_i(a),V_i(b)]}$ gives
  ${[U_i(a+1),V_i(b)]=z_{N+1}}$ by \Cref{prop:central-involutions}.
\end{proof}

\subsection{Finite quotients}
\label{subsec:finite-quotients}

For an integer ${m\geq1}$ let ${S_m=\mathbb{F}_2[t]/(t^m-1)}$. Division by the
monic polynomial ${t^m-1}$ shows that $1,t,\dots,t^{m-1}$ is a basis of $S_m$
over $\mathbb{F}_2$, and $t$ is a unit with inverse $t^{m-1}$; the ring $S_m$
need not be reduced. We write $E_{ab}$ for the matrix units in $M_5(S_m)$ and
$\mathbf{1}$ for the identity matrix.

\begin{proposition}
  \label{prop:finite-quotients}
  For every ${m\geq1}$ there is a homomorphism
  ${\pi_m\colon G\to\operatorname{GL}_5(S_m)}$ with
  \[
    \begin{gathered}
      \pi_m(T_\ell)=D_\ell,\qquad\pi_m(U_i(r))=\mathbf{1}+t^rE_{1,i+1},\\
      \pi_m(V_i(r))=\mathbf{1}+t^rE_{i+1,5},\qquad
      \pi_m(W_{ij}(r))=\mathbf{1}+t^rE_{i+1,j+1},
    \end{gathered}
  \]
  where $D_\ell$ is the diagonal matrix with entry $t$ in position ${\ell+1}$ and
  entries $1$ elsewhere. Its image $F_m$ is finite, and
  ${\pi_m(z_N)=\mathbf{1}+t^NE_{15}}$ for every ${N\in\ZZ}$. The elements
  $\pi_m(z_0),\dots,\pi_m(z_{m-1})$ generate a central subgroup $Z_m$ of $F_m$
  isomorphic to $(\ZZ/2)^m$.
\end{proposition}

\begin{proof}
  Every matrix unit ${A=E_{ab}}$ with ${a\neq b}$ satisfies ${A^2=0}$, so
  ${(\mathbf{1}+xA)^2}={\mathbf{1}+2xA}={\mathbf{1}}$ for ${x\in S_m}$; in particular,
  the matrices above are invertible and the relations ${U_i(r)^2=1}$ hold. The
  diagonal matrices commute, and ${DE_{ab}D^{-1}}={d_ad_b^{-1}E_{ab}}$ for a
  diagonal matrix $D$ with entries $d_a$; for ${D=D_\ell}$ the factors for
  $E_{1,i+1}$, $E_{i+1,5}$ and $E_{i+1,j+1}$ are $t^{-\delta_{\ell i}}$,
  $t^{\delta_{\ell i}}$ and $t^{\delta_{\ell i}-\delta_{\ell j}}$, which gives
  the relations~\eqref{eq:group-conjugation}. If $A$ and $B$ are matrix units
  with ${AB=BA=0}$, then ${\mathbf{1}+xA}$ and ${\mathbf{1}+yB}$ commute; since
  ${E_{ab}E_{cd}=\delta_{bc}E_{ad}}$, this applies to the pairs in the
  commutation relations of~\eqref{eq:group-relations}, by the restrictions on
  the indices: for instance, for $W_{ij}$ and $U_h$ with ${h\neq i}$ the
  products are ${E_{i+1,j+1}E_{1,h+1}=0}$ and
  ${E_{1,h+1}E_{i+1,j+1}}={\delta_{hi}E_{1,j+1}}={0}$. For distinct $a$, $b$ and $c$,
  the matrices ${A=E_{ab}}$ and ${B=E_{bc}}$ satisfy ${BA=0}$ and ${AB=E_{ac}}$, and
  $AB$ has zero product with $A$ and $B$ on both sides; hence
  \[
    [\mathbf{1}+xA,\mathbf{1}+yB]=(\mathbf{1}+xA)(\mathbf{1}+yB)
    (\mathbf{1}+xA)(\mathbf{1}+yB)=\mathbf{1}+xyAB.
  \]
  With ${(a,b,c)=(1,i+1,j+1)}$ and ${(i+1,j+1,5)}$ this gives the last two
  relations of~\eqref{eq:group-relations}, and with ${(1,i+1,5)}$ it gives
  ${\pi_m(z_N)=\mathbf{1}+t^{a+b}E_{15}=\mathbf{1}+t^NE_{15}}$. The image $F_m$ is
  finite since $M_5(S_m)$ is finite.

  The matrix unit $E_{15}$ has zero product on both sides with every matrix unit
  $E_{ab}$, ${a\neq b}$, occurring above, and every $D_\ell$ has equal entries
  in positions $1$ and $5$, so $\pi_m(z_N)$ is central in $F_m$. Since
  ${E_{15}^2=0}$, we have
  \[
    \prod_{q=0}^{m-1}\pi_m(z_q)^{c_q}
    =\mathbf{1}+\Bigl(\sum_{q=0}^{m-1}c_qt^q\Bigr)E_{15}
  \]
  for ${c_q\in\set{0,1}}$, which is the identity only if all $c_q$ vanish.
  Together with ${\pi_m(z_q)^2=\mathbf{1}}$ this gives
  ${Z_m\cong(\ZZ/2)^m}$.
\end{proof}

\begin{corollary}
  \label{coro:z-independent}
  The elements $z_N$, ${N\in\ZZ}$, generate a subgroup of the centre of $G$
  isomorphic to $\bigoplus_{N\in\ZZ}\ZZ/2$. In particular, the centre of $G$ is
  not finitely generated.
\end{corollary}

\begin{proof}
  Let ${N_1<\dots<N_a}$ be integers and choose ${m>N_a-N_1}$. The residues of the
  $N_b$ modulo $m$ are distinct, so \Cref{prop:finite-quotients} shows that
  every nonempty product of the $z_{N_b}$ has nontrivial image in $F_m$. Since
  the $z_N$ are central involutions by \Cref{prop:central-involutions}, they
  generate the stated subgroup. A subgroup of a finitely generated abelian
  group is finitely generated, which this subgroup is not.
\end{proof}

The group $G$ is finitely presented, has an automorphism shifting the
independent central involutions, and has finite quotients $F_m$ separating every
finite subfamily of them.

\subsection{Selection data with unbounded extinction}
\label{subsec:unbounded-extinction}

Let ${R=\mathbb{C}G}$ be the complex group algebra of the group $G$ of
\Cref{def:selection-group}, let ${\alpha\colon R\to R}$ be the linear extension
of $\alpha_G$, and put
\[
  e=\tfrac12(1+z_0)\in R,\qquad\Psi={}_\alpha(eR).
\]

\begin{lemma}
  \label{lemma:selection-iterates}
  Let $R$ be a $\kk$-algebra, ${e\in R}$ a central idempotent and $\alpha$ an
  automorphism of $R$, and let $H$ be the selection functor of
  ${\Psi={}_\alpha(eR)}$. Put ${e_q=\alpha^q(e)}$ and ${p_j=e_0e_1\cdots e_{j-1}}$,
  with ${p_0=1}$. Then for every finite-dimensional left $R$-module $Y$ and every
  ${j\geq0}$ there is an isomorphism
  \[
    H^jY\cong{}_{\alpha^j}(p_jY),
  \]
  natural in $Y$, where ${r\in R}$ acts on the subspace $p_jY$ of $Y$ as
  $\alpha^j(r)$.
\end{lemma}

\begin{proof}
  The elements $e_q$ are central idempotents, since $\alpha$ is an
  automorphism. For ${j=0}$ there is nothing to prove. If
  ${H^jY\cong{}_{\alpha^j}(p_jY)}$, then by \Cref{ex:twisted-corner} the module
  $H^{j+1}Y$ is the image of $e$ acting on ${}_{\alpha^j}(p_jY)$, that is, the
  subspace ${\alpha^j(e)p_jY=p_{j+1}Y}$, with $r$ acting as
  ${\alpha^j(\alpha(r))=\alpha^{j+1}(r)}$. The identifications are restrictions
  of the identity of $Y$, hence natural.
\end{proof}

\begin{theorem}
  \label{thm:selection-unbounded}
  The pair $(R,\Psi)$ is selection data over $\mathbb{C}$ with $R$ finitely
  presented. For every integer ${m\geq1}$ there is a nonzero finite-dimensional
  left $R$-module $Y_m$ of extinction time exactly $m$. In particular,
  $(R,\Psi)$ has unbounded extinction.
\end{theorem}

\begin{proof}
  By \Cref{prop:group-presentation}, the group $G$ has a presentation with
  fifteen generators $g_1,\dots,g_{15}$ and finitely many relators
  $w_1,\dots,w_b$. The algebra with generators $x_1,y_1,\dots,x_{15},y_{15}$
  and relations ${x_iy_i=y_ix_i=1}$ and ${w_c(x,y)=1}$, in which $g_i^{-1}$ is
  replaced by $y_i$, is isomorphic to $R$: the homomorphism to $R$ has
  an inverse given by the linear extension of the group homomorphism from $G$
  to its group of units sending $g_i$ to $x_i$. Hence, $R$ is finitely
  presented. The map $\alpha$ is an automorphism with inverse the linear
  extension of $\alpha_G^{-1}$, and $e$ is a central idempotent by
  \Cref{prop:central-involutions}; so $(R,\Psi)$ is selection data, by
  \Cref{ex:twisted-corner}. Moreover,
  ${e_q=\alpha^q(e)=\frac12(1+z_q)}$ by \Cref{prop:shift-automorphism}.

  Fix ${m\geq1}$, and let $F_m$ and $Z_m$ be as in
  \Cref{prop:finite-quotients}, with ${\overline{z}_q=\pi_m(z_q)}$. Since
  ${Z_m\cong(\ZZ/2)^m}$ with basis $\overline{z}_0,\dots,\overline{z}_{m-1}$,
  there is a character ${\chi_m\colon Z_m\to\mathbb{C}^\times}$ with
  ${\chi_m(\overline{z}_q)=1}$ for ${0\leq q<m-1}$ and
  ${\chi_m(\overline{z}_{m-1})=-1}$.
  Put
  \[
    p_{\chi_m}=\frac{1}{2^m}\sum_{z\in Z_m}\chi_m(z)^{-1}z\in\mathbb{C}F_m,
    \qquad Y_m=(\mathbb{C}F_m)\,p_{\chi_m},
  \]
  regarded as a left $R$-module through $\pi_m$. The coefficient of the
  identity in $p_{\chi_m}$ is $2^{-m}$, so ${Y_m\neq0}$, and $Y_m$ is
  finite-dimensional. For ${h\in Z_m}$, the substitution ${u=hz}$ gives
  ${hp_{\chi_m}=\chi_m(h)p_{\chi_m}}$, and since $Z_m$ is central in $F_m$,
  every ${h\in Z_m}$ acts on $Y_m$ as the scalar $\chi_m(h)$. Hence, $e_q$ acts on
  $Y_m$ as the identity for ${0\leq q<m-1}$ and as zero for ${q=m-1}$. By
  \Cref{lemma:selection-iterates}, $H^jY_m$ is the subspace $p_jY_m$, which is
  ${Y_m\neq0}$ for ${j<m}$ and zero for ${j\geq m}$, since then $p_j$ contains the
  factor $e_{m-1}$.
\end{proof}

\begin{remark}
  \label{rem:selection-visible-rank}
  In the notation of \Cref{prop:rank-obstruction}, we have
  ${\chi_Y(\mathsf{T}^j[R])=\dim p_jY}$ by \Cref{lemma:selection-iterates},
  and the values of these functions on the
  modules $Y_1,\dots,Y_s$ form a triangular matrix with nonzero diagonal, so
  $\mathcal V$ is infinite-dimensional. The module $Y_m$ is the module
  induced from the character $\chi_m$ of $Z_m$, of dimension $[F_m:Z_m]$. The
  integer $m$ enters only through the
  choice of the finite quotient $F_m$ and of the values of $\chi_m$ on the
  central involutions; $R$, $e$ and $\alpha$ are independent of $m$, as required
  in \Cref{sec:simulation}.
\end{remark}

\section{Bimodule complexes from selection data}
\label{sec:realisation}

This section turns selection data into a complex of bimodules over a
finite-dimensional algebra of global dimension at most two. Let $(R,\Psi)$ be
selection data in the sense of \Cref{def:selection-data}, with $R$ finitely
presented, and put ${H=\Psi\otimes_R-}$. \Cref{thm:realisation} provides a
finite-dimensional algebra $B$, an exact fully faithful functor $M$ from
finite-dimensional $R$-modules to $B$-modules, and a bounded complex $P$ of
$B$-bimodules such that
\[
  P\Lotimes[B]M(Y)\simeq M(HY)\oplus M(HY)[3]
\]
for every finite-dimensional $R$-module $Y$. The argument has four steps.
\begin{enumerate}
  \item The algebra $B$ is constructed from a presentation of $R$ by generators
    and relations of degree at most two, and $R$ acts on an object $E_0$ of a
    Verdier quotient $\mathcal{Q}$ of the bounded homotopy category of
    projective right $B$-modules (\Cref{subsec:encoding,subsec:action}).
  \item The bimodule $\Psi$ determines a direct summand $U$ of a finite sum of
    copies of $E_0$ in the idempotent completion of $\mathcal{Q}$. The object
    $U$ need not lie in $\mathcal{Q}$, but ${U\oplus U[3]}$ does
    (\Cref{subsec:odd-double}).
  \item A representation of the quiver with relations of $B$ in
    $\mathcal{Q}$, with ${U\oplus U[3]}$ at every vertex, lifts to the homotopy
    category up to homotopies prescribed on the relations
    (\Cref{subsec:lifting}).
  \item Such a homotopy representation is replaced by a complex of
    bimodules (\Cref{subsec:rectification}), and the extension between the
    two cohomology modules of ${P\Lotimes[B]M(Y)}$ vanishes because
    ${\operatorname{gldim}B\leq2}$ (\Cref{subsec:realisation-theorem}).
\end{enumerate}
The argument uses choices that are not made explicit in steps~(2) and~(3); see
\Cref{rem:uncontrolled-length}.

Throughout this section $\kk$ is a field, and all algebras, bimodules,
categories and functors are $\kk$-linear; bimodules are $\kk$-central, that is,
the two actions of $\kk$ on them agree. For a finite-dimensional algebra $B$ we
write $\operatorname{proj}B^{\op}$ for the category of finitely generated
projective right $B$-modules and $\operatorname{mod}B$ for the category of
finitely generated left $B$-modules. For a $\kk$-algebra $R$, not necessarily
finite-dimensional, $\operatorname{mod}_{\mathrm{fd}}R$ denotes the category of
finite-dimensional left $R$-modules.

\subsection{Right fractions}
\label{subsec:fractions}

Let $\mathcal{K}$ be an essentially small triangulated category and
$\mathcal{S}$ a full triangulated subcategory. Let $\mathcal{W}$ be the class of
morphisms $u$ of
$\mathcal{K}$ that occur in a triangle ${X'\xrightarrow{u}X\to Z\to X'[1]}$ with
$Z$ isomorphic to an object of $\mathcal{S}$. By~\cite[Tag~05RG,
Lemma~13.6.6]{Stacks}, $\mathcal{W}$ is a multiplicative system, and the
Verdier quotient is defined as
${\mathcal{K}/\mathcal{S}=\mathcal{W}^{-1}\mathcal{K}}$, with quotient functor
${q\colon\mathcal{K}\to\mathcal{K}/\mathcal{S}}$~\cite[Tag~05RI,
Definition~13.6.7]{Stacks}. Since the categories of left and of right
fractions of a multiplicative system are canonically
isomorphic~\cite[Tag~04VL, Lemma~4.27.19]{Stacks}, we may compute in
$\mathcal{K}/\mathcal{S}$ with right fractions~\cite[Tag~04VH,
Lemma~4.27.11]{Stacks}: a morphism ${qX\to qY}$ is the class of a pair $(f,u)$ of
morphisms ${f\colon X'\to Y}$ and ${u\colon X'\to X}$ with ${u\in\mathcal{W}}$. The
composition rule gives ${(f,u)\circ q(u)=q(f)}$, so that the class of $(f,u)$ is
$q(f)q(u)^{-1}$, and $q(u)$ is invertible for every
${u\in\mathcal{W}}$~\cite[Tag~04VK, Lemma~4.27.16]{Stacks}.

\begin{lemma}
  \label{lemma:fractions}
  In the situation above, let $X$ be an object of $\mathcal{K}$.
  \begin{enumerate}
    \item\label{it:fractions-denominator} For finitely many morphisms
      ${\gamma_i\colon qX\to qY_i}$ in $\mathcal{K}/\mathcal{S}$, there exist a
      morphism ${u\colon X'\to X}$ in $\mathcal{W}$ and morphisms
      ${f_i\colon X'\to Y_i}$ in $\mathcal{K}$ such that
      ${\gamma_i=q(f_i)q(u)^{-1}}$ for every $i$.
    \item\label{it:fractions-kill} For finitely many morphisms
      ${g_i\colon X\to Y_i}$ in $\mathcal{K}$ with ${q(g_i)=0}$, there exists a
      morphism ${v\colon X'\to X}$ in $\mathcal{W}$ such that ${g_iv=0}$ for every
      $i$.
  \end{enumerate}
\end{lemma}

\begin{proof}
  \eqref{it:fractions-denominator} By induction on the number of morphisms, it
  suffices to give two roofs $(f_1,u_1)$ and $(f_2,u_2)$ from $X$, with
  ${u_i\colon X_i\to X}$, a common denominator. The Ore condition for the right
  multiplicative system $\mathcal{W}$~\cite[Tag~04VC,
  Definition~4.27.1]{Stacks} gives morphisms ${t\colon X'\to X_2}$ in
  $\mathcal{W}$ and ${g\colon X'\to X_1}$ with ${u_1g=u_2t}$. The morphism
  ${u\coloneqq u_2t}$ lies in $\mathcal{W}$, and ${q(g)=q(u_1)^{-1}q(u)}$ is
  invertible. Hence,
  \[
    q(f_1)q(u_1)^{-1}=q(f_1g)q(u)^{-1}\qquad\text{and}\qquad
    q(f_2)q(u_2)^{-1}=q(f_2t)q(u)^{-1}.
  \]
  For the empty family take ${u=\id[X]}$.

  \eqref{it:fractions-kill} Replacing the family by the single morphism
  ${g=(g_i)_i\colon X\to\bigoplus_iY_i}$, we may assume that there is one
  morphism. The roofs $(g,\id[X])$ and $(0,\id[X])$ represent the same morphism
  of $\mathcal{K}/\mathcal{S}$, so that~\cite[Tag~04VJ, Lemma~4.27.14]{Stacks}
  gives ${v\colon X'\to X}$ in $\mathcal{W}$ with ${gv=0}$.
\end{proof}

\subsection{The algebra \texorpdfstring{$B$}{B}}
\label{subsec:encoding}

Every finitely presented algebra has a presentation by relations of degree at
most two. Indeed, given a finite presentation, introduce for each monomial
$y_1y_2\cdots y_n$ of length ${n\geq3}$ occurring in a relation new generators
$w_2,\dots,w_{n-1}$ with the relations ${w_2=y_1y_2}$ and ${w_{j+1}=w_jy_{j+1}}$,
and replace the monomial by $w_{n-1}y_n$. Sending $w_j$ to $y_1\cdots y_j$ and
each old generator to its class defines mutually inverse homomorphisms between
the two presented algebras.

\begin{construction}
  \label{cons:encoding-algebra}
  Let $R$ be a $\kk$-algebra with a presentation
  \[
    R=\kk\langle x_1,\dots,x_d\rangle/(p_1,\dots,p_t)
  \]
  in which every $p_j$ is a noncommutative polynomial of degree at most two;
  constant and linear terms are allowed. Let $Q$ be the quiver with vertices
  $0$, $1$ and~$2$, arrows ${s,a_1,\dots,a_d\colon0\to1}$ and
  ${s',a'_1,\dots,a'_d\colon1\to2}$, and no other arrows. Write
  $\mathcal{A}_{01}$ and $\mathcal{A}_{12}$ for the two sets of arrows. The
  \emph{homogenisation} $\widetilde{p}_j$ of $p_j$ is the linear combination of
  paths of length two obtained by the substitutions
  \[
    x_hx_l\longmapsto a'_ha_l,\qquad x_h\longmapsto s'a_h,\qquad
    1\longmapsto s's.
  \]
  Let $J_2$ be the subspace of the span of the paths of length two spanned by
  the elements ${a'_hs-s'a_h}$ for ${1\leq h\leq d}$ and $\widetilde{p}_j$ for
  ${1\leq j\leq t}$, and let $\mathcal{R}$ be a basis of $J_2$. We write
  \[
    \rho=\sum_{a\in\mathcal{A}_{01},\,b\in\mathcal{A}_{12}}\rho_{ba}\,ba
    \qquad(\rho\in\mathcal{R},\ \rho_{ba}\in\kk).
  \]
  Put ${B=\kk Q/(J_2)}$, with
  vertex idempotents $e_0$, $e_1$ and $e_2$. For a finite-dimensional left
  $R$-module $Y$, let $M(Y)$ be the representation of $Q$ with the space $Y$
  at each vertex, the identity of $Y$ on $s$ and $s'$, and the action of $x_h$
  on $a_h$ and~$a'_h$.
\end{construction}

The algebra $B$ depends on the presentation. Its defining relations span $J_2$.
In the quotient category of \Cref{subsec:action}, $s$ and $s'$ become invertible
and $x_h$ acts as $s^{-1}a_h$.

\begin{lemma}
  \label{lemma:directed-gldim}
  Let $\Lambda$ be a finite-dimensional algebra with a complete set of
  orthogonal idempotents $\varepsilon_0,\dots,\varepsilon_{n-1}$, where
  ${n\geq1}$, such that ${\varepsilon_i\Lambda\varepsilon_i=\kk\varepsilon_i}$ for
  every $i$ and ${\varepsilon_j\Lambda\varepsilon_i=0}$ for ${j<i}$. Then every left
  and every right $\Lambda$-module has projective dimension at most ${n-1}$.
\end{lemma}

\begin{proof}
  Let $N$ be a left $\Lambda$-module with ${\varepsilon_iN=0}$ for ${i<p}$. The map
  \[
    \bigoplus_{i=p}^{n-1}\Lambda\varepsilon_i\otimes_\kk\varepsilon_iN
    \longrightarrow N,\qquad c\otimes v\longmapsto cv,
  \]
  is surjective since ${v=\sum_i\varepsilon_iv}$, and its source is projective.
  The hypotheses give ${\varepsilon_p\Lambda\varepsilon_i=0}$ for ${i>p}$ and
  ${\varepsilon_p\Lambda\varepsilon_p=\kk\varepsilon_p}$, so the map is an
  isomorphism after multiplication by $\varepsilon_p$, and its kernel $N'$
  satisfies ${\varepsilon_iN'=0}$ for ${i<p+1}$. A module annihilated by
  $\varepsilon_0,\dots,\varepsilon_{n-2}$ is a direct sum of copies of
  ${\Lambda\varepsilon_{n-1}=\kk\varepsilon_{n-1}}$, hence projective. Starting
  with ${p=0}$ and repeating the construction ${n-1}$ times gives a projective
  resolution of length at most ${n-1}$. The opposite algebra, with the
  idempotents in the reverse order, satisfies the same hypotheses, which gives
  the statement for right modules.
\end{proof}

\begin{proposition}
  \label{prop:encoding-algebra}
  In \Cref{cons:encoding-algebra}, the algebra $B$ is finite-dimensional,
  \[
    \dim_\kk B=3+2(d+1)+(d+1)^2-\dim_\kk J_2,
  \]
  and every left and every right $B$-module has projective dimension at most
  two. Moreover, $M(Y)$ is a left $B$-module for every finite-dimensional left
  $R$-module $Y$, and $M$ extends to an exact fully faithful functor
  ${M\colon\operatorname{mod}_{\mathrm{fd}}R\to\operatorname{mod}B}$ with
  ${e_iM(Y)=Y}$ for every vertex $i$. In particular, ${M(Y)=0}$ only if ${Y=0}$.
\end{proposition}

\begin{proof}
  The quiver has three paths of length zero, ${2(d+1)}$ arrows, ${(d+1)^2}$ paths
  of length two and no longer paths. No arrow ends at~$0$ and no arrow starts
  at~$2$, so the product of an element of $J_2$ with an arrow on either side is
  zero, and the ideal generated by $J_2$ is $J_2$ itself. This gives the
  dimension formula. The idempotents $e_0,e_1,e_2$ satisfy
  ${e_iBe_i=\kk e_i}$ and ${e_jBe_i=0}$ for ${j<i}$, since every arrow increases the
  vertex index; so \Cref{lemma:directed-gldim} applies with ${n=3}$.

  The element ${a'_hs-s'a_h}$ acts on $M(Y)$ as ${x_h-x_h=0}$. The element
  $\widetilde{p}_j$ acts as $p_j(x_1,\dots,x_d)$, which is zero on~$Y$. A
  homomorphism of $R$-modules ${Y\to Y'}$, placed at the three vertices, is a
  homomorphism ${M(Y)\to M(Y')}$. Conversely, a homomorphism ${M(Y)\to M(Y')}$
  consists of three linear maps; compatibility with $s$ and $s'$ forces them
  to be equal, and compatibility with the arrows $a_h$ forces the common map
  to commute with every $x_h$. Hence, $M$ is fully faithful. It is exact
  because exactness of a sequence of $B$-modules is tested at each vertex.
\end{proof}

\subsection{The action on the quotient category}
\label{subsec:action}

In the situation of \Cref{cons:encoding-algebra}, put
\[
  \mathcal{K}\coloneqq K^b(\operatorname{proj}B^{\op}),\qquad
  E_i\coloneqq e_iB\quad(i=0,1,2).
\]
An arrow ${a\colon i\to j}$ of $Q$ lies in $e_jBe_i$, and left multiplication by
$a$ is a homomorphism of right modules ${E_i\to E_j}$, which we denote by $a$ as
well; composition of these homomorphisms is multiplication in~$B$. Let
$\mathcal{S}$ be the smallest thick subcategory of $\mathcal{K}$ containing
the cones of ${s\colon E_0\to E_1}$ and of ${s'\colon E_1\to E_2}$, let
$\mathcal{W}$ be the class of morphisms with cone in $\mathcal{S}$, and let
\[
  q\colon\mathcal{K}\longrightarrow\mathcal{Q}\coloneqq\mathcal{K}/\mathcal{S}
\]
be the Verdier quotient. We write a morphism of $\mathcal{K}$ and its image
under $q$ by the same letter. In $\mathcal{Q}$ the morphisms $s$ and $s'$ are
invertible.

For a finite-dimensional left $R$-module $Y$, the \emph{evaluation} at $Y$ is
the functor
\[
  \operatorname{ev}_Y\colon\mathcal{K}\longrightarrow
  \DerCat[b]{\operatorname{mod}\kk},\qquad T\longmapsto T\otimes_BM(Y).
\]
Termwise tensor product preserves homotopies, mapping cones and shifts, so
$\operatorname{ev}_Y$ is an exact functor. Since a bounded complex of
projective right modules is K-flat, $\operatorname{ev}_Y(T)$ computes
${T\Lotimes[B]M(Y)}$. Under the isomorphisms ${E_i\otimes_BM(Y)}\cong{ e_iM(Y)}={Y}$,
given by ${c\otimes v\mapsto cv}$, the morphism $\operatorname{ev}_Y(a)$ is the
action of the arrow $a$ on $M(Y)$.

\begin{proposition}
  \label{prop:quotient-action}
  In the situation above, there is a unique homomorphism of unital algebras
  \[
    \theta\colon R\longrightarrow\operatorname{End}_{\mathcal{Q}}(E_0)
    \qquad\text{with}\qquad\theta(x_h)=s^{-1}a_h\quad(1\leq h\leq d).
  \]
  For every finite-dimensional left $R$-module $Y$, the functor
  $\operatorname{ev}_Y$ factors uniquely through an exact functor
  ${\operatorname{ev}_Y\colon\mathcal{Q}\to\DerCat[b]{\operatorname{mod}\kk}}$,
  and $\operatorname{ev}_Y(\theta(r))$ is the action of $r$ on
  ${Y=\operatorname{ev}_Y(E_0)}$ for every ${r\in R}$.
\end{proposition}

\begin{proof}
  In $\mathcal{Q}$ the relation ${a'_hs=s'a_h}$ gives ${a'_h=s'a_hs^{-1}}$, and
  therefore
  \begin{gather*}
    (s's)^{-1}a'_ha_l=s^{-1}(s')^{-1}s'a_hs^{-1}a_l=(s^{-1}a_h)(s^{-1}a_l),\\
    (s's)^{-1}s'a_h=s^{-1}a_h,\qquad(s's)^{-1}s's=\id[E_0].
  \end{gather*}
  Multiplying the relation ${\widetilde{p}_j=0}$ on the left by $(s's)^{-1}$
  thus gives
  \[
    p_j(s^{-1}a_1,\dots,s^{-1}a_d)=0,
  \]
  constant term included. Hence,
  the homomorphism from the free algebra sending $x_h$ to $s^{-1}a_h$ factors
  through~$R$; it is unique since the $x_h$ generate $R$.

  The cones of $s$ and $s'$ are sent by $\operatorname{ev}_Y$ to cones of
  identity maps, which are zero. The full subcategory of objects sent to zero
  by the exact functor $\operatorname{ev}_Y$ is closed under shifts, cones
  and direct summands, so it contains $\mathcal{S}$, and the universal property
  of the Verdier quotient~\cite[Tag~05RJ, Lemma~13.6.8]{Stacks} gives the
  exact factorisation. Finally, $\operatorname{ev}_Y(s^{-1}a_h)$ is
  ${\id^{-1}\circ x_h=x_h}$, so ${\operatorname{ev}_Y\circ\theta}$ and the action
  of $R$ on $Y$ agree on generators.
\end{proof}

The argument uses the action $\theta$ and the evaluation functors; it does not
require an identification of $\mathcal Q$ with a derived category of $R$.

\subsection{A summand and its odd shifts}
\label{subsec:odd-double}

Let $\mathcal{T}$ be an additive category. Its \emph{idempotent completion}
$\operatorname{Kar}\mathcal{T}$ has as objects the pairs $(Z,p)$ of an object
$Z$ of $\mathcal{T}$ and an idempotent ${p\in\operatorname{End}(Z)}$; a morphism
${(Z,p)\to(Z',p')}$ is a morphism ${u\colon Z\to Z'}$ with ${u=p'up}$, and the
identity of $(Z,p)$ is $p$. The functor ${Z\mapsto(Z,\id[Z])}$ is fully faithful,
and we regard $\mathcal{T}$ as a full subcategory of
$\operatorname{Kar}\mathcal{T}$. If $\mathcal{T}$ is triangulated, the shift
extends to $\operatorname{Kar}\mathcal{T}$ by ${(Z,p)[1]=(Z[1],p[1])}$. We do not
use a triangulated structure on $\operatorname{Kar}\mathcal{T}$.

\begin{lemma}
  \label{lemma:cone-splitting}
  Let $\mathcal{T}$ be a triangulated category and
  \[
    X_1\xrightarrow{f}X_2\xrightarrow{j}X_3\xrightarrow{v}X_1[1]
  \]
  a triangle in $\mathcal{T}$. Suppose that there are decompositions
  ${X_1\cong I\oplus X_1'}$ and ${X_2}\cong{ I\oplus X_2'}$ in
  $\operatorname{Kar}\mathcal{T}$ in which $f$ has the components $\id[I]$
  on $I$ and zero elsewhere. Then
  ${X_3}\cong{ X_2'\oplus X_1'[1]}$ in $\operatorname{Kar}\mathcal{T}$.
\end{lemma}

\begin{proof}
  Let ${Z=(Z_0,p)}$ be an object of $\operatorname{Kar}\mathcal{T}$. The
  sequence obtained by applying $\Hom{Z}{-}$ to the triangle is the image of
  the idempotent endomorphism $\Hom{p}{-}$ of the long exact sequence obtained
  by applying $\Hom{Z_0}{-}$, and a direct summand of an exact sequence is
  exact. Hence, $\Hom{Z}{-}$ takes the triangle to a long exact sequence.

  Let ${\iota\colon X_2'\to X_2}$ and ${\kappa\colon X_1'[1]\to X_1[1]}$ be the
  inclusions, let ${\pi_2\colon X_2\to X_2'}$ and ${\pi\colon X_1[1]\to X_1'[1]}$
  be the projections, and put ${j'=j\iota}$ and ${v'=\pi v}$. Since ${jf=0}$ and
  ${f[1]v=0}$, the morphism $j$ vanishes on $I$ and the component of $v$ in
  $I[1]$ is zero, so that ${j=j'\pi_2}$ and ${v=\kappa v'}$. We claim that
  \[
    0\longrightarrow\Hom{Z}{X_2'}\xrightarrow{j'_*}\Hom{Z}{X_3}
    \xrightarrow{v'_*}\Hom{Z}{X_1'[1]}\longrightarrow0
  \]
  is exact for every $Z$. If ${j'g=0}$, then ${\iota g=fh}$ for some $h$; the
  component of $fh$ in $X_2'$ is zero, so ${g=0}$. If ${v'g=0}$, then ${vg=0}$, so
  ${g=jh=j'\pi_2h}$ for some $h$. If ${g\colon Z\to X_1'[1]}$, then
  ${f[1]\kappa g=0}$, so ${\kappa g=vh}$ for some $h$ and ${g=\pi\kappa g=v'h}$.

  Taking ${Z=X_1'[1]}$ gives ${\sigma\colon X_1'[1]\to X_3}$ with ${v'\sigma=\id}$.
  The morphism
  \[
    (j',\sigma)\colon X_2'\oplus X_1'[1]\longrightarrow X_3
  \]
  induces a
  bijection on $\Hom{Z}{-}$ for every $Z$: a morphism ${g\colon Z\to X_3}$ is the
  image of $(g',v'g)$, where ${j'g'=g-\sigma v'g}$, and the injectivity follows by
  applying $v'_*$, since ${v'j'=\pi vj\iota=0}$, and then the injectivity of
  $j'_*$. By the Yoneda lemma,
  $(j',\sigma)$ is an isomorphism.
\end{proof}

\begin{proposition}
  \label{prop:odd-double}
  Let $\mathcal{T}$ be an essentially small triangulated category, $E$ an
  object of $\mathcal{T}$ and $p$ an idempotent endomorphism of $E$, and put
  ${U=(E,p)\in\operatorname{Kar}\mathcal{T}}$. Then there exist objects $C$ and
  $V$ of $\mathcal{T}$ such that
  \[
    C\cong U\oplus U[1]\quad\text{and}\quad V\cong U\oplus U[3]
    \quad\text{in }\operatorname{Kar}\mathcal{T},
  \]
  and ${[C]=[V]=0}$ in $K_0(\mathcal{T})$.
\end{proposition}

\begin{proof}
  Let $C$ be a cone of ${\id-p\colon E\to E}$. In the decomposition
  ${E\cong(E,\id-p)\oplus U}$ the morphism ${\id-p}$ is the identity on the first
  summand and zero on the second, so \Cref{lemma:cone-splitting} gives
  ${C\cong U\oplus U[1]}$. Since $\mathcal{T}$ is a full subcategory of
  $\operatorname{Kar}\mathcal{T}$, the morphism
  \[
    g\colon C[1]\cong U[1]\oplus U[2]
    \xrightarrow{\left(\begin{smallmatrix}0&0\\\id&0\end{smallmatrix}\right)}
    U\oplus U[1]\cong C,
  \]
  which identifies the two copies of $U[1]$, is a morphism of $\mathcal{T}$.
  Let $V$ be a cone of~$g$. By \Cref{lemma:cone-splitting}, with ${I=U[1]}$, we
  have ${V\cong U\oplus U[3]}$. The triangles defining $C$ and $V$ give
  ${[C]=[E]-[E]=0}$ and ${[V]=[C]-[C[1]]=2[C]=0}$, where ${[C[1]]=-[C]}$ follows
  from the triangle ${C\to0\to C[1]\xrightarrow{\id}C[1]}$.
\end{proof}

Any object of $\mathcal{T}$ isomorphic to ${U\oplus U[3]}$ in
$\operatorname{Kar}\mathcal{T}$ is isomorphic to $V$ in $\mathcal{T}$, since
the inclusion of $\mathcal{T}$ is fully faithful, and so has class zero as well.
The object $U$ itself need not be isomorphic to an object of $\mathcal{T}$; for
${p=0}$ or ${p=\id}$ it is. The equality ${[V]=0}$ is in $K_0(\mathcal T)$.
Cancellation between summands with shifts of opposite parity also occurs in the
iterates of \Cref{coro:realisation-iterates}, whose classes vanish for ${j\geq1}$,
and in the second iterate of the functor of \Cref{sec:simulation}; see
\Cref{rem:k0-invisibility}.

\begin{remark}
  \label{rem:why-three}
  The cones of \Cref{prop:odd-double} produce ${U\oplus U[1]}$ and
  ${U\oplus U[3]}$. If the diagram of \Cref{cons:diagram-V} were defined using
  ${U\oplus U[c]}$, the
  two cohomology modules of ${P\Lotimes[B]M(Y)}$ in the proof of
  \Cref{thm:realisation} would be in degrees $0$ and ${-c}$, and the extension
  class between them would lie in $\operatorname{Ext}^{c+1}_B$, which vanishes
  when ${c+1>\operatorname{gldim}B}$. Since ${\operatorname{gldim}B\leq2}$, the
  shift $3$ is the smallest positive odd integer $c$ for which this argument
  applies; for ${c=1}$ it would need the vanishing of an
  $\operatorname{Ext}^2_B$, which the bound on the global dimension does not
  provide.
\end{remark}

We now apply \Cref{prop:odd-double} to selection data. Recall from
\Cref{def:selection-data} that $\Psi$ is an $R$-bimodule that is finitely
generated and projective as a right $R$-module, and ${H=\Psi\otimes_R-}$. We
write elements of $R^n$ as columns, so that $M_n(R)$ acts on $R^n$ from the
left and ${\operatorname{End}_{R^{\op}}(R^n)=M_n(R)}$, with composition the
product of matrices.

\begin{lemma}
  \label{lemma:selection-matrix}
  Let $(R,\Psi)$ be selection data. Then there exist an integer ${n\geq1}$, an
  idempotent ${\varepsilon\in M_n(R)}$ and a homomorphism of algebras
  ${\rho\colon R\to\varepsilon M_n(R)\varepsilon}$ with ${\rho(1)=\varepsilon}$
  such that ${\Psi\cong\varepsilon R^n}$ as bimodules, where $R$ acts on
  $\varepsilon R^n$ from the left through $\rho$. For every finite-dimensional
  left $R$-module $Y$ there is an isomorphism, natural in $Y$,
  \[
    HY\cong\varepsilon_Y(Y^n),\quad\text{on which ${r\in R}$ acts as
    $\rho(r)_Y$,}
  \]
  where $m_Y$ denotes the operator on $Y^n$ given by a matrix ${m\in M_n(R)}$.
  In particular, $H$ is exact and ${\dim_\kk HY\leq n\dim_\kk Y}$.
\end{lemma}

\begin{proof}
  Choose a split surjection ${\pi\colon R^n\to\Psi}$ of right modules with
  section $\sigma$, and put ${\varepsilon=\sigma\pi\in M_n(R)}$. Then $\sigma$
  identifies $\Psi$ with $\varepsilon R^n$, and the endomorphisms of the right
  module $\varepsilon R^n$, extended by zero on ${(1-\varepsilon)R^n}$, form
  $\varepsilon M_n(R)\varepsilon$. The left action of $R$ on $\Psi$ thus
  becomes a homomorphism $\rho$ with ${\rho(1)=\varepsilon}$; it is
  $\kk$-linear because $\Psi$ is $\kk$-central. The isomorphism
  ${R^n\otimes_RY\cong Y^n}$, ${(r_i)_i\otimes y\mapsto(r_iy)_i}$, restricts to
  ${\varepsilon R^n\otimes_RY\cong\varepsilon_Y(Y^n)}$, and it carries the left
  action of $r$ to $\rho(r)_Y$. After forgetting the left action of $R$, the
  functor $H$ is a direct summand of the exact functor ${R^n\otimes_R-}$ with
  values in vector spaces; since exactness is detected on underlying vector
  spaces, $H$ is exact.
\end{proof}

Now suppose that $R$ is finitely presented, fix a presentation as in
\Cref{cons:encoding-algebra}, and let $B$, $\mathcal{Q}$ and $\theta$ be as in
\Cref{subsec:action}. Choose $n$, $\varepsilon$ and $\rho$ as in
\Cref{lemma:selection-matrix}. Applying $\theta$ to the entries of a matrix
defines a homomorphism of algebras
\[
  \theta_n\colon
  M_n(R)\longrightarrow\operatorname{End}_{\mathcal{Q}}(E_0^{\oplus n}).
\]
Put
\[
  U\coloneqq\bigl(E_0^{\oplus n},\theta_n(\varepsilon)\bigr)\in
  \operatorname{Kar}\mathcal{Q}.
\]
For ${r\in R}$ the endomorphism $\theta_n(\rho(r))$ of $E_0^{\oplus n}$
satisfies
\[
  \theta_n(\rho(r))=\theta_n(\varepsilon)\theta_n(\rho(r))
  \theta_n(\varepsilon),
\]
so it is an endomorphism of $U$, and
${r\mapsto\theta_n(\rho(r))}$ is a homomorphism of unital algebras
${R\to\operatorname{End}(U)}$, since ${\theta_n(\varepsilon)=\id[U]}$.

By \Cref{prop:odd-double} there is an object $V$ of $\mathcal{Q}$ with
${V\cong U\oplus U[3]}$ in $\operatorname{Kar}\mathcal{Q}$; we fix such an
isomorphism and transport to $V$ the action of $R$ on ${U\oplus U[3]}$ given by
${\theta_n\circ\rho}$ on both summands. Write ${r_V\in\operatorname{End}(V)}$ for
the action of ${r\in R}$; the map ${r\mapsto r_V}$ is a homomorphism of unital
algebras.

\begin{construction}
  \label{cons:diagram-V}
  A \emph{$B$-diagram} in a $\kk$-linear category $\mathcal{C}$ consists of
  objects $W_0$, $W_1$ and $W_2$ and a morphism ${t_a\colon W_i\to W_j}$ for
  every arrow ${a\colon i\to j}$ of $Q$, such that
  ${\sum_{b,a}\rho_{ba}t_bt_a=0}$ for every ${\rho\in\mathcal{R}}$. Let $\mathbb{V}$
  be the $B$-diagram in $\mathcal{Q}$
  with ${W_0=W_1=W_2=V}$, with ${t_s=t_{s'}=\id[V]}$, and with
  ${t_{a_h}=t_{a'_h}=(x_h)_V}$ for ${1\leq h\leq d}$.
\end{construction}

The relations hold in $\mathbb{V}$ since they become ${p_j((x_1)_V,\dots,
(x_d)_V)}={0}$ and ${(x_h)_V-(x_h)_V=0}$. A bounded complex $P$ of $B$-bimodules that
is termwise finitely generated and projective as a right module defines a
$B$-diagram ${(e_iP,a\cdot-)}$ in $\mathcal{K}$, and then in $\mathcal{Q}$. We
lift $\mathbb{V}$ to $\mathcal{K}$ and replace the lift by a complex of
bimodules.

\begin{lemma}
  \label{lemma:evaluate-V}
  For every finite-dimensional left $R$-module $Y$, the functor
  $\operatorname{ev}_Y$ extends to $\operatorname{Kar}\mathcal{Q}$, and there
  is an isomorphism ${\operatorname{ev}_Y(V)\cong HY\oplus HY[3]}$ in
  $\DerCat[b]{\operatorname{mod}\kk}$ under which $\operatorname{ev}_Y(r_V)$ is
  the action of $r$ on $HY$ on both summands, for every ${r\in R}$.
\end{lemma}

\begin{proof}
  Idempotents split in $\DerCat[b]{\operatorname{mod}\kk}$: every bounded
  complex of vector spaces is isomorphic to its cohomology with zero
  differential, morphisms between such complexes have only components of
  degree zero, and an idempotent is then a family of idempotent linear maps,
  whose images are direct summands. Hence, $\operatorname{ev}_Y$ extends to
  $\operatorname{Kar}\mathcal{Q}$, with $(Z,p)$ sent to the image of
  $\operatorname{ev}_Y(p)$. By \Cref{prop:quotient-action}, the functor
  $\operatorname{ev}_Y$ sends $E_0^{\oplus n}$ to $Y^n$ and $\theta_n(m)$ to
  $m_Y$ for ${m\in M_n(R)}$. Hence, it sends $U$ to ${\varepsilon_Y(Y^n)\cong HY}$
  and $\theta_n(\rho(r))$ to the action of $r$, by
  \Cref{lemma:selection-matrix}. Since $\operatorname{ev}_Y$ commutes with the
  shift, the statement follows.
\end{proof}

\subsection{Lifts of \texorpdfstring{$\mathbb{V}$}{V}}
\label{subsec:lifting}

\begin{proposition}
  \label{prop:lifting}
  In the situation of \Cref{cons:diagram-V}, there exist bounded complexes
  $\widetilde{V}_0$, $\widetilde{V}_1$ and $\widetilde{V}_2$ of finitely
  generated projective right $B$-modules, chain maps
  ${f_a\colon\widetilde{V}_i\to\widetilde{V}_j}$ for the arrows ${a\colon i\to j}$
  of $Q$, homogeneous maps ${h_\rho\colon\widetilde{V}_0\to\widetilde{V}_2}$ of
  degree ${-1}$ of right modules for ${\rho\in\mathcal{R}}$, and isomorphisms
  ${\phi_i\colon q\widetilde{V}_i\to V}$ in $\mathcal{Q}$ such that
  \begin{equation}
    \label{eq:lifting}
    \begin{gathered}
      \phi_j\circ q(f_a)=t_a\circ\phi_i\quad\text{for every arrow }
      a\colon i\to j,\\
      dh_\rho+h_\rho d=f_\rho\coloneqq\sum_{b,a}\rho_{ba}f_bf_a
      \quad\text{for every }\rho\in\mathcal{R}.
    \end{gathered}
  \end{equation}
\end{proposition}

\begin{proof}
  The objects of $\mathcal{Q}$ are those of $\mathcal{K}$, so $V$ is an object
  of $\mathcal{K}$ and ${qV=V}$. We proceed from vertex~$2$ to vertex~$0$, so that
  the denominators introduced at a vertex do not affect the arrows already
  lifted.

  Apply \Cref{lemma:fractions}\eqref{it:fractions-denominator} to the
  morphisms ${t_b\colon qV\to qV}$ for ${b\in\mathcal{A}_{12}}$. This gives
  ${u_1\colon\widetilde{V}_1\to V}$ in $\mathcal{W}$ and morphisms
  ${f_b\colon\widetilde{V}_1\to V}$ with ${t_b=q(f_b)q(u_1)^{-1}}$. Put
  ${\widetilde{V}_2=V}$, ${\phi_2=\id[V]}$ and ${\phi_1=q(u_1)}$; then
  ${\phi_2q(f_b)=t_b\phi_1}$. Next apply
  \Cref{lemma:fractions}\eqref{it:fractions-denominator} to the morphisms
  ${\phi_1^{-1}t_a\colon qV\to q\widetilde{V}_1}$ for ${a\in\mathcal{A}_{01}}$.
  This gives ${u_0\colon V'_0\to V}$ in $\mathcal{W}$ and morphisms
  ${f'_a\colon V'_0\to\widetilde{V}_1}$, and with ${\phi'_0=q(u_0)}$ the first
  identity of~\eqref{eq:lifting} holds for all
  arrows. For ${\rho\in\mathcal{R}}$ put
  \[
    g_\rho=\sum_{b,a}\rho_{ba}f_bf'_a\colon V'_0\longrightarrow\widetilde{V}_2.
  \]
  Then
  \[
    \phi_2q(g_\rho)=\Bigl(\sum_{b,a}\rho_{ba}t_bt_a\Bigr)\phi'_0=0,
  \]
  so ${q(g_\rho)=0}$. By \Cref{lemma:fractions}\eqref{it:fractions-kill} there
  is ${v\colon\widetilde{V}_0\to V'_0}$ in $\mathcal{W}$ with ${g_\rho v=0}$ for
  every $\rho$. Put ${f_a=f'_av}$ for ${a\in\mathcal{A}_{01}}$ and
  ${\phi_0=\phi'_0q(v)}$. The first identity of~\eqref{eq:lifting} still holds,
  and ${\sum_{b,a}\rho_{ba}f_bf_a=g_\rho v=0}$ in $\mathcal{K}$; no further
  compatibility arises, since no arrow ends at~$0$.

  Finally, choose chain maps representing the morphisms $f_a$ of $\mathcal{K}$.
  For each ${\rho\in\mathcal{R}}$ the chain map $f_\rho$ is then null-homotopic,
  and we choose a homotopy $h_\rho$ with ${dh_\rho+h_\rho d=f_\rho}$.
\end{proof}

All the data of \Cref{prop:lifting} are chosen from $\mathbb V$, and do not
depend on a module~$Y$. Applying $\operatorname{ev}_Y$ to
the first identity of~\eqref{eq:lifting} and using \Cref{lemma:evaluate-V}, we
obtain for every $Y$ isomorphisms
${\operatorname{ev}_Y(\widetilde{V}_i)\cong HY\oplus HY[3]}$ under which
$\operatorname{ev}_Y(f_s)$ and $\operatorname{ev}_Y(f_{s'})$ are identities and
$\operatorname{ev}_Y(f_{a_h})$ and $\operatorname{ev}_Y(f_{a'_h})$ are the
action of $x_h$ on both summands.

\begin{remark}
  \label{rem:uncontrolled-length}
  The construction uses choices not given by explicit formulas at four points:
  the
  splitting morphisms in the proof of \Cref{lemma:cone-splitting}, the roofs
  representing the morphisms $t_a$ in \Cref{lemma:fractions}, the Ore squares
  producing common denominators, and the denominator $v$ for which ${g_\rho v=0}$
  for every ${\rho\in\mathcal R}$. A proof that a cone lies in the thick
  subcategory
  $\mathcal{S}$ does not by itself provide a bounded number of cones, shifts
  and direct summands witnessing it, and none of these steps specifies the
  complexes $\widetilde{V}_i$ or bounds their length. Once the
  $\widetilde{V}_i$ and the $f_a$ are given as complexes and chain maps, the
  homotopies $h_\rho$ are solutions of finite systems of linear equations, and
  \Cref{subsec:rectification} is explicit: if all $\widetilde{V}_i$ are
  concentrated in degrees $[a,b]$, the complex $P$ of
  \Cref{prop:rectification} is concentrated in degrees ${[a-2,b]}$. The argument
  therefore does not determine the integer $l$ used in \Cref{sec:simulation}.
  This is the
  absence of a bound in the argument, not a proof that no bound or algorithm
  exists; see \Cref{rem:what-is-explicit}.
\end{remark}

\subsection{Rectification}
\label{subsec:rectification}

The following construction applies to every algebra ${B}={\kk Q/(J_2)}$ for which
$Q$
has vertices $0$, $1$ and $2$, finite sets of arrows $\mathcal{A}_{01}$ from
$0$ to $1$ and $\mathcal{A}_{12}$ from $1$ to $2$ and no other arrows, and
$J_2$ is a subspace of the span of the paths of length two, with basis
$\mathcal{R}$.

\begin{proposition}
  \label{prop:rectification}
  Let $B$ be as above, and suppose given bounded complexes
  $\widetilde{V}_0$, $\widetilde{V}_1$ and $\widetilde{V}_2$ of finitely
  generated projective right $B$-modules, chain maps $f_a$ and homogeneous
  maps $h_\rho$ of degree ${-1}$ of right modules satisfying the second identity
  of~\eqref{eq:lifting}. Then there are a bounded complex $P$ of
  finite-dimensional $B$-bimodules whose terms are projective right modules,
  homotopy equivalences ${\iota_i\colon\widetilde{V}_i\to e_iP}$ of complexes of
  right modules, and, for every arrow ${a\colon i\to j}$, a homotopy
  ${a\iota_i\simeq\iota_jf_a}$, where $a$ acts on $P$ by left multiplication.
\end{proposition}

\begin{proof}
  All unadorned tensor products in this proof are over $\kk$. On a tensor
  product ${Be_j\otimes\widetilde{V}_i}$, the algebra $B$ acts from the left on
  the first factor and from the right on the second, and the differential is
  ${\id\otimes d}$. Consider the complexes of bimodules
  \[
    L_0=\bigoplus_{i}Be_i\otimes\widetilde{V}_i,\qquad
    L_1=\bigoplus_{a\colon i\to j}Be_j\otimes\widetilde{V}_i,\qquad
    L_2=\bigoplus_{\rho\in\mathcal{R}}Be_2\otimes\widetilde{V}_0,
  \]
  with differentials $d_0$, $d_1$ and $d_2$, and the bimodule maps
  ${\partial_1\colon L_1\to L_0}$ and ${\partial_2\colon L_2}\to{ L_1}$ of degree
  zero and ${\eta\colon L_2\to L_0}$ of degree ${-1}$ defined by
  \begin{equation}
    \label{eq:rectification-maps}
    \begin{aligned}
      \partial_1(c\otimes v)_a&=(ca\otimes v)_i-(c\otimes f_av)_j
        &&\text{for }a\colon i\to j,\\
      \partial_2(c\otimes v)_\rho&=\sum_{b,a}\rho_{ba}\bigl((cb\otimes v)_a
        +(c\otimes f_av)_b\bigr),\\
      \eta(c\otimes v)_\rho&=(c\otimes h_\rho v)_2,
    \end{aligned}
  \end{equation}
  where a subscript names the summand containing an element: in $L_0$ the
  summand of a vertex, in $L_1$ that of an arrow. These maps are
  bimodule homomorphisms because paths act on the first factor and $f_a$ and
  $h_\rho$ are homomorphisms of right modules on the second. Since the $f_a$
  are chain maps, ${d_0\partial_1=\partial_1d_1}$ and
  ${d_1\partial_2=\partial_2d_2}$. In the composite
  \[
    \partial_1\partial_2(c\otimes v)_\rho=\sum_{b,a}\rho_{ba}\bigl(
    (cba\otimes v)_0-(cb\otimes f_av)_1+(cb\otimes f_av)_1
    -(c\otimes f_bf_av)_2\bigr)=-(c\otimes f_\rho v)_2
  \]
  the two middle terms cancel, and the first term vanishes since
  ${\sum\rho_{ba}ba=\rho=0}$ in~$B$. Together with~\eqref{eq:lifting} this gives
  \begin{equation}
    \label{eq:eta-homotopy}
    d_0\eta+\eta d_2=-\partial_1\partial_2.
  \end{equation}
  Put
  \[
    P\coloneqq L_0\oplus L_1[1]\oplus L_2[2],\qquad
    d_P=\begin{pmatrix}d_0&\partial_1&\eta\\0&-d_1&\partial_2\\0&0&d_2
    \end{pmatrix}.
  \]
  In total degree $m$ the three components lie in degrees $m$, ${m+1}$ and
  ${m+2}$ of $L_0$, $L_1$ and $L_2$, so all entries of $d_P$ have degree one; the
  diagonal signs are those of the shifts. The entries of $d_P^2$ are $d_0^2$,
  $d_1^2$, $d_2^2$, ${d_0\partial_1-\partial_1d_1}$,
  ${\partial_2d_2-d_1\partial_2}$ and ${d_0\eta+\partial_1\partial_2+\eta d_2}$,
  and all of them vanish by the identities above and~\eqref{eq:eta-homotopy}.
  As a right module, ${Be_j\otimes\widetilde{V}_i^m}$ is a finite direct sum of
  copies of $\widetilde{V}_i^m$, so the terms of $P$ are finite-dimensional and
  projective on the right, and $P$ is bounded.

  For a vertex $i$, let ${\iota_i\colon\widetilde{V}_i\to e_iP}$ send $v$ to
  ${(e_i\otimes v)_i\in e_iL_0}$. It is a chain map and, in each degree, a split
  injection of right modules. We show that the quotient
  ${C_i=e_iP/\iota_i(\widetilde{V}_i)}$ is acyclic. Filter $e_iP$ by the
  subcomplexes $F^r$ consisting of the summands whose second factor is
  $\widetilde{V}_j$ with ${j\geq r}$; the differential preserves this filtration,
  since $d$ and the path terms preserve $j$, the terms involving $f_a$ increase
  it, and $\eta$ increases it from $0$ to~$2$. The induced filtration of $C_i$
  has the following subquotients. The subquotient for ${j=i}$ is zero, since the
  only summand of $e_iP$ with second factor $\widetilde{V}_i$ is
  ${e_iBe_i\otimes\widetilde{V}_i=\kk e_i\otimes\widetilde{V}_i}$. The
  subquotients for ${j>i}$ are zero, since ${e_iBe_j=0}$. For $(j,i)$ equal to
  $(0,1)$ or $(1,2)$, the subquotient is the total complex of the double
  complex ${K\otimes\widetilde{V}_j}$, where $K$ is the complex
  \[
    0\longrightarrow\kk\{a\colon j\to i\}\xrightarrow{a\mapsto a}
    e_iBe_j\longrightarrow0
  \]
  in degrees ${-1}$ and~$0$, which is exact since all relations have length two.
  For ${(j,i)=(0,2)}$ it is the total complex of ${K\otimes\widetilde{V}_0}$ with
  \[
    K\colon\ 0\longrightarrow\kk\mathcal{R}
    \xrightarrow{\rho\mapsto\sum\rho_{ba}(b,a)}
    \kk\{(b,a)\}\xrightarrow{(b,a)\mapsto ba}e_2Be_0\longrightarrow0
  \]
  in degrees ${-2}$, ${-1}$ and~$0$, where $(b,a)$ runs through
  ${\mathcal{A}_{12}\times\mathcal{A}_{01}}$. This complex is exact because
  $\mathcal{R}$ is linearly independent with span $J_2$ and $e_2Be_0$ is the
  quotient of the span of the paths of length two by $J_2$; here it matters
  that $\mathcal{R}$ is a basis and not merely a spanning set. An exact complex
  $K$ of vector spaces has a contraction $\sigma$ with
  ${\partial_K\sigma+\sigma\partial_K=\id}$, and ${\sigma\otimes\id}$ contracts the
  total complex of ${K\otimes\widetilde{V}_j}$: in horizontal degree $m$ its
  differential is
  \[
    \partial_K\otimes\id+(-1)^m\id\otimes d,
  \]
  and the terms of
  ${(\sigma\otimes\id)(\id\otimes d)}$ and ${(\id\otimes d)(\sigma\otimes\id)}$
  appear with the opposite signs ${(-1)^m}$ and ${(-1)^{m-1}}$. These signs agree
  with those of $d_P$. Hence, all subquotients are acyclic, and so is $C_i$, by
  induction along the finite filtration and the long exact cohomology sequence.

  The terms of $C_i$ are projective right modules, since $\iota_i$ is split in
  each degree, and a bounded acyclic complex of projective modules is
  contractible. Choose a splitting ${e_iP=\widetilde{V}_i\oplus C_i}$ of graded
  right modules, in which $d_P$ has the form
  $\left(\begin{smallmatrix}d&t\\0&d_C\end{smallmatrix}\right)$ with
  ${dt+td_C=0}$, and a contraction $\sigma$ of $C_i$. Then
  ${r(x,c)=x-t\sigma c}$ is a chain map with ${r\iota_i=\id}$, and ${K(x,c)=(0,\sigma
  c)}$ satisfies ${d_PK+Kd_P=\id-\iota_ir}$. Hence, $\iota_i$ is a homotopy
  equivalence.

  For an arrow ${a\colon i\to j}$, let ${H_a\colon\widetilde{V}_i\to e_jP}$ be the
  map of degree ${-1}$ sending $v$ to ${(e_j\otimes v)_a\in e_jL_1[1]}$. The
  component of $d_PH_a(v)$ in $L_1[1]$ is ${-(e_j\otimes dv)_a=-H_a(dv)}$, and its
  component in $L_0$ is
  \[
    \partial_1(e_j\otimes v)_a=(a\otimes v)_i-(e_j\otimes f_av)_j.
  \]
  Hence,
  \[
    d_PH_a+H_ad=a\iota_i-\iota_jf_a.\qedhere
  \]
\end{proof}

The three columns of $P$ are indexed by the vertices, the arrows and the
relations of $Q$. Since every relation starts at $0$ and ends at $2$, no
relation can be composed with a further arrow, and no homotopies beyond $h_\rho$
are required.

\subsection{The realisation theorem}
\label{subsec:realisation-theorem}

\begin{theorem}
  \label{thm:realisation}
  Let $(R,\Psi)$ be selection data over $\kk$ in which $R$ is finitely
  presented, and put ${H=\Psi\otimes_R-}$. Let $B$ and $M$ be the algebra
  and the functor of \Cref{cons:encoding-algebra} for a presentation of $R$ by
  relations of degree at most two. Then there exists a bounded complex $P$ of
  finite-dimensional $B$-bimodules, whose terms are projective as right
  modules, such that for every finite-dimensional left $R$-module $Y$
  \begin{equation}
    \label{eq:realisation}
    P\Lotimes[B]M(Y)\cong M(HY)\oplus M(HY)[3]\quad\text{in }
    \DerCat[b]{\operatorname{mod}B}.
  \end{equation}
  The complex $P$ does not depend on~$Y$.
\end{theorem}

\begin{proof}
  Let $P$ be the complex obtained by applying \Cref{prop:rectification} to the
  data of \Cref{prop:lifting}; it depends only on $(R,\Psi)$, the presentation
  and the choices made there. Fix $Y$ and put ${T_Y=P\otimes_BM(Y)}$, which
  computes ${P\Lotimes[B]M(Y)}$ since $P$ is a bounded complex of projective right
  modules. For each vertex $i$ we have ${e_iT_Y=\operatorname{ev}_Y(e_iP)}$, and
  the action of an arrow $a$ on $T_Y$ is $\operatorname{ev}_Y$ of left
  multiplication by~$a$.

  Applying $\operatorname{ev}_Y$ to the homotopy equivalences $\iota_i$ and the
  homotopies of \Cref{prop:rectification} identifies $\operatorname{ev}_Y(e_iP)$
  with $\operatorname{ev}_Y(\widetilde{V}_i)$, compatibly with the arrows, and
  the paragraph after \Cref{prop:lifting} identifies the latter with
  ${HY\oplus HY[3]}$. Taking cohomology, we obtain isomorphisms of vector spaces
  at each vertex that commute with the actions of all arrows, hence
  isomorphisms of left $B$-modules
  \[
    H^m(T_Y)\cong
    \begin{cases}
      M(HY)&\text{if }m\in\set{0,-3},\\
      0&\text{otherwise.}
    \end{cases}
  \]

  Put ${N_0=H^0(T_Y)}$ and ${N_{-3}=H^{-3}(T_Y)}$. The good truncation $T'$ of
  $T_Y$ to the degrees ${[-3,0]}$ is quasi-isomorphic to $T_Y$, and its
  cohomology in degree ${-3}$ is its degree ${-3}$ cycle module, which is a
  subcomplex isomorphic to $N_{-3}[3]$. The quotient complex has cohomology
  $N_0$ in degree zero only and terms in nonpositive degrees, so it is
  quasi-isomorphic to $N_0$. The short exact sequence of complexes gives a
  triangle
  \[
    N_{-3}[3]\longrightarrow T_Y\longrightarrow N_0\xrightarrow{\ \delta\ }
    N_{-3}[4]
  \]
  in $\DerCat[b]{\operatorname{mod}B}$. By \Cref{prop:encoding-algebra},
  $N_0$ has a projective resolution of length at most two, so that
  ${\delta\in\Hom{N_0}{N_{-3}[4]}=\operatorname{Ext}^4_B(N_0,N_{-3})=0}$. Hence,
  the identity of $N_0$ lifts to a morphism ${N_0\to T_Y}$, and together with
  the first morphism of the triangle it gives a morphism
  ${N_{-3}[3]\oplus N_0\to T_Y}$, which is a quasi-isomorphism by the long exact
  cohomology sequence. This proves~\eqref{eq:realisation}.
\end{proof}

\begin{corollary}
  \label{coro:realisation-iterates}
  In the situation of \Cref{thm:realisation}, put ${\mathcal{F}=P\Lotimes[B]-}$.
  For every finite-dimensional left $R$-module $Y$ and every integer
  ${j\geq0}$,
  \[
    \mathcal{F}^jM(Y)\cong\bigoplus_{r=0}^{j}
    \bigl(M(H^jY)[3r]\bigr)^{\oplus\binom{j}{r}}.
  \]
  In particular, ${\mathcal{F}^jM(Y)=0}$ if and only if ${H^jY=0}$.
\end{corollary}

\begin{proof}
  For ${j=0}$ there is nothing to prove. Suppose the formula holds for $j$. The
  module $H^jY$ is finite-dimensional by \Cref{lemma:selection-matrix}, and
  $\mathcal{F}$ commutes with finite direct sums and shifts, so
  \Cref{thm:realisation} applied to $H^jY$ gives
  \[
    \mathcal{F}^{j+1}M(Y)\cong\bigoplus_{r=0}^{j}
    \bigl(M(H^{j+1}Y)[3r]\oplus M(H^{j+1}Y)[3r+3]\bigr)^{\oplus\binom{j}{r}},
  \]
  in which the shift $3r$ occurs with multiplicity
  ${\binom{j}{r}+\binom{j}{r-1}=\binom{j+1}{r}}$. The last assertion follows
  from \Cref{prop:encoding-algebra}.
\end{proof}

\begin{remark}
  \label{rem:realisation-special-case}
  The main preprint~\cite[Sections~3 and~4]{OAI26findim} treats the case of a
  central idempotent
  ${e\in R}$ and an algebra endomorphism $\alpha$ of $R$, with
  ${\Psi={}_\alpha(eR)}$ and ${HY={}_\alpha(eY)}$. This is the case ${n=1}$,
  ${\varepsilon=e}$ and ${\rho(r)=e\alpha(r)e}$ of \Cref{lemma:selection-matrix};
  the centrality of $e$ makes $\rho$ multiplicative. The proof above uses
  neither the centrality of $\varepsilon$ nor any property of $\rho$ beyond
  its being a homomorphism of unital $\kk$-algebras into
  $\varepsilon M_n(R)\varepsilon$, and it does not use that $\Psi$ is
  projective as a left module.
\end{remark}

\section{Trivial extensions from bimodule complexes}
\label{sec:simulation}

\Cref{thm:realisation} provides a complex $P$ of bimodules whose derived tensor
functor iterates the selection functor, while \Cref{coro:extinction} requires
a single bimodule over an algebra of finite global dimension. A bounded complex
of bimodules is simulated, up to a shift, by the square of the derived tensor
functor of an ordinary bimodule over a larger algebra.

\subsection{An ordinary bimodule associated with a complex}
\label{subsec:simulation}

For an integer ${l\geq0}$, let $K_l$ be the path algebra of the quiver
\[
  0\xrightarrow{c_1}1\xrightarrow{c_2}\cdots\xrightarrow{c_l}l
\]
modulo all paths of length two, with vertex idempotents $1_0,\dots,1_l$, so that
${c_i\in1_iK_l1_{i-1}}$. Let $W$ be the simple right $K_l$-module at the
vertex~$l$. For ${i>0}$ the right module $1_iK_l$ has basis $1_i,c_i$ and radical
spanned by $c_i$, and ${1_0K_l=\kk1_0}$. Hence, left multiplication by $c_i$ maps
$1_{i-1}K_l$ onto the radical of $1_iK_l$ with kernel the radical of
$1_{i-1}K_l$, and
\begin{equation}
  \label{eq:chain-resolution}
  0\longrightarrow1_0K_l\xrightarrow{c_1\cdot}1_1K_l\xrightarrow{c_2\cdot}
  \cdots\xrightarrow{c_l\cdot}1_lK_l\longrightarrow W\longrightarrow0
\end{equation}
is a projective resolution of $W$; we write $R_W$ for the complex with
${R_W^n=1_{l+n}K_l}$ for ${-l\leq n\leq0}$.

\begin{proposition}
  \label{prop:simulation}
  Let $B$ be a finite-dimensional algebra whose left and right global
  dimensions are at most $d_L$ and $d_R$, and let $P$ be a bounded complex of
  finite-dimensional $B$-bimodules, projective as right modules, with ${P^j=0}$
  for ${j\notin[a,b]}$; put ${l=b-a\geq0}$. Put ${B_1=B\otimes_\kk K_l}$ and
  ${\Delta=B\times B_1}$, and let $\pi_0$ and $\pi_1$ be the central idempotents
  of $\Delta$ corresponding to the two factors. Define
  \begin{itemize}
    \item the $(B_1,B)$-bimodule $O$ by ${1_iO=P^{a+i}}$ for ${0\leq i\leq l}$, with
      ${c_i\colon1_{i-1}O\to1_iO}$ acting as the differential $d_P^{a+i-1}$ and
      with the two actions of $B$ on the terms of $P$;
    \item the $(B,B_1)$-bimodule ${Y=B\otimes_\kk W}$, with
      ${b'(b\otimes w)(b''\otimes c)=b'bb''\otimes wc}$;
  \end{itemize}
  and regard ${X=O\oplus Y}$ as a $\Delta$-bimodule with ${O=\pi_1O\pi_0}$ and
  ${Y=\pi_0Y\pi_1}$. Let ${\Phi=X\Lotimes[\Delta]-}$. Then for every complex $N$
  of left $B$-modules, regarded as a complex of $\Delta$-modules through the
  first factor,
  \[
    \Phi^2N\cong(P\Lotimes[B]N)[b]\quad\text{in }\DerCat{\operatorname{Mod}
    \Delta}.
  \]
  Moreover, the left and right global dimensions of $\Delta$ are at most
  ${d_L+l}$ and ${d_R+l}$.
\end{proposition}

\begin{proof}
  The bimodule $O$ is well defined: the relations ${c_{i+1}c_i=0}$ hold because
  ${d_P^2=0}$, and the actions of $B$ commute with those of the $c_i$ because the
  differential of $P$ is a bimodule map. As a right $B$-module, $O$ is the
  direct sum of the terms of $P$, hence projective. The complex
  ${R_Y=B\otimes_\kk R_W}$ is a resolution of $Y$ by $(B,B_1)$-bimodules whose
  terms ${B\otimes_\kk1_{l+n}K_l=(1\otimes1_{l+n})B_1}$ are projective right
  $B_1$-modules. In degree ${n\in[-l,0]}$ there is an isomorphism of
  $B$-bimodules
  \[
    (B\otimes_\kk1_{l+n}K_l)\otimes_{B_1}O\longrightarrow1_{l+n}O=P^{b+n},
    \qquad(\beta\otimes u)\otimes o\longmapsto\beta uo,
  \]
  with inverse ${o\mapsto(1\otimes1_{l+n})\otimes o}$, and under these
  isomorphisms the differential of ${R_Y\otimes_{B_1}O}$ becomes $d_P$, since
  the differential of $R_W$ is left multiplication by the arrows. The complex
  $P[b]$ has the same terms in the same degrees and the differential
  ${(-1)^bd_P}$; the maps ${(-1)^{bn}\id}$ on the terms of degree $n$ form an
  isomorphism of complexes of bimodules ${R_Y\otimes_{B_1}O\cong P[b]}$, since
  ${(-1)^{b(n+1)}=(-1)^b(-1)^{bn}}$.

  The complex ${R_X=O\oplus R_Y}$ is a bounded resolution of $X$ by
  $\Delta$-bimodules projective as right modules, so it computes $\Phi$ on all
  complexes. For a complex $N$ of left $B$-modules, one application gives
  ${O\otimes_BN}$, a complex of $B_1$-modules, and a second gives
  \[
    R_Y\otimes_{B_1}(O\otimes_BN)\cong(R_Y\otimes_{B_1}O)\otimes_BN
    \cong P[b]\otimes_BN,
  \]
  a complex of $B$-modules; the associativity isomorphism involves no sign
  since the order of the factors is unchanged. Since $P[b]$ is a bounded complex
  of projective right modules, ${P[b]\otimes_BN}$ computes
  ${(P\Lotimes[B]N)[b]}$.

  For the global dimensions, the resolution~\eqref{eq:chain-resolution},
  ending at any vertex, shows that the right simple $K_l$-modules have
  projective dimension at most $l$, and similarly for the left simple modules,
  whose successive syzygies are simple at the following vertices. The ideal
  ${I=J_B\otimes K_l+B\otimes J_K}$ of $B_1$, where $J_B$ and $J_K$ denote the
  radicals, is nilpotent, being the sum of two commuting nilpotent ideals, and
  ${B_1/I\cong(B/J_B)^{l+1}}$ is semisimple. Hence, $I$ is the radical of $B_1$
  and the simple left $B_1$-modules are the modules ${S\otimes_\kk L}$ for simple
  modules $S$ over $B$ and $L$ over $K_l$. The tensor product of projective
  resolutions of $S$ and $L$ is a projective resolution of ${S\otimes_\kk L}$ of
  length at most ${d_L+l}$. Since every module has a finite filtration by
  semisimple modules, the left global dimension of $B_1$ is at most ${d_L+l}$;
  the same argument applies on the right, and the bounds pass to the product
  ${\Delta=B\times B_1}$.
\end{proof}

The vertices of $K_l$ correspond to the degrees of $P$, the arrows act on $O$ by
the differential of $P$, and ${Y\Lotimes[B_1]O\cong P[b]}$. The construction uses
${l+2}$ copies of $B$; the resulting number of simple modules and the bounds on
global dimension grow linearly with $l$.

\subsection{The main theorem}
\label{subsec:main-theorem}

\begin{theorem}
  \label{thm:main}
  Let $(R,\Psi)$ be selection data over a field $\kk$ in which $R$ is finitely
  presented, and let ${H=\Psi\otimes_R-}$. Then there is a finite-dimensional
  algebra ${A=\Delta\ltimes X}$ over $\kk$, with $\Delta$ of finite global
  dimension, and an integer ${l\geq0}$ with the following property: for every
  finite-dimensional left $R$-module $Y$ of finite extinction time ${t\geq1}$,
  there is a finite-dimensional left $A$-module $N$ with
  \[
    2t-2\leq\operatorname{pd}_AN\leq(l+2)+(2t-1)(l+3).
  \]
  In particular, if $(R,\Psi)$ has unbounded extinction, then
  ${\operatorname{findim}A=\infty}$.
\end{theorem}

\begin{proof}
  Fix a presentation of $R$ by relations of degree at most two, the algebra $B$
  and the functor $M$ of \Cref{cons:encoding-algebra}, and the
  complex $P$ of \Cref{thm:realisation}. Choose integers ${a\leq b}$ with
  ${P^j=0}$ for ${j\notin[a,b]}$, put ${l=b-a}$, and let $\Delta$, $X$ and $\Phi$ be
  as in \Cref{prop:simulation}; put ${A=\Delta\ltimes X}$. By
  \Cref{prop:encoding-algebra}, $B$ has left and right global dimension at
  most two, so $\Delta$ has left and right global dimension at most ${l+2}$ by
  \Cref{prop:simulation}. All these choices are made before the module $Y$ is
  chosen.

  Let $Y$ be a finite-dimensional left $R$-module of extinction time ${t\geq1}$,
  and let $N$ be $M(Y)$, regarded as a $\Delta$-module through the first factor
  and then as an $A$-module. By \Cref{prop:simulation}, ${\Phi^{2r}N\cong
  (\mathcal{F}^rM(Y))[rb]}$ for every ${r\geq0}$, where
  ${\mathcal{F}=P\Lotimes[B]-}$, and \Cref{coro:realisation-iterates} gives
  \[
    \Phi^{2r}N\cong\bigoplus_{j=0}^r\bigl(M(H^rY)[rb+3j]\bigr)^{\oplus
    \binom{r}{j}}.
  \]
  Hence, ${\Phi^{2t}N\simeq0}$, since ${H^tY=0}$, and ${\Phi^{2t-2}N\not\simeq0}$,
  since ${H^{t-1}Y\neq0}$ and hence ${M(H^{t-1}Y)\neq0}$ by
  \Cref{prop:encoding-algebra}. By \Cref{coro:extinction}, with
  ${d_L=d_R=l+2}$ and with $2t$ in place of $t$, the module $N$ has finite
  projective dimension at most ${(l+2)+(2t-1)(l+3)}$, and by
  \Cref{prop:pd-formula} at least ${2t-2}$. The last assertion follows since
  every finite-dimensional $A$-module is finitely generated.
\end{proof}

\begin{corollary}[{\cite[Theorem~1.1]{OAI26findim}}]
  \label{coro:main}
  There is a finite-dimensional complex algebra $A$ with
  ${\operatorname{findim}A=\infty}$. More precisely, there are an integer
  ${l\geq0}$ and a finite-dimensional complex algebra ${A=\Delta\ltimes X}$ with
  ${3(l+2)}$ simple modules, where $\Delta$ has global dimension at most ${l+2}$,
  such that for every ${m\geq1}$ some finite-dimensional $A$-module $N_m$
  satisfies ${2m-2\leq\operatorname{pd}_AN_m<\infty}$.
\end{corollary}

\begin{proof}
  Apply \Cref{thm:main} to the selection data of
  \Cref{thm:selection-unbounded} and to the modules $Y_m$ of extinction time
  $m$. The simple $A$-modules are the simple $\Delta$-modules, by the proof of
  \Cref{prop:pd-formula}. The algebra $B$ has three simple modules,
  one for each vertex, and the simple modules of ${B_1=B\otimes K_l}$ are the
  ${3(l+1)}$ modules ${S\otimes L}$ of the proof of \Cref{prop:simulation}.
\end{proof}

\begin{remark}
  \label{rem:what-is-explicit}
  The group $G$ and its
  presentation (\Cref{subsec:group}), the selection data
  (\Cref{subsec:unbounded-extinction}), the algebra $B$ with its
  dimension (\Cref{prop:encoding-algebra}), the rectification
  (\Cref{prop:rectification}) once its input is given, and the simulation
  (\Cref{prop:simulation}) are given by formulas. The exception is the input of
  the rectification: the complexes $\widetilde{V}_i$, the chain maps $f_a$ and
  the homotopies $h_\rho$ of \Cref{prop:lifting} exist by the calculus of
  fractions, but the proof gives no bound on the lengths of the
  $\widetilde{V}_i$ (\Cref{rem:uncontrolled-length}). Consequently, the integer
  $l$, the number ${3(l+2)}$ of simple modules of $A$ and the dimension of $A$
  are not determined by the argument. Whether the lifting can be made
  effective, with a bound on $l$, is not settled here. Applying the procedure of
  \Cref{subsec:encoding} to the chosen presentation of $\mathbb{C}G$ introduces
  auxiliary generators for the relators of degree larger than two, in addition
  to the thirty symbols for the fifteen generators of $G$ and their inverses.
  For a resulting presentation with $d$ generators, $B$ has ${2(d+1)}$ arrows.
\end{remark}

\begin{remark}
  \label{rem:simulation-bound}
  The main preprint bounds the global dimensions of $\Delta$ by ${3l+2}$, using
  the ordering of the vertices~\cite[Proposition~5.1]{OAI26findim}; the bound ${l+2}$ of \Cref{prop:simulation} comes from the tensor product
  description of the radical of $B_1$. Neither bound is needed beyond
  finiteness, except for the explicit upper bound in \Cref{thm:main}.
\end{remark}

\section{Tate cohomology and the Auslander--Reiten conjecture}
\label{sec:conversion}

By \Cref{thm:ar-to-findim}, a counterexample to the Auslander--Reiten
conjecture~\cite{AR75a} yields an algebra of infinite little finitistic
dimension. The preprint~\cite{OAI26ar} constructs a counterexample from data in
the stable
category of a symmetric algebra $E$: a bimodule $F$ and a stable morphism
${v\colon S\to F\otimes_ES}$. \Cref{thm:conversion} gives this construction over
an arbitrary field; the preprint treats characteristic two, in which the signs
disappear. The theorem reduces the vanishing of the self-extensions of a module
over the triangular
algebra $\left(\begin{smallmatrix}E&0\\F&E\end{smallmatrix}\right)$ to the
bijectivity of a map between Tate cohomology groups over~$E$.

\subsection{Complete resolutions}
\label{subsec:complete-resolutions}

Let $A$ be a finite-dimensional algebra. For complexes $\mathbb{P}$ and
$\mathbb{W}$ of $A$-modules, the Hom complex has the terms and differential
\[
  \Hom[A]{\mathbb{P}}{\mathbb{W}}^a=\prod_{n\in\ZZ}\Hom[A]{\mathbb{P}^n}
  {\mathbb{W}^{n+a}},\qquad\partial f=d_{\mathbb{W}}f-(-1)^afd_{\mathbb{P}};
\]
its cocycles of degree $a$ are the chain maps ${\mathbb{P}\to\mathbb{W}[a]}$,
up to the signs of the shift, and its coboundaries are the null-homotopic ones.
For an exact complex $\mathbb{P}$ put
${C^n(\mathbb{P})=\operatorname{coker}(\mathbb{P}^{n-1}\to\mathbb{P}^n)}$; the
differential induces injections ${C^n(\mathbb{P})\to\mathbb{P}^{n+1}}$ and short
exact sequences
\begin{equation}
  \label{eq:cokernel-sequences}
  0\longrightarrow C^n(\mathbb{P})\longrightarrow\mathbb{P}^{n+1}
  \longrightarrow C^{n+1}(\mathbb{P})\longrightarrow0.
\end{equation}
An exact complex $\mathbb{P}$ of finitely generated projective modules is
\emph{totally acyclic} if $\Hom[A]{\mathbb{P}}{A}$ is exact; then
$\Hom[A]{\mathbb{P}}{L}$ is exact for every projective module $L$, since
${\Hom[A]{\mathbb{P}}{A^{(I)}}=\Hom[A]{\mathbb{P}}{A}^{(I)}}$ by finite
generation. A finitely generated module is \emph{Gorenstein-projective} if it
is isomorphic to $C^0(\mathbb{P})$ for a totally acyclic complex $\mathbb{P}$,
which is then called a \emph{complete resolution} of it.

\begin{lemma}
  \label{lemma:complete-resolutions}
  Let $\mathbb{P}$ be a totally acyclic complex and $\mathbb{W}$ an exact
  complex of finitely generated projective $A$-modules, and put
  ${M=C^0(\mathbb{P})}$ and ${N=C^0(\mathbb{W})}$. Then for every integer $a$
  there are isomorphisms
  \[
    H^a\Hom[A]{\mathbb{P}}{\mathbb{W}}\cong
    \underline{\operatorname{Hom}}_A(C^{-a}(\mathbb{P}),N)\cong
    \underline{\operatorname{Hom}}_A(M,C^a(\mathbb{W})),
  \]
  compatible with composition and with additive functors that preserve the
  complexes involved, and for ${a>0}$ these groups are isomorphic to
  $\operatorname{Ext}^a_A(M,N)$. Moreover, ${\operatorname{Ext}^a_A(M,L)=0}$ for
  every ${a>0}$ and every projective module~$L$.
\end{lemma}

\begin{proof}
  Let ${\epsilon_n\colon\mathbb{P}^n\to C^n(\mathbb{P})}$ and
  ${\jmath_n\colon C^n(\mathbb{P})\to\mathbb{P}^{n+1}}$ be the induced maps. If
  $L$ is projective, every map ${u\colon C^n(\mathbb{P})\to L}$ extends across
  $\jmath_n$: the map $u\epsilon_n$ is a cocycle of $\Hom[A]{\mathbb{P}}{L}$,
  hence a coboundary, so ${u\epsilon_n=bd^n_{\mathbb{P}}}$ for some
  ${b\colon\mathbb{P}^{n+1}\to L}$, and ${b\jmath_n=u}$.

  A chain map ${f\colon\mathbb{P}\to\mathbb{W}}$ induces a map ${M\to N}$.
  Conversely, given ${u\colon M\to N}$, lift $u\epsilon_0$ to
  ${f^0\colon\mathbb{P}^0\to\mathbb{W}^0}$, and construct $f^n$ for ${n<0}$ by
  lifting through the surjections ${\mathbb{W}^{n}\to\operatorname{im}
  d^{n}_{\mathbb{W}}}$, using projectivity of the $\mathbb{P}^n$. For ${n\geq0}$,
  once $f^n$ is chosen, $d^n_{\mathbb{W}}f^n$ vanishes on the image of
  $d^{n-1}_{\mathbb{P}}$ and so factors through $C^n(\mathbb{P})$; the
  extension property gives $f^{n+1}$ with
  ${f^{n+1}d^n_{\mathbb{P}}=d^n_{\mathbb{W}}f^n}$.

  If ${f=d_{\mathbb{W}}h+hd_{\mathbb{P}}}$, then the induced map is induced by
  $h^1d^0_{\mathbb{P}}$ and factors through $\mathbb{P}^1$. Conversely, suppose
  the induced map factors as ${M\xrightarrow{\alpha}L\xrightarrow{\beta}N}$ with
  $L$ projective. Extending $\alpha$ to $\mathbb{P}^1$ and lifting $\beta$ to
  $\mathbb{W}^0$ gives ${h^1\colon\mathbb{P}^1\to\mathbb{W}^0}$ whose boundary
  induces $\beta\alpha$; subtracting it, we may assume that $f$ induces zero.
  Then $f^0$ takes values in the image of $d^{-1}_{\mathbb{W}}$, and we
  construct $h$ with ${f=d_{\mathbb{W}}h+hd_{\mathbb{P}}}$ by two recursions,
  starting with ${h^1=0}$ and a lift $h^0$ of $f^0$ through
  $d^{-1}_{\mathbb{W}}$. For ${n<0}$, the map ${f^n-h^{n+1}d^n_{\mathbb{P}}}$ takes
  values in ${\ker d^n_{\mathbb{W}}=\operatorname{im}d^{n-1}_{\mathbb{W}}}$ and
  lifts to $h^n$ by projectivity of $\mathbb{P}^n$. For ${n\geq1}$, the map
  ${f^n-d^{n-1}_{\mathbb{W}}h^n}$ vanishes on the image of
  $d^{n-1}_{\mathbb{P}}$, hence factors through $C^n(\mathbb{P})$, and the
  extension property gives $h^{n+1}$. Since the Hom complex is a product, no
  finiteness condition arises. This proves the case ${a=0}$; the compatibility
  with composition and with additive functors follows from the construction.

  The case of general $a$ follows by applying the case ${a=0}$ to the pairs
  ${(\mathbb{P}[-a],\mathbb{W})}$ and $(\mathbb{P},\mathbb{W}[a])$, since
  ${C^0(\mathbb{P}[-a])=C^{-a}(\mathbb{P})}$ and ${C^0(\mathbb{W}[a])=
  C^a(\mathbb{W})}$. For ${a>0}$, the truncation ${\cdots\to\mathbb{P}^{-1}\to
  \mathbb{P}^0}$ is a projective resolution of $M$, and it computes
  $\operatorname{Ext}^a_A(M,N)$ as the maps ${C^{-a}(\mathbb{P})\to N}$ modulo
  those that extend to $\mathbb{P}^{-a+1}$; by the extension property, these
  are exactly the maps factoring through a projective module. Applying the
  same description with $N$ replaced by a projective module $L$ gives
  ${\operatorname{Ext}^a_A(M,L)=0}$, since every map ${C^{-a}(\mathbb{P})\to L}$
  extends.
\end{proof}

Now let $E$ be a finite-dimensional symmetric algebra. Then
${E\cong DE}$ as bimodules, so finitely generated projective $E$-modules are
injective, and every finitely generated module $M$ embeds into a finitely
generated projective module, by dualising a projective cover of the right
module $DM$. Splicing a projective resolution of $M$ with a coresolution by
projective modules gives an exact complex $\mathbb{P}$ of finitely generated
projective modules with ${C^0(\mathbb{P})=M}$, and $\mathbb{P}$ is totally
acyclic, since $\Hom[E]{-}{L}$ is exact for projective $L$. By
\Cref{lemma:complete-resolutions}, two such complexes for $M$ are homotopy
equivalent, the functor $C^0$ identifies the homotopy category of these
complexes with the stable category $\operatorname{\underline{mod}}E$, the
shift $[1]$ corresponds to the cosyzygy by~\eqref{eq:cokernel-sequences}, and
\[
  \widehat{\operatorname{Ext}}{}^a_E(M,N)=H^a\Hom[E]{\mathbb{P}_M}
  {\mathbb{P}_N}
\]
for complete resolutions $\mathbb{P}_M$ and $\mathbb{P}_N$. Let $F$ be a
finitely generated $E$-bimodule that is projective as a left and as a right
module, and write ${FM=F\otimes_EM}$. The functor ${F\otimes_E-}$ is exact, since
$F_E$ is projective, and preserves projective modules, since ${}_EF$ is
projective; hence $F\mathbb{P}$ is a complete resolution of $FM$. In particular,
${F\otimes_E-}$ induces an exact functor on
$\operatorname{\underline{mod}}E$ with ${F(M[a])\cong(FM)[a]}$, and for a stable
morphism ${g\colon M\to N[a]}$ we write ${F(g)\colon FM\to(FN)[a]}$.

\subsection{Triangular algebras}
\label{subsec:triangular}

For $E$ and $F$ as above, let
\[
  \Lambda=\begin{pmatrix}E&0\\F&E\end{pmatrix},\qquad
  (a,f,b)(a',f',b')=(aa',fa'+bf',bb'),
\]
where $(a,f,b)$ denotes the matrix with diagonal entries $a$, $b$ and lower
left entry $f$. The two diagonal idempotents split a left $\Lambda$-module
into components $M$ and $N$, and the action of the lower left corner gives an
$E$-linear map ${\eta\colon FM\to N}$; conversely, every such triple
$(M,N,\eta)$ is a $\Lambda$-module, with ${(a,f,b)(m,n)=(am,\eta(f\otimes
m)+bn)}$. A morphism ${(M,N,\eta)\to(M',N',\eta')}$ is a pair of $E$-linear maps
$(\alpha,\beta)$ with ${\beta\eta=\eta'F(\alpha)}$, and a sequence of triples is
exact if and only if both component sequences are exact. Define exact functors
\[
  \mathsf{L}_1(M)=(M,FM,\id),\qquad\mathsf{L}_2(N)=(0,N,0).
\]

\begin{lemma}
  \label{lemma:triangular}
  In the situation above, the following hold.
  \begin{enumerate}
    \item\label{it:triangular-hom} There are natural isomorphisms
      ${\Hom[\Lambda]{\mathsf{L}_iM}{\mathsf{L}_iN}}\cong{\Hom[E]{M}{N}}$ for
      ${i=1,2}$, ${\Hom[\Lambda]{\mathsf{L}_2M}{\mathsf{L}_1N}}\cong
      {\Hom[E]{M}{FN}}$ and ${\Hom[\Lambda]{\mathsf{L}_1M}{\mathsf{L}_2N}}={0}$.
    \item\label{it:triangular-proj} We have ${\Lambda=\mathsf{L}_1(E)\oplus
      \mathsf{L}_2(E)}$, so $\mathsf{L}_1$ and $\mathsf{L}_2$ preserve finitely
      generated projective modules, and the first component of every
      finitely generated projective $\Lambda$-module is a projective
      $E$-module.
    \item\label{it:triangular-acyclic} If $\mathbb{P}$ is a totally acyclic
      complex of $E$-modules, then $\mathsf{L}_1\mathbb{P}$ and
      $\mathsf{L}_2\mathbb{P}$ are totally acyclic complexes of
      $\Lambda$-modules.
  \end{enumerate}
\end{lemma}

\begin{proof}
  \eqref{it:triangular-hom} A morphism ${\mathsf{L}_1M\to(M',N',\eta')}$ is
  determined by its first component $\alpha$, the second being
  $\eta'F(\alpha)$, and a morphism ${\mathsf{L}_2M\to(M',N',\eta')}$ is an
  arbitrary map ${M\to N'}$. This gives the first three isomorphisms. For the
  last, the compatibility condition forces the second component to vanish,
  since the structure map of $\mathsf{L}_1M$ is the identity.

  \eqref{it:triangular-proj} The first column of $\Lambda$ is
  ${(E,F,\id)=\mathsf{L}_1(E)}$ and the second is ${(0,E,0)=\mathsf{L}_2(E)}$.
  The first component of a direct summand of $\Lambda^r$ is a direct summand of
  $E^r$.

  \eqref{it:triangular-acyclic} The terms are finitely generated projective by
  \eqref{it:triangular-proj}, and both complexes are exact. By
  \eqref{it:triangular-hom} and \eqref{it:triangular-proj},
  \begin{gather*}
    \Hom[\Lambda]{\mathsf{L}_1\mathbb{P}}{\Lambda}\cong\Hom[E]{\mathbb{P}}{E},\\
    \Hom[\Lambda]{\mathsf{L}_2\mathbb{P}}{\Lambda}\cong
    \Hom[E]{\mathbb{P}}{F}\oplus\Hom[E]{\mathbb{P}}{E},
  \end{gather*}
  which are exact since $\mathbb{P}$ is totally acyclic and ${}_EF$ is
  projective.
\end{proof}

\subsection{A module with no positive self-extensions}
\label{subsec:conversion-theorem}

\begin{theorem}
  \label{thm:conversion}
  Let $E$ be a finite-dimensional symmetric algebra over a field $\kk$, let $F$
  be a finite-dimensional $E$-bimodule that is projective as a left and as a
  right module, let $S$ be a finite-dimensional left $E$-module, and let
  ${v\colon S\to FS}$ be a morphism in $\operatorname{\underline{mod}}E$. For
  ${a\in\ZZ}$ consider the linear map
  \[
    \delta^a\colon\widehat{\operatorname{Ext}}{}^a_E(S,S)^{\oplus2}
    \longrightarrow\widehat{\operatorname{Ext}}{}^a_E(S,FS),\qquad
    (g,j)\longmapsto F(g)v-v[a]j.
  \]
  Suppose that $\delta^0$ is surjective with nonzero kernel and that $\delta^a$
  is bijective for every ${a>0}$. Choose a homomorphism ${v_0\colon S\to FS}$
  representing $v$ and an injective homomorphism ${i\colon S\to Q}$ into a
  finitely generated projective module $Q$, and put
  \[
    \begin{gathered}
    Z_2=(FS\oplus Q)/\set{(v_0(s),i(s))}[s\in S],\\
    \iota\colon FS\to Z_2,\ t\mapsto[t,0],\qquad Z=(S,Z_2,\iota).
    \end{gathered}
  \]
  Then $Z$ is a finite-dimensional, nonprojective and Gorenstein-projective
  $\Lambda$-module, and
  \[
    \operatorname{Ext}^a_\Lambda(Z,Z)=0=\operatorname{Ext}^a_\Lambda(Z,\Lambda)
    \qquad\text{for all }a>0.
  \]
  In particular, $Z$ is a counterexample to the Auslander--Reiten conjecture
  for~$\Lambda$.
\end{theorem}

\begin{proof}
  \emph{The complete resolution.} Choose a complete resolution $\mathbb{P}$ of
  $S$ over $E$ with ${\mathbb{P}^1=Q}$ and ${d^0_{\mathbb{P}}=i\epsilon_0}$, where
  ${\epsilon_0\colon\mathbb{P}^0\to S}$ is the projection; this is possible by
  starting the coresolution of $S$ with~$i$. By
  \Cref{lemma:complete-resolutions} there is a chain map ${f\colon\mathbb{P}\to
  F\mathbb{P}}$ inducing $v_0$ on $C^0$. By \Cref{lemma:triangular}, $f$ defines
  a chain map ${V\colon\mathsf{L}_2\mathbb{P}\to\mathsf{L}_1\mathbb{P}}$. Let
  $\mathbb{P}_Z$ be its mapping cone:
  \[
    \mathbb{P}_Z^n=\mathsf{L}_1\mathbb{P}^n\oplus\mathsf{L}_2\mathbb{P}^{n+1},
    \qquad d_{\mathbb{P}_Z}=\begin{pmatrix}\mathsf{L}_1(d)&V\\0&
    -\mathsf{L}_2(d)\end{pmatrix}.
  \]
  The sequence
  \[
    0\longrightarrow\mathsf{L}_1\mathbb{P}\longrightarrow\mathbb{P}_Z
    \longrightarrow\mathsf{L}_2\mathbb{P}[1]\longrightarrow0
  \]
  is split in each degree, and its outer terms
  are totally acyclic by \Cref{lemma:triangular}; the long exact cohomology
  sequences of this sequence and of its dual show that $\mathbb{P}_Z$ is
  totally acyclic.

  \emph{Its cokernel.} The first component of $C^0(\mathbb{P}_Z)$ is
  ${C^0(\mathbb{P})=S}$. Its second component is the cokernel of
  \[
    F\mathbb{P}^{-1}\oplus\mathbb{P}^0\longrightarrow F\mathbb{P}^0\oplus Q,
    \qquad(x,y)\longmapsto(F(d)x+f^0y,-i\epsilon_0y),
  \]
  which is ${(FS\oplus Q)/\set{(v_0(s),-i(s))}[s\in S]}$, since ${F\otimes_E-}$ is
  exact and $f$ induces $v_0$; the structure map is induced by the inclusion of
  $FS$. The automorphism ${(t,q)\mapsto(t,-q)}$ of ${FS\oplus Q}$ induces an
  isomorphism ${C^0(\mathbb{P}_Z)\cong Z}$. Hence, $Z$ is Gorenstein-projective,
  and ${\operatorname{Ext}^a_\Lambda(Z,\Lambda)=0}$ for ${a>0}$ by
  \Cref{lemma:complete-resolutions}. The module $Z$ is finite-dimensional
  since $S$, $FS$ and $Q$ are.

  \emph{The endomorphism complex.} Put ${\mathcal{H}=\Hom[E]{\mathbb{P}}
  {\mathbb{P}}}$ and ${\mathcal{K}=\Hom[E]{\mathbb{P}}{F\mathbb{P}}}$, so that
  ${H^a(\mathcal{H})=\widehat{\operatorname{Ext}}{}^a_E(S,S)}$ and
  ${H^a(\mathcal{K})=\widehat{\operatorname{Ext}}{}^a_E(S,FS)}$, and let
  $\boldsymbol{\delta}$ be the chain map
  \[
    \boldsymbol{\delta}\colon\mathcal{H}^{\oplus2}\longrightarrow\mathcal{K},
    \qquad(g,j)\longmapsto F(g)f-fj;
  \]
  it induces $\delta^a$ in cohomology. Since there
  are no nonzero maps from $\mathsf{L}_1$-terms to $\mathsf{L}_2$-terms, an
  element $\Phi$ of degree $n$ of $\Hom[\Lambda]{\mathbb{P}_Z}{\mathbb{P}_Z}$
  has unique coordinates ${(g,j,h)\in\mathcal{H}^n\oplus\mathcal{H}^n\oplus
  \mathcal{K}^{n-1}}$, defined by
  \[
    \Phi=\begin{pmatrix}\mathsf{L}_1(g)&(-1)^nh\\0&(-1)^n\mathsf{L}_2(j)
    \end{pmatrix},
  \]
  where $h$ stands for the map ${\mathsf{L}_2\mathbb{P}\to\mathsf{L}_1
  \mathbb{P}}$ corresponding to it under \Cref{lemma:triangular}. Expanding
  \[
    \partial\Phi=d_{\mathbb{P}_Z}\Phi-(-1)^n\Phi d_{\mathbb{P}_Z},
  \]
  the diagonal
  entries are $\mathsf{L}_1(\partial g)$ and ${(-1)^{n+1}\mathsf{L}_2(\partial
  j)}$, and the upper right entry is
  \[
    (-1)^n\bigl(dh+(-1)^nhd+fj-F(g)f\bigr)
    =(-1)^{n+1}\bigl(\boldsymbol{\delta}(g,j)-\partial h\bigr).
  \]
  Hence,
  \[
    \partial(g,j,h)=\bigl(\partial g,\partial j,\boldsymbol{\delta}(g,j)-
    \partial h\bigr),
  \]
  and there is a degreewise split short exact sequence of complexes
  \[
    0\longrightarrow\mathcal{K}[-1]\longrightarrow
    \Hom[\Lambda]{\mathbb{P}_Z}{\mathbb{P}_Z}\longrightarrow
    \mathcal{H}^{\oplus2}\longrightarrow0,
  \]
  given by ${h\mapsto(0,0,h)}$ and ${(g,j,h)\mapsto(g,j)}$. A cocycle $(g,j)$
  lifts to $(g,j,0)$, whose boundary is $(0,0,\boldsymbol{\delta}(g,j))$, so
  the connecting homomorphism is $\delta^a$, and for every $a$ there is an exact
  sequence
  \begin{equation}
    \label{eq:conversion-sequence}
    0\longrightarrow\operatorname{coker}\delta^{a-1}\longrightarrow
    H^a\Hom[\Lambda]{\mathbb{P}_Z}{\mathbb{P}_Z}\longrightarrow
    \ker\delta^a\longrightarrow0.
  \end{equation}
  For ${a>0}$ the middle term is $\operatorname{Ext}^a_\Lambda(Z,Z)$ by
  \Cref{lemma:complete-resolutions}. For ${a=1}$ the outer terms vanish since
  $\delta^0$ is surjective and $\delta^1$ is injective, and for ${a\geq2}$ since
  $\delta^{a-1}$ and $\delta^a$ are bijective.

  \emph{Nonprojectivity.} Since ${\ker\delta^0\neq0}$, the stable endomorphism
  ring of $S$ is nonzero, so $S$ is not projective. The first component of $Z$
  is $S$, so $Z$ is not projective by \Cref{lemma:triangular}.
\end{proof}

\begin{remark}
  \label{rem:conversion-mechanism}
  The proof uses the symmetry of $E$ only to ensure that finitely generated
  projective modules are injective and that every module embeds into a
  projective one; it uses neither Tate duality nor the simplicity of~$S$. The
  sequence~\eqref{eq:conversion-sequence} holds without the hypotheses on
  $\delta$. In particular, if $\delta^1$ is injective, then
  ${\operatorname{Ext}^1_\Lambda(Z,Z)\cong\operatorname{coker}\delta^0}$. An
  application of \Cref{thm:conversion} therefore needs a bimodule $F$ for
  which $\widehat{\operatorname{Ext}}{}^a_E(S,FS)$ is a copy of
  $\widehat{\operatorname{Ext}}{}^a_E(S,S)^{\oplus2}$ for ${a>0}$, compatibly
  with the two actions of the Tate algebra of $S$ in the definition of
  $\delta^a$, while in degree zero $\delta^0$ is surjective but not injective;
  in the construction below, $\widehat{\operatorname{Ext}}{}^0_E(S,FS)$ is
  one-dimensional. The bimodule of \Cref{sec:ar-route} is obtained as a
  fibre of a map from two twisted copies of $E$ to a bimodule whose evaluation
  at $S$ has Tate cohomology only in degrees $0$ and $3$.
\end{remark}

\section{A counterexample to the Auslander--Reiten conjecture}
\label{sec:ar-route}

Following~\cite{OAI26ar}, we construct a bimodule and a stable morphism
satisfying the hypotheses of \Cref{thm:conversion}, and deduce a second algebra
of infinite little finitistic dimension, with nine simple modules. Throughout,
\[
  \kk=\mathbb{F}_2(q,H_1,H_2),
\]
the field of rational functions in three variables over $\mathbb{F}_2$; in
characteristic two all signs disappear, and ${\delta^a(g,j)=F(g)v+v[a]j}$ in
\Cref{thm:conversion}. The construction uses a ten-dimensional algebra $C$ with
a simple module $s$ whose resolution has differentials depending on successive
powers of $q$ (\Cref{subsec:algebra-C}); its trivial extension ${T=C\ltimes DC}$,
over which $s$ has polynomial Ext algebra on a generator of degree three
(\Cref{subsec:algebra-T}), realised by an explicit Hochschild cocycle
(\Cref{subsec:cocycle}); the tensor square ${E=T\otimes_\kk T}$, a bimodule
$\mathcal{C}$ built from two cones on that cocycle, and twisted lifts
(\Cref{subsec:tensor-square,subsec:two-cones,subsec:lifts}); and the fibre $F$
of a map from two twisted copies of $E$ to $\mathcal{C}[3]$
(\Cref{subsec:fibre}). Finite identities, such as associativity of $C$ and
the cocycle identities, are checked by exact computations recorded in
\Cref{app:computations}; every statement about all degrees is proved in the
text.

\subsection{The algebra \texorpdfstring{$C$}{C}}
\label{subsec:algebra-C}

Let $C$ be the $\kk$-algebra with basis
\[
  e,x,y,z\in eCe,\qquad u,v\in eCf,\qquad t,j\in fCe,\qquad f,n\in fCf,
\]
where $e$ and $f$ are orthogonal idempotents with ${e+f=1}$, multiplication by
$e$ and $f$ is prescribed by the corners, and the products of two of the
remaining eight basis vectors are zero except for
\[
  xy=qz,\quad yx=z,\quad xu=yu=v,\quad tx=j,\quad ty=q^2j,\quad ut=y+qx,
\]
\[
  vt=qz,\quad uj=z,\quad tu=(1+q)n,\quad nt=qj,\quad un=v.
\]
Let $N$ be the span of $x,y,z,u,v,t,j,n$, let $s$ be the simple left module
${C/(N+Ce)}$ at $f$, and let $s^{\mathrm{r}}$ be the simple right module at~$f$.

\begin{lemma}
  \label{lemma:algebra-C}
  The algebra $C$ is associative. It is graded by
  ${\deg u=\deg t=1}$, ${\deg x=\deg y=\deg n=2}$, ${\deg v=\deg j=3}$, ${\deg z=4}$
  and ${\deg e=\deg f=0}$, its radical is $N$, and ${N^5=0}$.
\end{lemma}

\begin{proof}
  Every product above is homogeneous for the grading. A product of three
  basis vectors vanishes under both bracketings if one of them is an
  idempotent that does not match the corners, or if its total degree exceeds
  four, since there is no basis vector of degree larger than four. Among the
  remaining triples of radical basis vectors, the only ones with compatible
  corners are $xut$, $yut$, $utx$, $uty$, $utu$, $unt$, $txu$, $tyu$, $tut$,
  $tun$ and $ntu$, and their products under both bracketings are $qz$, $qz$,
  $z$, $q^2z$, ${(1+q)v}$, $qz$, $0$, $0$, ${q(1+q)j}$, $0$ and $0$; for instance
  ${(ut)y=(y+qx)y=q^2z=u(q^2j)=u(ty)}$. The associators of all $1000$ triples
  of basis vectors are also checked by an exact computation
  (\Cref{app:computations}). The span $N$ of the basis vectors of positive
  degree is a nilpotent ideal with ${N^5=0}$ and ${C/N\cong\kk\times\kk}$, so $N$
  is the radical.
\end{proof}

For ${i\geq0}$ put ${\ell_i=x+q^iy\in eCe}$. The left modules $Ce$ and $Cf$ have
the bases $e,x,y,z,t,j$ and $f,u,v,n$, and right multiplication by $u$ and by
$\ell_i$ acts as follows:
\begin{equation}
  \label{eq:C-right-multiplication}
  \begin{array}{c|cccccc}
    a&e&x&y&z&t&j\\\hline
    au&u&v&v&0&(1+q)n&0\\
    a\ell_i&\ell_i&q^{i+1}z&z&0&(1+q^{i+2})j&0
  \end{array}
\end{equation}

\begin{proposition}
  \label{prop:C-resolution}
  The complex
  \[
    \mathbb{R}\colon\quad\cdots\longrightarrow Ce\xrightarrow{\cdot\ell_2}Ce
    \xrightarrow{\cdot\ell_1}Ce\xrightarrow{\cdot\ell_0}Ce
    \xrightarrow{\cdot u}Cf,
  \]
  with $Cf$ in degree zero and right multiplications as differentials, is a
  minimal projective resolution of $s$. Consequently,
  ${\operatorname{Ext}^a_C(s,s)=0}$ for ${a>0}$, and ${\RHom[C]{s}{C}\cong
  s^{\mathrm{r}}[-2]}$ in the derived category of right $C$-modules.
\end{proposition}

\begin{proof}
  Since $q$ is transcendental over $\mathbb{F}_2$, the elements ${1+q}$ and
  ${1+q^{i+2}}$ are nonzero. By~\eqref{eq:C-right-multiplication}, the image of
  ${\cdot u}$ is the span of $u$, $v$ and $n$, which is the radical of $Cf$, and
  its kernel is spanned by $\ell_0$, $z$ and $j$. A vector in the kernel of
  ${\cdot\ell_i}$ has zero coefficients at $e$ and $t$, and coefficients at $x$
  and $y$ related by ${a_y=q^{i+1}a_x}$; hence
  \[
    \ker(\cdot\ell_i)=\langle\ell_{i+1},z,j\rangle,\qquad
    \operatorname{im}(\cdot\ell_i)=\langle\ell_i,z,j\rangle,
  \]
  for every ${i\geq0}$, and the complex is exact. All differentials have image
  in $N$ times their targets, so each surjection onto a syzygy is a projective
  cover, and the resolution is minimal. Since ${\Hom[C]{Ce}{s}=es=0}$ and
  ${\Hom[C]{Cf}{s}=\kk}$, the asserted vanishing of self-extensions follows.

  Identifying ${\Hom[C]{Ce}{C}=eC}$ and ${\Hom[C]{Cf}{C}=fC}$ as right modules,
  precomposition with right multiplication by $b$ becomes left multiplication
  by $b$, and $\Hom[C]{\mathbb{R}}{C}$ is the complex
  \[
    fC\xrightarrow{u\cdot}eC\xrightarrow{\ell_0\cdot}eC\xrightarrow{\ell_1\cdot}
    eC\longrightarrow\cdots
  \]
  in degrees $0,1,2,\dots$, with ${fC=\langle f,t,j,n\rangle}$ and
  ${eC=\langle e,x,y,z,u,v\rangle}$. Left multiplication by $u$ sends
  $f,t,j,n$ to the independent vectors ${u,y+qx,z,v}$, and left multiplication
  by $\ell_i$ sends $e,x,y,z,u,v$ to ${\ell_i,q^iz,qz,0,(1+q^i)v,0}$. For ${i=0}$
  its kernel is ${\langle qx+y,z,u,v\rangle}$, the image of ${u\cdot}$, and its
  image is $\langle\ell_0,z\rangle$; for ${i\geq1}$ its kernel is
  $\langle\ell_{i-1},z,v\rangle$ and its image is $\langle\ell_i,z,v\rangle$.
  Hence, the only cohomology is in degree two, spanned by the class of $v$,
  on which $f$ acts as the identity and the radical acts as zero, since
  ${vt=qz}$ is a boundary. Good truncations identify the complex with
  ${s^{\mathrm{r}}[-2]}$.
\end{proof}

The parameter $q$ makes the resolution aperiodic, while the self-extensions of
$s$ vanish in positive degrees: on $e\langle\ell_i,z,j\rangle$ the elements $x$
and $y$ act with ${x=q^{i+1}y}$ and ${y\neq0}$, a ratio preserved by isomorphisms of
$C$-modules, so the syzygies $\langle\ell_i,z,j\rangle$ are pairwise
nonisomorphic. The module $s$ is not a counterexample to the Auslander--Reiten
conjecture over $C$, since ${\operatorname{Ext}^2_C(s,C)\neq0}$; over ${T=C\ltimes
DC}$, its Ext algebra is polynomial on a generator of degree three
(\Cref{thm:T-ext}).

\subsection{The trivial extension \texorpdfstring{$T$}{T}}
\label{subsec:algebra-T}

Let ${T=C\ltimes DC}$, where the bimodule structure of $DC$ is
${(a\phi)(c)=\phi(ca)}$ and ${(\phi a)(c)=\phi(ac)}$, and regard $s$ as a
$T$-module through ${T\to C}$. For a basis vector $c$ of $C$ we write $c^*$ for
the dual basis vector of $DC$.

\begin{lemma}
  \label{lemma:T-symmetric}
  The algebra $T$ is a twenty-dimensional symmetric algebra, with symmetrising
  form ${\operatorname{tr}(a,\phi)=\phi(1)}$. Its radical is ${N\oplus DC}$, with
  sixth power zero, and ${T/\operatorname{rad}T\cong\kk\times\kk}$.
\end{lemma}

\begin{proof}
  The algebra $T$ is symmetric, with the stated form, by
  \Cref{subsec:trivial-extensions}. With ${\deg c^*=5-\deg c}$, the algebra $T$ is
  graded, with all basis vectors
  other than $e$ and $f$ in degrees between one and five; their span is a
  nilpotent ideal with sixth power zero and quotient ${\kk\times\kk}$.
\end{proof}

\begin{theorem}
  \label{thm:T-ext}
  There is ${\tau\in\operatorname{Ext}^3_T(s,s)}$ such that
  ${\operatorname{Ext}^*_T(s,s)=\kk[\tau]}$, a polynomial algebra under the
  Yoneda product. Every nonzero element of $\operatorname{Ext}^3_T(s,s)$ is
  such a generator. Moreover, $s$ is not projective and
  ${\underline{\operatorname{End}}_T(s)=\kk}$.
\end{theorem}

\begin{proof}
  Let $\mathbb{R}$ be the resolution of \Cref{prop:C-resolution}. For a
  finitely generated projective left $C$-module $V$, the map
  ${DC\otimes_CV\to D\Hom[C]{V}{C}}$, ${\phi\otimes v\mapsto(g\mapsto\phi(g(v)))}$,
  is a natural isomorphism of left $C$-modules, as one sees for ${V=C}$. Hence,
  ${DC\otimes_C\mathbb{R}\cong D\Hom[C]{\mathbb{R}}{C}}$, which has cohomology
  ${D(s^{\mathrm{r}})=s}$ in degree ${-2}$ by \Cref{prop:C-resolution}; since $DC$
  acts on it as zero, it is quasi-isomorphic to $s[2]$ as a complex of
  $T$-modules. Tensoring the sequence ${0\to DC\to T\to C\to0}$ of
  $(T,C)$-bimodules with the complex $\mathbb{R}$ of projective, hence flat,
  $C$-modules gives a short exact sequence of complexes of $T$-modules, and
  hence a triangle
  \begin{equation}
    \label{eq:T-triangle}
    s[2]\longrightarrow\mathcal{I}\longrightarrow s\xrightarrow{\ \eta\ }s[3],
    \qquad\mathcal{I}=T\otimes_C\mathbb{R},
  \end{equation}
  in ${\DerCat[-]{\operatorname{mod}T}}$. The terms ${T\otimes_CCe=Te}$ and
  ${T\otimes_CCf=Tf}$ of $\mathcal{I}$ are projective, and $\mathcal{I}$ is
  bounded above, so adjunction gives
  ${\Hom{\mathcal{I}}{s[a]}\cong H^a\Hom[C]{\mathbb{R}}{s}}$, which is $\kk$ for
  ${a=0}$ and zero otherwise.

  Put ${V^a=\Hom{s}{s[a]}}$, which is $\operatorname{Ext}^a_T(s,s)$ for ${a\geq0}$
  and zero for ${a<0}$. Applying $\Hom{-}{s[a]}$ to~\eqref{eq:T-triangle} gives
  exact sequences
  \[
    \Hom{\mathcal{I}}{s[a-1]}\longrightarrow V^{a-3}
    \xrightarrow{\ \beta\mapsto\beta[3]\circ\eta\ }V^a\longrightarrow
    \Hom{\mathcal{I}}{s[a]}.
  \]
  For ${a=1}$ this gives ${V^1=0}$, and for ${a\geq2}$ the outer terms vanish, so
  that composition with $\eta$ is an isomorphism ${V^{a-3}\to V^a}$. Since
  ${V^0=\kk}$ and ${V^{-1}=V^{-2}=0}$, we obtain ${V^{3m}=\kk\eta^m}$ for ${m\geq0}$
  and ${V^a=0}$ in all other degrees, and no power of $\eta$ vanishes. Hence,
  ${\operatorname{Ext}^*_T(s,s)=\kk[\eta]}$, and every nonzero multiple of $\eta$
  is a generator. Since ${\operatorname{Ext}^3_T(s,s)\neq0}$, the module $s$ is
  not projective, and a nonzero endomorphism of the one-dimensional module $s$
  factoring through a projective module would make $s$ a direct summand of a
  projective module; hence ${\underline{\operatorname{End}}_T(s)=\kk}$.
\end{proof}

\subsection{A Hochschild cocycle}
\label{subsec:cocycle}

The construction of the bimodule $F$ needs a representative of $\tau$ by a
bimodule map, which the abstract computation of \Cref{thm:T-ext} does not
provide. Let ${I=\kk e\oplus\kk f\subseteq T}$ and
${\mathfrak{r}=\operatorname{rad}
T}$, so that ${T=I\oplus\mathfrak{r}}$. The \emph{bar resolution} of $T$ relative
to $I$ is the complex $\mathbb{B}$ of $T$-bimodules with
\[
  \mathbb{B}^{-n}=T\otimes_I\mathfrak{r}^{\otimes_In}\otimes_IT\quad(n\geq0),
  \qquad\mathbb{B}^n=0\quad(n>0),
\]
whose elements we write $a_0[a_1|\cdots|a_n]a_{n+1}$, with the differential
given by the sum of the ${n+1}$ products of adjacent entries, and augmentation
${\mathbb{B}^0\to T}$ given by multiplication. Its terms are finite direct sums
of the projective bimodules ${Ti\otimes_\kk jT}$ for ${i,j\in\set{e,f}}$. The map
\[
  a_0[a_1|\cdots|a_n]a_{n+1}\longmapsto[\overline{a}_0|a_1|\cdots|a_n]a_{n+1},
\]
where $\overline{a}_0$ is the component of $a_0$ in $\mathfrak{r}$, together
with ${a\mapsto1[\,]a}$ in augmented degree, is a contraction of the augmented
complex, and it is right linear; hence $\mathbb{B}$ is a projective bimodule
resolution of $T$, and ${\mathbb{B}\otimes_Ts}$ is a projective resolution of
$s$.

Let ${p\colon\mathfrak{r}^{\otimes_I3}\to T}$ be the $I$-bilinear map given by
the table of $179$ coefficients in \Cref{app:computations}, reproduced from the
appendix of~\cite{OAI26ar}: an entry $abcw$ in the row $q^d$ means that the
coefficient of the basis vector $w$ in $p(a,b,c)$ is $q^d$, where capital
letters denote dual basis vectors, so that for example ${J=j^*}$; all other
coefficients are zero. For ${\lambda\in\kk^\times}$ let $h_\lambda$ be the
automorphism ${(c,\phi)\mapsto(c,\lambda\phi)}$ of $T$, and let
${z_\lambda\colon\mathfrak{r}\to\mathfrak{r}}$ be the linear map given on basis
vectors by
\[
  z_\lambda(a)=\begin{cases}
    \lambda qa&\text{if }a\in\set{u,v,t,e^*,x^*,y^*,z^*,u^*,v^*},\\
    0&\text{otherwise}.
  \end{cases}
\]

\begin{lemma}
  \label{lemma:cocycle}
  The map $p$ has the following properties.
  \begin{enumerate}
    \item\label{it:cocycle} It is a Hochschild cocycle: for all
      ${a,b,c,d\in\mathfrak{r}}$,
      \[
        ap(b,c,d)+p(ab,c,d)+p(a,bc,d)+p(a,b,cd)+p(a,b,c)d=0,
      \]
      where $p$ is extended by zero to inputs with a factor in $I$.
    \item\label{it:weight} Every nonzero coefficient of $p(a,b,c)$ at a basis
      vector $w$, for basis vectors $a,b,c$, has exactly one more dual basis
      vector among $a,b,c$ than $w$ has.
    \item\label{it:boundary} For all ${a,b\in\mathfrak{r}}$ and
      ${\lambda\in\kk^\times}$,
      \[
        \sum_{w}p(a,b,h_\lambda(w))w^*=az_\lambda(b)+z_\lambda(ab)+
        z_\lambda(a)h_\lambda^{-1}(b),
      \]
      where $w$ runs through the radical basis vectors and $w^*$ denotes the
      dual of $w$ under the trace, so that ${(c^*)^*=c}$.
    \item\label{it:evaluation} ${p(t,x,j^*)=0}$ and ${p(t,y,j^*)=q^3f}$.
  \end{enumerate}
\end{lemma}

\begin{proof}
  These are finitely many polynomial identities in $q$ and $\lambda$, and
  they are checked by an exact computation over $\mathbb{F}_2[q,\lambda]$ on
  all inputs; see \Cref{app:computations}.
\end{proof}

Define a map ${p\colon\mathbb{B}\to\mathbb{B}[3]}$ by
\[
  p(a_0[a_1|\cdots|a_n]a_{n+1})=a_0[a_1|\cdots|a_{n-3}]\,
  p(a_{n-2},a_{n-1},a_n)a_{n+1}\quad(n\geq3),
\]
and zero on $\mathbb{B}^{-n}$ for ${n<3}$. In ${dp+pd}$, every product of adjacent
entries among the first ${n-3}$ bars appears twice, and the remaining five terms
form the left side of \Cref{lemma:cocycle}\eqref{it:cocycle}; hence $p$ is a
chain map of bimodule complexes.

\begin{proposition}
  \label{prop:cocycle-generator}
  The class $\tau$ of the chain map ${p\otimes_Ts\colon\mathbb{B}\otimes_Ts\to
  (\mathbb{B}\otimes_Ts)[3]}$ in $\operatorname{Ext}^3_T(s,s)$ is nonzero, and
  therefore generates $\operatorname{Ext}^*_T(s,s)$.
\end{proposition}

\begin{proof}
  The class is represented by the cocycle
  ${c\colon\mathbb{B}^{-3}\otimes_Ts\to s}$ obtained by composing with the
  augmentation. The element
  ${\zeta=q^2[t|x|j^*]+[t|y|j^*]}$ of ${s^{\mathrm{r}}\otimes_T\mathbb{B}^{-3}
  \otimes_Ts}$ is a cycle: its outer terms vanish because the radical
  annihilates simple modules, and since ${tx=j}$, ${ty=q^2j}$, ${xj^*=t^*}$ and
  ${yj^*=q^2t^*}$, the four inner terms cancel in pairs. A coboundary vanishes on
  every cycle, whereas ${c(\zeta)=q^2p(t,x,j^*)+p(t,y,j^*)}$ acts on $s$ as
  ${q^3\neq0}$ by \Cref{lemma:cocycle}\eqref{it:evaluation}. Hence, ${\tau\neq0}$,
  and \Cref{thm:T-ext} applies.
\end{proof}

\begin{lemma}
  \label{lemma:twist}
  For ${\lambda\in\kk^\times}$ the functor ${T_{h_\lambda}\otimes_T-}$, where
  $T_{h_\lambda}$ is $T$ with the right action twisted by $h_\lambda$, fixes
  $s$ up to a canonical isomorphism and acts on $\operatorname{Ext}^{3m}_T
  (s,s)$ as multiplication by $\lambda^{-m}$, for every ${m\geq0}$.
\end{lemma}

\begin{proof}
  The map ${T_{h_\lambda}\otimes_TM\to{}_{h_\lambda^{-1}}M}$,
  ${a\otimes m\mapsto h_\lambda^{-1}(a)m}$, is an isomorphism, natural in the
  left module $M$, onto $M$ with the action restricted along $h_\lambda^{-1}$.
  Since $h_\lambda$ is the identity modulo $DC$, which acts as zero on $s$, the
  twisted module ${}_{h_\lambda^{-1}}s$ is $s$. Applying $h_\lambda^{-1}$ to all
  entries of ${\mathbb{B}\otimes_Ts}$ defines an isomorphism from this
  resolution to its twist which induces the identity on $s$. Under it, the
  cocycle $c$ of \Cref{prop:cocycle-generator} becomes ${c\circ h_\lambda^{-1}}$.
  The only coefficients of $p$ that survive in $c$ are those at $f$ or $e$,
  which have no dual factor, and by \Cref{lemma:cocycle}\eqref{it:weight} they
  occur only on inputs with exactly one dual factor; on those inputs
  ${c\circ h_\lambda^{-1}=\lambda^{-1}c}$. Hence, the functor sends $\tau$ to
  $\lambda^{-1}\tau$, and, being exact, it preserves Yoneda products, so it
  sends $\tau^m$ to $\lambda^{-m}\tau^m$.
\end{proof}

\subsection{The tensor square}
\label{subsec:tensor-square}

Let ${E=T\otimes_\kk T}$ and ${S=s\otimes_\kk s}$, and put
${H^a}={\widehat{\operatorname{Ext}}{}^a_E(S,S)}$.

\begin{proposition}
  \label{prop:E-ext}
  The algebras $E$ and ${E^e=E\otimes_\kk E^{\op}}$ are symmetric, and
  \[
    \operatorname{Ext}^*_E(S,S)=\kk[\tau_1,\tau_2],\qquad
    \tau_1=\tau\otimes1,\ \tau_2=1\otimes\tau,
  \]
  is a polynomial algebra on two generators of degree three. Moreover, $S$ is
  not projective and ${\underline{\operatorname{End}}_E(S)=\kk}$.
\end{proposition}

\begin{proof}
  The tensor product of the symmetrising forms of the factors is a symmetric
  associative form, nondegenerate because its Gram matrix is the Kronecker
  product of two invertible matrices; the same applies to $E^e$. If
  ${\mathbb{Q}\to s}$ is a projective resolution with finitely generated terms,
  then the total complex ${\mathbb{Q}\otimes_\kk\mathbb{Q}}$ is a projective
  resolution of $S$ by finitely generated $E$-modules, as one sees by
  splitting $\mathbb{Q}$, as a complex of vector spaces, into its cohomology
  and contractible complexes. The isomorphism of complexes
  \[
    \Hom[E]{\mathbb{Q}\otimes\mathbb{Q}}{S}\cong\Hom[T]{\mathbb{Q}}{s}\otimes
    \Hom[T]{\mathbb{Q}}{s}
  \]
  gives $\operatorname{Ext}^{3m}_E(S,S)$ the basis
  $\tau_1^i\tau_2^{m-i}$, ${0\leq i\leq m}$, and all other positive degrees
  vanish. The classes $\tau_1$ and $\tau_2$ are represented by lifts of $\tau$
  acting on different tensor factors, which commute up to a sign, hence
  commute in characteristic two, so
  their monomials are the external products ${\tau^i\otimes\tau^{m-i}}$ and
  there are no relations. The last assertion follows as in \Cref{thm:T-ext}.
\end{proof}

For ${m\geq0}$ let ${P_m\subseteq\kk[\tau_1,\tau_2]}$ be the space of homogeneous
polynomials of degree $m$ in the generators, so that ${H^{3m}=P_m}$, and put
${P_m=0}$ for ${m<0}$.

\begin{proposition}
  \label{prop:E-tate}
  The Tate groups of $S$ are
  \[
    H^{3m}=P_m,\qquad H^{-3m-1}\cong DP_m\quad(m\geq0),\qquad
    H^a=0\quad\text{otherwise},
  \]
  where the isomorphism is that of \Cref{thm:tate-duality}. For ${m\geq1}$, left
  multiplication by $\tau_i$ from $H^{-3m-1}$ to $H^{-3m+2}$ is
  the transpose ${DP_m}\to{ DP_{m-1}}$ of left multiplication by $\tau_i$ from
  $P_{m-1}$ to $P_m$; it is surjective with kernel spanned by the dual of
  $\tau_{3-i}^m$. Multiplication by $\tau_i$ on ${H^{-1}=DP_0}$ is zero.
\end{proposition}

\begin{proof}
  The first assertion follows from \Cref{prop:E-ext} and Tate duality,
  \Cref{thm:tate-duality}, with ${X=Y=S}$: ${H^{-1-a}\cong DH^a}$. The
  compatibility of the duality pairing with composition~\cite[Section~2,
  (2.8)]{Lin13}, recalled after \Cref{thm:tate-duality}, identifies
  multiplication by $\tau_i$ on the negative groups with the transpose of
  multiplication by $\tau_i$ on the positive ones. Multiplication by $\tau_i$
  from $P_{m-1}$ to $P_m$ is injective with cokernel spanned by the class of
  $\tau_{3-i}^m$, which gives the description of its transpose. On $H^{-1}$
  the product lands in ${H^2=0}$.
\end{proof}

\subsection{The bimodule \texorpdfstring{$\mathcal{C}$}{C}}
\label{subsec:two-cones}

Let $\mathcal{L}$ be the mapping cone of
${p[-3]\colon\mathbb{B}[-3]\to\mathbb{B}}$,
a complex of projective $T$-bimodules with ${\mathcal{L}^j=\mathbb{B}^{j-2}\oplus
\mathbb{B}^j}$, and let
\[
  K=\mathcal{L}\otimes_\kk\mathcal{L},
\]
a complex of projective $E$-bimodules, that is, of projective $E^e$-modules,
bounded above with finitely generated terms. The complex $\mathcal{L}$ has
cohomology $T$ in degrees $0$ and $2$ and no other cohomology, and $K$ has a
filtration with four subquotients ${\mathbb{B}\otimes\mathbb{B}}$,
${(\mathbb{B}\otimes\mathbb{B})[-2]}$ twice, and ${(\mathbb{B}\otimes\mathbb{B})
[-4]}$, where ${\mathbb{B}\otimes\mathbb{B}}$ is a projective bimodule resolution
of $E$; in particular $K$ is exact in negative degrees. For ${N\geq0}$ put
\[
  C_N=\operatorname{coker}(K^{-N-1}\to K^{-N}).
\]

\begin{lemma}
  \label{lemma:two-cones}
  For ${N\geq1}$, the bimodule $C_N$ is projective as a left and as a right
  $E$-module, and ${C_{N+1}\cong C_N[-1]}$ in $\operatorname{\underline{mod}}E^e$.
  Let $\mathcal{C}$ be a finitely generated $E$-bimodule, projective on both
  sides, representing $C_N[N]$, which does not depend on $N$. Then
  ${M\mapsto\mathcal{C}M=\mathcal{C}\otimes_EM}$ induces an exact functor on
  $\operatorname{\underline{mod}}E$, and there are triangles
  \[
    S[-3]\xrightarrow{\tau_1}S\longrightarrow\mathcal{C}_1S
    \xrightarrow{\pi_1}S[-2],\qquad
    (\mathcal{C}_1S)[-3]\xrightarrow{\tau_2}\mathcal{C}_1S\longrightarrow
    \mathcal{C}S\xrightarrow{\pi_2}(\mathcal{C}_1S)[-2]
  \]
  in $\operatorname{\underline{mod}}E$, where $\mathcal{C}_1S$ is the
  evaluation at $S$ of a bimodule, projective on both sides, representing
  \[
    \operatorname{coker}\bigl((\mathcal{L}\otimes_\kk\mathbb{B})^{-2}\to
    (\mathcal{L}\otimes_\kk\mathbb{B})^{-1}\bigr)[1],
  \]
  and the second $\tau_2$ is induced by $p$ acting on the second factor.
\end{lemma}

\begin{proof}
  The cokernel functor $C_N$ is exact on degreewise split short exact
  sequences of complexes that are exact in degree ${-N-1}$: if $b$ in the
  subcomplex becomes a boundary $dv$ in the middle, the image of $v$ in the
  quotient is a cycle, hence a boundary $d\bar{w}$; replacing $v$ by
  ${v-dw}$ for a lift $w$ of $\bar{w}$ shows that $b$ is a boundary in the
  subcomplex. Hence, $C_N$ has a filtration whose factors are the cokernels
  \[
    Z_j=\operatorname{coker}\bigl((\mathbb{B}\otimes\mathbb{B})^{-j-1}\to
    (\mathbb{B}\otimes\mathbb{B})^{-j}\bigr)
  \]
  for ${j=N}$, ${N+2}$, ${N+2}$ and ${N+4}$. The
  augmented resolution ${\mathbb{B}\otimes\mathbb{B}\to E}$ splits as a complex
  of left and as a complex of right $E$-modules, since $E$ is projective on
  each side, so the $Z_j$ are projective on each side, and so is $C_N$, since
  the filtration splits on each side. Exactness of $K$ in degree ${-N}$ gives an
  exact sequence ${0\to C_{N+1}\to K^{-N}\to C_N\to0}$ with $K^{-N}$ projective
  over $E^e$, so ${C_{N+1}\cong C_N[-1]}$.

  A cosyzygy of a bimodule $M$ that is projective on each side, taken in
  $E^e$-modules, is again projective on each side: $E^e$ is symmetric, so $M$
  embeds into a projective $E^e$-module, and since modules projective on one
  side over the symmetric algebra $E$ are injective on that side, the embedding
  splits on each side. Hence, $\mathcal{C}$ can be chosen projective on both
  sides. For such a bimodule, ${\mathcal{C}\otimes_E-}$ is exact and preserves
  projective modules, and it induces an exact functor on the stable category.
  Evaluation at $S$ of the degreewise split cone sequences, which split as
  sequences of right modules, gives the two triangles; the first comes from the
  cone in the first factor of $K$, and the second from the action of $p$ on the
  second factor, as $K$ is the cone of the map induced by ${p[-3]}$ on
  ${\mathcal{L}\otimes\mathbb{B}}$.
\end{proof}

Write ${\pi=\pi_1[-2]\circ\pi_2\colon\mathcal{C}S\to S[-4]}$ for the composite
of the two projections, and put ${W^a=\widehat{\operatorname{Ext}}{}^a_E
(S,\mathcal{C}S)}$.

\begin{proposition}
  \label{prop:two-cone-profile}
  For every ${a\in\ZZ}$,
  \[
    W^a\cong\begin{cases}\kk&a\in\set{0,3},\\0&\text{otherwise,}\end{cases}
  \]
  and $\pi[3]$ induces an isomorphism ${W^3\to H^{-1}}$.
\end{proposition}

\begin{proof}
  Put ${U^a=\widehat{\operatorname{Ext}}{}^a_E(S,\mathcal{C}_1S)}$. The first
  triangle of \Cref{lemma:two-cones} gives exact sequences
  \[
    0\longrightarrow\operatorname{coker}(\tau_1\colon H^{a-3}\to H^a)
    \longrightarrow U^a\xrightarrow{\pi_1}\ker(\tau_1\colon H^{a-2}\to
    H^{a+1})\longrightarrow0.
  \]
  By \Cref{prop:E-ext,prop:E-tate}, multiplication by $\tau_1$ is injective
  on $\bigoplus_{m\geq0}P_m$ with cokernel $\kk[\tau_2]$, and on the negative
  part it is surjective with kernel spanned by the duals
  ${\varepsilon_m\in H^{-3m-1}}$ of $\tau_2^m$. Since the two possible
  contributions lie in degrees that are incongruent modulo three,
  \[
    U^{3m}=\kk x_m,\qquad U^{1-3m}=\kk y_m\quad(m\geq0),\qquad
    U^a=0\quad\text{otherwise},
  \]
  where $x_m$ is the image of $\tau_2^m$ and ${\pi_1(y_m)=\varepsilon_m}$. The
  action of $\tau_2$ on the second tensor factor commutes with all maps of the
  first triangle,
  so ${\tau_2x_m=x_{m+1}}$, ${\tau_2y_m=y_{m-1}}$ for ${m\geq1}$ and ${\tau_2y_0=0}$.
  The second triangle gives exact sequences
  \[
    0\longrightarrow\operatorname{coker}(\tau_2\colon U^{a-3}\to U^a)
    \longrightarrow W^a\xrightarrow{\pi_2}\ker(\tau_2\colon U^{a-2}\to
    U^{a+1})\longrightarrow0.
  \]
  The cokernel is nonzero only for ${a=0}$, where it is $\kk x_0$, and the kernel
  is nonzero only for ${a-2=1}$, where it is $\kk y_0$. Hence, ${W^0\cong\kk}$,
  ${W^3\cong\kk}$ via $\pi_2$ onto ${U^1=\kk y_0}$, and ${W^a=0}$ otherwise; since
  ${\pi_1(y_0)=\varepsilon_0}$ spans $H^{-1}$, the composite $\pi[3]$ is an
  isomorphism ${W^3\to H^{-1}}$.
\end{proof}

\subsection{Twisted lifts}
\label{subsec:lifts}

For ${\lambda\in\kk^\times}$ let $E_\lambda$ be the $E$-bimodule $E$ with the
right action twisted by ${h_\lambda\otimes h_\lambda}$. By \Cref{lemma:twist}
and \Cref{prop:E-ext}, the functor ${E_\lambda\otimes_E-}$ fixes $S$ and acts on
$H^{3m}$ as multiplication by $\lambda^{-m}$, for every ${m\geq0}$. Let
${\beta_0\in H^{-1}}$ be the class of the map ${S\to\Omega S}$,
${f\otimes f\mapsto f^*\otimes f^*}$, where ${\Omega S=\operatorname{rad}(E)
(f\otimes f)}$ is the kernel of the projective cover ${E(f\otimes f)\to S}$.

\begin{proposition}
  \label{prop:lifts}
  For every ${\lambda\in\kk^\times}$ there is a morphism ${g_\lambda\colon
  E_\lambda\to\mathcal{C}[3]}$ in $\operatorname{\underline{mod}}E^e$ such that
  ${\pi[3]\circ(g_\lambda\otimes_ES)=\lambda^2\beta_0}$, and ${\beta_0\neq0}$.
\end{proposition}

\begin{proof}
  Let $\mathcal{B}$ be the basis of $T$ consisting of the basis of $C$ and its
  dual basis, and let $w^*$ denote the trace dual of ${w\in\mathcal{B}}$. The
  Casimir element ${\sum_ww\otimes w^*}$ of ${T\otimes_\kk T}$ corresponds to the
  identity under ${x\otimes y\mapsto(t\mapsto x\operatorname{tr}(yt))}$, so it
  commutes with the elements of $T$. Applying $h_\lambda$ to the first factor
  and projecting to ${T\otimes_IT=\mathbb{B}^0}$ gives
  \[
    \xi_\lambda=\sum_{w\in\mathcal{B}}h_\lambda(w)\otimes_Iw^*,\qquad
    h_\lambda(a)\xi_\lambda=\xi_\lambda a\quad(a\in T),
  \]
  so ${b(x)=x\xi_\lambda}$ defines a bimodule map ${b\colon T_{h_\lambda}\to
  \mathbb{B}^0}$. Its composite with the augmentation is $\sum_wh_\lambda(w)w^*$,
  which is zero: for a basis vector $w$ of $C$, the products $ww^*$ and $w^*w$
  are the duals of the idempotents at the two ends of $w$, and each idempotent
  is the end of an even number of basis vectors of $C$ on each side, namely
  $6$ for $e$ and $4$ for $f$. Separating the same identity according to the
  left idempotent ${r\in\set{e,f}}$ of $w$ and removing the idempotent term gives
  ${\sum_{w}h_\lambda(w)w^*=r^*}$, the sum over the radical basis vectors $w$
  with left idempotent $r$; these identities are also checked in
  \Cref{app:computations}.

  Let $\mathbb{Q}$ be $\mathbb{B}$ with the right action twisted by $h_\lambda$,
  a resolution of $T_{h_\lambda}$, and define maps out of $\mathbb{Q}$ on
  generators $[a_1|\cdots|a_n]$, extended by ${\psi(a_0[\cdots]a_{n+1})}={
  a_0\psi([\cdots])h_\lambda^{-1}(a_{n+1})}$, which is the rule for bimodule
  maps from $\mathbb{Q}$. Let ${D_b\colon\mathbb{Q}\to\mathbb{B}}$ be $b$
  composed with the augmentation in degree zero and zero elsewhere, and put
  \begin{align*}
    B_n([a_1|\cdots|a_n])&=\sum_w[a_1|\cdots|a_n|h_\lambda(w)]w^*,\\
    G_n([a_1|\cdots|a_n])&=[a_1|\cdots|a_{n-1}]z_\lambda(a_n),
  \end{align*}
  with $w$ running through the radical basis vectors and ${G_0=0}$, maps of
  degrees ${-1}$ and ${+1}$. Projecting ${a\xi_\lambda=\xi_\lambda h_\lambda^{-1}(a)}$
  to the radical in the first factor gives, for a radical $a$ with right
  idempotent $r$,
  \[
    \sum_wah_\lambda(w)\otimes_Iw^*+\sum_wh_\lambda(w)\otimes_I
    w^*h_\lambda^{-1}(a)=a\otimes_Ir^*.
  \]
  Using this identity and the previous one, all terms of ${dB+Bd-D_b}$ cancel:
  the products among the first bars appear twice, the two terms involving the
  last bar of the source add up to $[a_1|\cdots|a_n]r^*$, and the product of
  the appended bar with its coefficient gives the same element. Hence,
  ${dB+Bd=0}$ in negative degrees. In degree zero, for ${r\in\set{e,f}}$,
  \[
    dB_0(r[\,]r)=\sum_wh_\lambda(w)\otimes_Iw^*+r\otimes_Ir^*=r\xi_\lambda=
    D_b(r[\,]r),
  \]
  with $w$ running through the radical basis vectors with left idempotent $r$,
  so ${dB+Bd=D_b}$. In ${dG+Gd}$ the products among the first bars cancel, and what
  remains is
  \[
    [a_1|\cdots|a_{n-2}]\bigl(a_{n-1}z_\lambda(a_n)+z_\lambda(a_{n-1}a_n)+
    z_\lambda(a_{n-1})h_\lambda^{-1}(a_n)\bigr),
  \]
  which is $pB$ by
  \Cref{lemma:cocycle}\eqref{it:boundary}; for ${n\leq1}$ both sides vanish.
  Hence, ${pB=dG+Gd}$.

  The complex ${\mathbb{Q}\otimes_\kk\mathbb{Q}}$ resolves $E_\lambda$. Define a
  map ${\Phi\colon\mathbb{Q}\otimes\mathbb{Q}\to K[3]}$ with component ${B\otimes
  D_b}$ into the summand ${(\mathbb{B}\otimes\mathbb{B})[-4][3]}$ of the underlying
  graded bimodule, component
  ${G\otimes D_b}$ into the summand corresponding to the cone in the second tensor
  factor, and zero
  components into the other two summands. The component of ${d\Phi-\Phi d}$ in
  ${(\mathbb B\otimes\mathbb B)[-4][3]}$ is ${D_b\otimes D_b}$; in the summand
  corresponding to the second tensor factor, the component ${(dG+Gd)\otimes D_b}$
  is cancelled by the term ${pB\otimes D_b}$ of the cone differential, and the
  term coming from the first factor vanishes because ${pD_b=0}$. Hence, $\Phi$ is
  a chain map in all degrees below zero, where ${D_b\otimes D_b}$ vanishes. For
  ${n\geq4}$ both complexes are exact in degree ${-n}$, so $\Phi$ induces a map
  of cokernels ${C^{-n}(\mathbb{Q}\otimes\mathbb{Q})\to C_{n-3}}$, that is, a
  stable morphism ${E_\lambda[-n]\to\mathcal{C}[3-n]}$, and shifting by $n$ gives
  $g_\lambda$. Its composite with the projection onto ${(\mathbb B\otimes\mathbb
  B)[-4][3]}$ is the
  comparison map of the bimodule map
  \[
    E_\lambda\longrightarrow\Omega_{E^e}E\subseteq(\mathbb{B}\otimes
    \mathbb{B})^0,\qquad1\longmapsto\xi_\lambda\otimes\xi_\lambda,
  \]
  which takes values in the kernel of the augmentation, since $b$ does; the
  augmentation is a projective cover, since its source is
  ${\bigoplus_rEr\otimes_\kk rE}$ over the four primitive vertex idempotents $r$
  of $E$ and the induced map on tops is an isomorphism.

  Evaluating at $S$, all terms of ${\xi_\lambda\otimes_Ts}$ vanish except the
  one with ${w^*=f}$, for which ${h_\lambda(w)=\lambda f^*}$; hence
  ${\pi[3]\circ(g_\lambda\otimes S)}$ is the class of
  ${f\otimes f\mapsto\lambda^2(f^*\otimes f^*)}$, which is $\lambda^2\beta_0$.
  Finally, ${\beta_0\neq0}$: a map ${S\to\Omega S}$ factoring through a projective
  module factors through a free module $E^r$, its components ${S\to E}$ take
  values in the left socle of $E$, and a map ${E\to\Omega S}$ is right
  multiplication by a radical element; for $u$ in the left socle and $j$ in the
  radical we have ${\operatorname{tr}(auj)=\operatorname{tr}(jau)=0}$ for all
  $a$, so ${uj=0}$, and every such factorisation is zero, whereas
  ${f^*\otimes f^*\neq0}$.
\end{proof}

\subsection{The fibre and the counterexample}
\label{subsec:fibre}

Let ${w\in W^3}$ be the element with ${\pi[3](w)=\beta_0}$, which exists and is
nonzero by \Cref{prop:two-cone-profile,prop:lifts}. For a finitely generated
$E^e$-module $M$, put
\[
  \Sigma M=\operatorname{coker}\bigl(M\to\Hom[\kk]{E^e}{M}\bigr),
\]
where $E^e$ acts on $\Hom[\kk]{E^e}{M}$ by ${(r\varphi)(t)=\varphi(tr)}$ and
${m\mapsto(t\mapsto tm)}$ is an injective homomorphism. Its target is
projective-injective, since
\[
  \Hom[\kk]{E^e}{M}\cong D(E^e)\otimes_\kk M\cong E^e\otimes_\kk M,
\]
with $M$ regarded as a vector space. If $M$ is projective on each side, so is
$\Sigma M$, as in the proof of \Cref{lemma:two-cones}, and $\Sigma M$ represents
$M[1]$. The bimodule ${\mathcal{Y}=\Sigma^4C_1}$ represents $\mathcal{C}[3]$.
Choose bimodule homomorphisms ${g_i\colon E_{H_i}\to\mathcal{Y}}$ representing
$H_i^{-2}g_{H_i}$, for ${i=1,2}$, so that both evaluate at $S$ to $w$. Let
${Q_{\mathcal{Y}}=E^e\otimes_\kk\mathcal{Y}\to\mathcal{Y}}$ be the free cover
${r\otimes y\mapsto ry}$, and define $F$ by the exact sequence
\begin{equation}
  \label{eq:fibre}
  0\longrightarrow F\longrightarrow E_{H_1}\oplus E_{H_2}\oplus Q_{\mathcal{Y}}
  \xrightarrow{(g_1,g_2,\cdot)}\mathcal{Y}\longrightarrow0,
\end{equation}
with projections ${\rho_i\colon F\to E_{H_i}}$.

\begin{proposition}
  \label{prop:fibre}
  The bimodule $F$ is finite-dimensional and projective as a left and as a
  right $E$-module. There is a unique stable morphism ${v\colon S\to FS}$ with
  ${\rho_1v=\rho_2v=\id[S]}$, and for ${V^a=\widehat{\operatorname{Ext}}{}^a_E
  (S,FS)}$ the map
  \[
    \delta^a\colon(H^a)^{\oplus2}\longrightarrow V^a,\qquad
    (g,j)\longmapsto F(g)v+v[a]j,
  \]
  is bijective for every ${a>0}$, while $\delta^0$ is surjective with kernel
  spanned by $(1,1)$.
\end{proposition}

\begin{proof}
  All terms of~\eqref{eq:fibre} are finite-dimensional, and $\mathcal{Y}$ is
  projective on each side, so the sequence splits on each side and $F$ is
  projective on each side. Tensoring with $S$ preserves exactness, the free
  bimodule evaluates to a projective module, and ${E_{H_i}\otimes_ES\cong S}$;
  hence there is a triangle
  \[
    FS\xrightarrow{(\rho_1,\rho_2)}S\oplus S\xrightarrow{(w,w)}
    \mathcal{C}S[3]\longrightarrow FS[1]
  \]
  in $\operatorname{\underline{mod}}E$, and the $\rho_i$ induce natural
  transformations ${F\otimes_E-\to E_{H_i}\otimes_E-}$ of exact functors on the
  stable category. Applying ${\widehat{\operatorname{Ext}}{}^a_E(S,-)}$ gives
  exact sequences
  \[
    (H^{a-1})^{\oplus2}\longrightarrow W^{a+2}\longrightarrow V^a
    \xrightarrow{(\rho_1,\rho_2)}(H^a)^{\oplus2}\longrightarrow W^{a+3}.
  \]
  For ${a=0}$, since ${W^2=0}$ and ${W^3=\kk w}$, the group $V^0$ maps
  isomorphically onto the kernel of ${(c,d)\mapsto(c+d)w}$, the diagonal line;
  this gives the unique $v$, which spans $V^0$. For ${a=1}$, the map
  ${(H^0)^{\oplus2}\to W^3}$ is surjective and ${W^4=0}$, and for ${a\geq2}$ both
  $W^{a+2}$ and $W^{a+3}$ vanish; hence ${(\rho_1,\rho_2)\colon V^a\to
  (H^a)^{\oplus2}}$ is an isomorphism for all ${a>0}$.

  By naturality,
  \[
    \rho_i[a]\circ F(g)\circ v=(E_{H_i}\otimes_E g)\circ\rho_iv=
    E_{H_i}\otimes_Eg\qquad\text{and}\qquad\rho_i[a]\circ v[a]\circ j=j.
  \]
  In degree zero this
  gives ${\delta^0(c,d)=(c+d)v}$, which is surjective with kernel $\kk(1,1)$. For
  ${a=3m}$ with ${m>0}$, composing $\delta^a$ with the isomorphism $(\rho_1,\rho_2)$
  gives
  \[
    (g,j)\longmapsto(H_1^{-m}g+j,\ H_2^{-m}g+j),
  \]
  by \Cref{lemma:twist,prop:E-ext}; its scalar ${2\times2}$ matrix has determinant
  ${H_1^{-m}+H_2^{-m}\neq0}$, since
  $H_1^m$ and $H_2^m$ are distinct monomials. For positive $a$ not divisible
  by three both sides vanish.
\end{proof}

\begin{theorem}[{\cite[Theorem~1.1]{OAI26ar}}]
  \label{thm:ar-counterexample}
  Let
  \[
    \Lambda=\begin{pmatrix}E&0\\F&E\end{pmatrix},
  \]
  and let $Z$ be the $\Lambda$-module of \Cref{thm:conversion} for the data of
  \Cref{prop:fibre}. Then $\Lambda$ is a finite-dimensional algebra over
  ${\kk=\mathbb{F}_2(q,H_1,H_2)}$ with eight simple modules, all one-dimensional,
  and $Z$ is a finite-dimensional, nonprojective, Gorenstein-projective
  $\Lambda$-module with ${\operatorname{Ext}^i_\Lambda(Z,Z\oplus\Lambda)=0}$ for
  all ${i\geq1}$. The same holds after extending scalars to any field
  containing~$\kk$.
\end{theorem}

\begin{proof}
  The first assertion about $Z$ is \Cref{thm:conversion} applied to
  \Cref{prop:fibre}. The radical of $E$ is ${J_E=\operatorname{rad}T\otimes T+
  T\otimes\operatorname{rad}T}$, a sum of two commuting nilpotent ideals with
  quotient $\kk^4$. The ideal of $\Lambda$ with diagonal entries in $J_E$ and
  lower left entry in $F$ is nilpotent: a product of $2n$ of its elements has
  at most one factor contributing through the lower left corner, flanked by
  ${2n-1}$ factors from $J_E$, and vanishes if ${J_E^n=0}$. Its quotient is
  $\kk^8$, so it is the radical, and $\Lambda$ has eight one-dimensional simple
  modules.

  For a field extension $K$ of $\kk$, a projective resolution of a finitely
  generated module $M$ by finitely generated projective modules gives
  \[
    K\otimes_\kk\operatorname{Ext}^i_\Lambda(M,N)\cong
    \operatorname{Ext}^i_{\Lambda_K}(M_K,N_K),
  \]
  since Hom from a finitely
  generated projective module commutes with extension of scalars. Hence, the
  vanishing of extensions persists. A complete resolution of $Z$ remains totally
  acyclic after extension of scalars, and a nonsplit projective cover sequence
  of $Z$ remains nonsplit, since its class in $\operatorname{Ext}^1$ stays
  nonzero; so $Z_K$ is Gorenstein-projective and nonprojective. The radical
  computation is unchanged.
\end{proof}

\subsection{An endomorphism algebra of infinite finitistic dimension}
\label{subsec:ar-consequences}

\begin{corollary}
  \label{coro:ar-findim}
  Let $Z'$ be an indecomposable nonprojective direct summand of the module $Z$
  of \Cref{thm:ar-counterexample}. Then the algebra
  ${\operatorname{End}_\Lambda(\Lambda\oplus Z')}$ has nine simple modules and
  infinite little finitistic dimension. The same construction gives such an
  algebra over every field containing~$\kk$.
\end{corollary}

\begin{proof}
  The module $Z'$ exists since $Z$ is nonprojective, and it inherits the
  vanishing of extensions, since Ext is additive. By \Cref{thm:ar-to-findim},
  the algebra ${\operatorname{End}_\Lambda(\Lambda\oplus Z')}$ has infinite
  little finitistic dimension. Since ${\Lambda/\operatorname{rad}\Lambda\cong
  \kk^8}$, the algebra $\Lambda$ is basic with eight indecomposable projective
  modules, and ${\Lambda\oplus Z'}$ has nine pairwise nonisomorphic
  indecomposable summands, which index the simple modules of its endomorphism
  algebra. For a field $K$ containing $\kk$, the same argument applies to
  $\Lambda_K$ and an indecomposable nonprojective summand of $Z_K$, by
  \Cref{thm:ar-counterexample}.
\end{proof}

\begin{remark}
  \label{rem:ar-sizes}
  The algebra $\Lambda$ is explicit. With the bimodule
  representatives used above, ${\dim_\kk E=400}$ and ${\dim_\kk E^e=160000}$, each
  application of $\Sigma$ multiplies dimensions by $159999$, and
  \[
    \dim_\kk F=800+159999^5\dim_\kk C_1,\qquad
    \dim_\kk\Lambda=2\dim_\kk E+\dim_\kk F,
  \]
  where ${\dim_\kk C_1=1\,623\,889\,344}$ is computed from the dimensions of the
  terms of $K$ and its cohomology (\Cref{app:computations}); hence
  ${\dim_\kk\Lambda\approx1.7\cdot10^{35}}$. A computation with minimal bimodule
  resolutions of the cones over a finite field, recorded in
  \Cref{app:computations}, determined that the two cokernels corresponding
  to $C_0$ and $C_1$ would have dimensions $784\,704$ and $1\,377\,984$, and
  was stopped before forming them, at a linear system with about
  ${3.6\cdot10^6}$ unknowns; it gives no lower bound for the dimension of a
  minimised construction. The algebra in \Cref{coro:ar-findim} has nine simple
  modules; the number ${3(l+2)}$ for \Cref{coro:main} is not determined.
\end{remark}

\section{Obstructions to smaller constructions}
\label{sec:obstructions}

The algebra of \Cref{coro:ar-findim} has nine simple modules, and the
triangular algebra used to obtain it has eight simple modules and dimension of
the order of $10^{35}$ (\Cref{rem:ar-sizes}). Most of this size comes
from the tensor square ${E=T\otimes T}$ and from the bimodule cosyzygies
over~$E^e$. The analogous construction over $T$ itself, with a single cone,
fails in degree zero
(\Cref{prop:one-factor,coro:one-factor}), and Tate duality also forces the
vanishing of the related Toda brackets (\Cref{thm:tate-obstruction}).

\subsection{One cone and two cones}
\label{subsec:one-cone}

In this subsection $A$ is a finite-dimensional symmetric algebra and $s$ a
simple $A$-module with ${\operatorname{End}_A(s)=\kk}$ and
${\operatorname{Ext}^*_A(s,s)=\kk[\tau]}$, a polynomial algebra on a generator
of degree three; for instance ${(A,s)=(T,s)}$ of \Cref{thm:T-ext}. Write
${H^a=\widehat{\operatorname{Ext}}{}^a_A(s,s)}$. By \Cref{thm:tate-duality},
${H^{3m}=\kk\tau^m}$ and ${H^{-3m-1}\cong D(\kk\tau^m)}$ for ${m\geq0}$, all other
$H^a$ vanish, and multiplication by $\tau$ is an isomorphism
${H^{-3m-1}\to H^{-3m+2}}$ for ${m\geq1}$, by the argument of
\Cref{prop:E-tate}. Let $\beta_0$ span $H^{-1}$.

\begin{proposition}
  \label{prop:one-cone}
  Let
  \[
    s[-3]\xrightarrow{\tau}s\xrightarrow{i}X\xrightarrow{\pi}s[-2]
  \]
  be a triangle in $\operatorname{\underline{mod}}A$, and put
  ${U^a=\widehat{\operatorname{Ext}}{}^a_A(s,X)}$. Then $U^0$ and $U^1$ are
  one-dimensional, $i$ induces an isomorphism ${H^0\to U^0}$, $\pi[1]$ induces an
  isomorphism ${U^1\to H^{-1}}$, and ${U^a=0}$ for ${a\notin\set{0,1}}$.
\end{proposition}

\begin{proof}
  The triangle gives exact sequences
  \[
    0\longrightarrow\operatorname{coker}(\tau\colon H^{a-3}\to H^a)
    \xrightarrow{\ i\ }U^a\xrightarrow{\ \pi\ }\ker(\tau\colon H^{a-2}\to
    H^{a+1})\longrightarrow0.
  \]
  Multiplication by $\tau$ is an isomorphism ${H^{3m-3}\to H^{3m}}$ for ${m\geq1}$
  and ${H^{-3m-1}}\to{ H^{-3m+2}}$ for ${m\geq1}$, so the only nonzero cokernel is
  $H^0$, in degree ${a=0}$, and the only nonzero kernel is $H^{-1}$, in degree
  ${a-2=-1}$.
\end{proof}

In the construction over ${T\otimes_\kk T}$ of \Cref{sec:ar-route}, the first cone
has
one-dimensional groups in the degrees $3m$ and ${1-3m}$ for ${m\geq0}$; after the
second cone, taken in the other tensor factor, only the groups in degrees zero
and three remain (\Cref{prop:two-cone-profile}). With a single cone the two
nonzero degrees are
adjacent, and the comparison map of \Cref{thm:conversion} cannot be surjective
in degree zero unless at least one of the compositions ${w_i[-1]\circ\beta_0}$ in
\Cref{prop:one-factor} is nonzero.

\begin{proposition}
  \label{prop:one-factor}
  In the situation of \Cref{prop:one-cone}, let $F$ be a finite-dimen\-sional
  $A$-bimodule, projective on each side, with a triangle
  \[
    Fs\xrightarrow{(\rho_1,\rho_2)}s\oplus s\xrightarrow{(w_1,w_2)}X[1]
    \longrightarrow Fs[1]
  \]
  in $\operatorname{\underline{mod}}A$, where ${w_1,w_2\in U^1}$ are not both
  zero. Suppose that ${w_1[-1]\circ\beta_0}={0}={w_2[-1]\circ\beta_0}$ in $U^0$. Then
  for every stable morphism ${v\colon s\to Fs}$ the map $\delta^0$ of
  \Cref{thm:conversion} is not surjective, and the module $Z$ constructed by the
  formula in \Cref{thm:conversion}, if defined, satisfies
  ${\operatorname{Ext}^1_\Lambda(Z,Z)\neq0}$.
\end{proposition}

\begin{proof}
  Put ${V^0=\widehat{\operatorname{Ext}}{}^0_A(s,Fs)}$. Applying
  ${\widehat{\operatorname{Ext}}{}^0_A(s,-)}$ to the triangle gives an exact
  sequence
  \[
    (H^{-1})^{\oplus2}\xrightarrow{(w_1[-1],w_2[-1])}U^0\longrightarrow V^0
    \longrightarrow(H^0)^{\oplus2}\xrightarrow{(w_1,w_2)}U^1.
  \]
  The first map vanishes by hypothesis, so ${U^0\cong\kk}$ injects into $V^0$.
  The last map has rank one, since ${H^0=\kk}$ and the $w_i$ are not both zero,
  so its kernel is one-dimensional. Hence, ${\dim V^0=2}$. Since ${H^0=\kk}$ acts by
  scalars, ${\delta^0(c,d)=(c-d)v}$ is a multiple of $v$, so the image of
  $\delta^0$ has dimension at most one and $\delta^0$ is not surjective. By the
  exact sequence~\eqref{eq:conversion-sequence} for ${a=1}$, the cokernel of
  $\delta^0$ embeds into $\operatorname{Ext}^1_\Lambda(Z,Z)$.
\end{proof}

\begin{corollary}
  \label{coro:one-factor}
  The hypothesis ${w_1[-1]\circ\beta_0=0=w_2[-1]\circ\beta_0}$ of
  \Cref{prop:one-factor} always holds. Hence, in the situation of
  \Cref{prop:one-cone}, the formula in \Cref{thm:conversion}, applied to a
  single cone as above, never produces a module $Z$ with
  ${\operatorname{Ext}^1_\Lambda(Z,Z)=0}$.
\end{corollary}

\begin{proof}
  Since ${\tau\neq0}$, the nondegenerate pairing of \Cref{thm:tate-duality},
  which is compatible with composition, makes the map ${H^{-4}\to H^{-1}}$,
  ${\gamma\mapsto\gamma[3]\circ\tau}$, nonzero, hence surjective onto the
  one-dimensional space $H^{-1}$. Choose $\gamma$ with
  ${\gamma[3]\circ\tau=\beta_0}$. For ${w\in U^1}$,
  \[
    w[-1]\circ\beta_0=(w[-4]\circ\gamma)[3]\circ\tau,
  \]
  and ${w[-4]\circ\gamma\in U^{-3}}$, which is zero by \Cref{prop:one-cone}.
\end{proof}

\subsection{A Tate-duality obstruction}
\label{subsec:tate-obstruction}

\begin{theorem}
  \label{thm:tate-obstruction}
  Let $A$ be a finite-dimensional symmetric algebra and $s$ a nonprojective
  simple $A$-module with ${\operatorname{End}_A(s)=\kk}$ and
  \[
    \operatorname{Ext}^1_A(s,s)=\operatorname{Ext}^2_A(s,s)=
    \operatorname{Ext}^4_A(s,s)=0.
  \]
  Then for all ${\tau\in\widehat{\operatorname{Ext}}{}^3_A(s,s)}$ and
  ${\beta\in\widehat{\operatorname{Ext}}{}^{-1}_A(s,s)}$ the Toda bracket
  $\langle\tau,\beta,\beta\rangle$ is defined and equal to~$\set{0}$.
\end{theorem}

\begin{proof}
  Write ${H^a}={\widehat{\operatorname{Ext}}{}^a_A(s,s)}$. By
  \Cref{thm:tate-duality}, ${H^{-1-a}\cong DH^a}$, so ${H^{-2}}={H^{-3}}={H^{-5}}={0}$,
  and ${H^{-1}\cong DH^0}$ is one-dimensional, since ${H^0}={
  \underline{\operatorname{End}}_A(s)}={\kk}$ for the nonprojective module $s$
  with ${\operatorname{End}_A(s)}={\kk}$. The bracket is defined, since
  ${\tau\beta\in H^2=0}$ and ${\beta^2\in H^{-2}=0}$; it lies in $H^0$, and its
  indeterminacy ${\tau H^{-3}+H^1\beta}$ is zero. If ${\tau=0}$, then $0$ is an
  element of the bracket, obtained with ${b=0}$ in the defining system. If
  ${\tau\neq0}$, the nondegenerate pairing of \Cref{thm:tate-duality}, which is
  compatible with composition, makes the map ${H^{-4}\to H^{-1}}$,
  ${\gamma}\mapsto{\gamma\tau}$, nonzero and hence surjective onto the
  one-dimensional space $H^{-1}$; choose ${\gamma\in H^{-4}}$ with
  ${\gamma\tau=\beta}$. The bracket $\langle\tau,\beta,\gamma\rangle$ is defined,
  since ${\beta\gamma}\in{ H^{-5}}={0}$, and it lies in ${H^{-3}=0}$. By
  \Cref{lemma:toda-juggling},
  \[
    \set{0}=\langle\tau,\beta,\gamma\rangle\circ\tau\subseteq
    \langle\tau,\beta,\gamma\tau\rangle=\langle\tau,\beta,\beta\rangle.
  \]
  Since the indeterminacy vanishes, ${\langle\tau,\beta,\beta\rangle=\set{0}}$.
\end{proof}

\begin{corollary}
  \label{coro:tate-obstruction}
  Let $A$ be a finite-dimensional symmetric algebra and $s$ a simple
  $A$-module. If ${\operatorname{End}_A(s)=\kk}$ and
  ${\operatorname{Ext}^*_A(s,s)=\kk[\tau]}$
  with $\tau$ of degree ${p\geq3}$, then ${\langle\tau,\beta,\beta\rangle=\set{0}}$
  for every ${\beta\in\widehat{\operatorname{Ext}}{}^{-1}_A(s,s)}$.
\end{corollary}

\begin{proof}
  For ${p=3}$ this is \Cref{thm:tate-obstruction}. For ${p>3}$ the bracket lies in
  ${\widehat{\operatorname{Ext}}{}^{p-3}_A(s,s)}={\operatorname{Ext}^{p-3}_A(s,s)}$,
  which vanishes since ${0<p-3<p}$.
\end{proof}

The tensor square ${S=s\otimes s}$ over ${E=T\otimes T}$ of \Cref{sec:ar-route}
also satisfies the hypotheses of \Cref{thm:tate-obstruction}, since
${\operatorname{Ext}^*_E(S,S)=\kk[\tau_1,\tau_2]}$ vanishes in degrees $1$, $2$
and~$4$; so the corresponding brackets vanish there as well. The construction
over ${T\otimes_\kk T}$ does not depend on such a bracket: it uses a second cone,
after which the only nonzero groups $W^a$ are in degrees $0$ and $3$, as in
\Cref{prop:two-cone-profile}, so that the term ${W^2=0}$ replaces the term $U^0$
of \Cref{prop:one-factor}. \Cref{thm:tate-obstruction} does not exclude
other constructions, for instance with simple modules whose stable
endomorphism ring is larger than $\kk$, with other brackets, or with a different
construction from stable morphisms to modules.

\begin{remark}
  \label{rem:weights}
  For trivial extensions ${T=C\ltimes DC}$, graded with $C$ in degree zero and
  $DC$ in degree one, there is a second explanation, by weights. The twists
  $h_\lambda$ of \Cref{subsec:cocycle} act on a homogeneous stable class of
  internal degree $\delta$ by $\lambda^{\delta}$, under the convention of
  \Cref{lemma:twist}; a bracket of homogeneous classes with zero indeterminacy
  and target concentrated in internal degree zero vanishes unless the
  internal degrees add up to zero. For the pair $(T,s)$ of
  \Cref{sec:ar-route}, let $\beta_0$ span
  $\widehat{\operatorname{Ext}}{}^{-1}_T(s,s)$. Then ${\delta(\tau)=-1}$ by
  \Cref{lemma:twist} and ${\delta(\beta_0)=1}$, since
  $\beta_0$ is dual to the identity under a form of internal degree one, so
  the weights do not cancel. This argument, and a six-dimensional example in
  which the weights cancel and a bracket of the form $\langle\tau,\beta,\beta\rangle$ is
  nonzero, are
  recorded with their verification in the repository of the project; see
  \Cref{app:verification}. The six-dimensional example has periodic, not
  polynomial, Ext algebra, in accordance with \Cref{thm:tate-obstruction}.
\end{remark}

\subsection{Computational evidence}
\label{subsec:obstruction-computations}

The following computations concern constructions with a single cone, or without
cones. Their scripts and parameters are described in \Cref{app:computations};
the results agree with \Cref{coro:one-factor} in the stated cases and degrees.
They are finite computations with the stated scope,
not proofs in all degrees.
\begin{enumerate}
  \item An example without cones:
    \[
      \Lambda_0=\begin{pmatrix}T&0\\T_{H_1}\oplus T_{H_2}&T\end{pmatrix},
    \]
    of dimension $80$ with four simple modules, and a nonprojective module $Z_0$
    of dimension ten. Over finite fields $\mathbb{F}_{2^{16}}$, in three runs
    and up to degree six, $\operatorname{Ext}^a_{\Lambda_0}(Z_0,Z_0)$ has
    dimensions $1,0,0,0,0,0$ for ${a=1,\dots,6}$ and
    ${\operatorname{Ext}^a_{\Lambda_0}(Z_0,\Lambda_0)=0}$; the same holds exactly
    over $\mathbb{F}_2(q,H_1,H_2)$ up to degree two. The surviving class is the
    cokernel of $\delta^0$, as predicted by~\eqref{eq:conversion-sequence}.
  \item The candidate constructed over $T$ from the cone of $\tau$, lifts of
    the twists and a fibre $F_1$, of dimension $352$, giving an algebra
    $\Lambda_1$ of dimension $392$ with four simple modules and a module $Z_1$
    of dimension $16$. Over $\mathbb{F}_{2^{16}}$, in two samples,
    $\operatorname{Ext}^a_{\Lambda_1}(Z_1,Z_1)$ has dimensions $1,0,0$ and
    ${\operatorname{Ext}^a_{\Lambda_1}(Z_1,\Lambda_1)=0}$ for ${a=1,2,3}$, and
    $\delta^0$ has rank one, as in \Cref{prop:one-factor}.
  \item Three further attempts to use a single cone over algebras of dimensions
    $6$, $12$ and $40$ fail because the relevant simple modules do not have
    polynomial Ext algebras on one generator of degree three; their Ext
    dimensions were computed over $\mathbb{F}_{2^{16}}$ up to degree twelve,
    using a tensor-product formula in degrees $2$ to $12$ for the
    forty-dimensional example.
\end{enumerate}


\appendix
\crefalias{section}{appendix}
\section{Computations}
\label{app:computations}

\emph{Certificates} check finitely many polynomial identities exactly,
over polynomial rings such as $\mathbb{F}_2[q]$ or $\mathbb{F}_2[q,\lambda]$;
they are part of the proofs of the statements that cite them, and
multilinearity extends them to all inputs. \emph{Experiments} compute
homological invariants in finitely many degrees, over finite fields or over
$\mathbb{F}_2(q,H_1,H_2)$; they are evidence with the stated scope and prove
nothing about other degrees or other fields. All scripts and their saved
outputs are in the directory \texttt{computations/} of the repository described
in \Cref{app:ai-declaration}. A rerun of the
scripts, recorded in the repository, executed 45 commands with Python~3.14 and
SageMath~10.10, and their mathematical output agreed with the saved results.
Three finite-field runs for $\Lambda_0$ were not repeated, and two stages for
$\Lambda_1$ and two stages of the attempt using minimal bimodule resolutions
over ${T\otimes_\kk T}$ were rerun from
saved intermediate results. The finite fields
$\mathbb{F}_{2^{16}}$ below are defined by the modulus
${x^{16}+x^5+x^3+x^2+1}$, and the parameters $q$, $H_1$ and $H_2$ are powers of a
generator, with exponents recorded in the saved outputs.

\subsection*{The cochain}

The following table, reproduced from the appendix of~\cite{OAI26ar}, defines the
cochain $p$ of \Cref{subsec:cocycle}. Lower-case letters denote the basis of $C$
of \Cref{subsec:algebra-C} and capital letters the dual basis of $DC$; an entry
$abcw$ in the row $q^d$ means that the coefficient of the basis vector $w$ in
$p(a,b,c)$ is $q^d$, and all other coefficients are zero. A script confirms that
the table below and the source table have the same $179$ entries
in the same order.

\begin{center}
\small
\begin{tabular}{r@{\quad}l}
 coefficient & input and output letters \\ \hline
 $q^2$ & \texttt{EuNT EyZX FNtU JtXZ NtEU NtXV XzJT ZxEY tEun tXuf} \\
 $q^2$ & \texttt{tXvn tXzj tuFn tuNf tvNn uFNT vNNT vtJu xEyz zJtx} \\
 $q$ & \texttt{EuUE EuVX EvVE ExXE EyYE EzZE FNnF FVvF FVxU FtTF} \\
 $q$ & \texttt{FtYU Ftun Ftxj JFtX JNtZ JUyX JVyZ JjEE JjTT JjXX} \\
 $q$ & \texttt{NFtU NNtV NUyU NVyV NjJN NjZV NnNN NnVV Ntuf Ntvn} \\
 $q$ & \texttt{Ntzj TtEE TtXX UuFF UuUU UvNF UvVU UxXU Uyun Uyxj} \\
 $q$ & \texttt{UzZU VvNN VvVV VxEU VxXV Vyuf Vyvn Vyzj VzZV XxEE} \\
 $q$ & \texttt{XxXX YuFT YuUY YvNT YvVY YxXY YyEE YyJJ YyTT YyXX} \\
 $q$ & \texttt{YyYY YyZZ YzZY ZxXZ ZzEE ZzJJ ZzTT ZzXX jJFF jJUU} \\
 $q$ & \texttt{jJjj jJnn jYyj jZvn jZzj nNFF nNUU ntTn tJjt tTNN} \\
 $q$ & \texttt{tTVV tTtt tYyt tZzt tzJn uFVX uFty uNFT uNUY uNte} \\
 $q$ & \texttt{uUyy uVye ujJu vNUE vNVX vNty vVyy vtTv xEuv xJFT} \\
 $q$ & \texttt{xJUY xJjx xJte xTtx xXuu xXvv xXyy xXzz xYyx xZye} \\
 $q$ & \texttt{xZzx xuFv xuNu xvNv yxJu yzJv zJNT zJVY zYyz zZuu} \\
 $q$ & \texttt{zZvv zZyy zZzz} \\
 $q^3$ & \texttt{EyJT JtEY tEyj tyJf} \\
 $1$ & \texttt{JNjX JVzX NNjU NVzU NjYU Njun Njxj Ntxt TNtX TVyX} \\
 $1$ & \texttt{Vyxt VzYU Vzun Vzxj YuNJ YuVZ jYun juNn jxJn txTn} \\
 $1$ & \texttt{uNNJ uNVZ uNjy uVzy xJNJ xJVZ xTNT xTVY yxTv zYuv} \\
 $1$ & \texttt{zuNv zxJv} \\
\end{tabular}
\end{center}

\subsection*{Certificates for \texorpdfstring{\Cref{sec:ar-route}}{Section 9}}

\begin{enumerate}
  \item \emph{The algebras $C$ and $T$}
  (\Cref{lemma:algebra-C,lemma:T-symmetric};
    \texttt{08-D-E/foundations\_}\allowbreak\texttt{certificate.py},
    over $\mathbb{F}_2[q]$). All
    $1000$ associators of basis vectors of $C$ and all $8000$ of $T$ vanish; the
    grading is respected by all products; the trace form of $T$ is symmetric,
    associative and has the stated dual bases; ${(ut)^2}={q(1+q)z}$. An earlier,
    independently written script (\texttt{02-ar-finite-data/}) checks the same
    identities over $\mathbb{F}_2(q)$ and $\mathbb{F}_{2^{16}}$ and computes the
    resolution of $s$ in degrees up to eight, in agreement with
    \Cref{prop:C-resolution}.
  \item \emph{The cochain} (\Cref{lemma:cocycle};
    \texttt{08-D-E/finite\_cochain\_certificate.py}, over $\mathbb{F}_2[q]$ with
    the coefficients of $\lambda^0$ and $\lambda^1$ compared separately, hence
    over $\mathbb{F}_2[q,\lambda]$). Every entry has radical inputs, matching
    corners and weight ${-1}$; the cocycle identity holds on all ${18^4=104\,976}$
    words of four radical basis vectors; the boundary identity holds on all
    $324$ pairs; ${p(t,x,j^*)=0}$, ${p(t,y,j^*)=q^3f}$, and the bar chain $\zeta$
    of \Cref{prop:cocycle-generator} is a cycle. The script
    \texttt{08-D-E/certificate\_receipt.py} confirms that its table equals the
    source table.
  \item \emph{Casimir identities} (\Cref{prop:lifts}; the same script). The sum
    \[
      \sum_wh_\lambda(w)w^*
    \]
    over the twenty basis vectors vanishes, and the sum
    over the radical basis vectors with left idempotent $r$ equals $r^*$, for
    ${r\in\set{e,f}}$. The other identities used in that proof are proved in the
    text.
\end{enumerate}

\subsection*{Dimensions}

The dimension of $C_1$ in \Cref{rem:ar-sizes} is computed exactly from the
corner dimensions of $T$ and $\operatorname{rad}T$, which give the dimensions
of the terms of the bar resolution, and from the cohomology $E$, $E^2$ and $E$
of $K$ in degrees $0$, $2$ and $4$ (\texttt{09-sizes/dimensions.py}). The
terms $K^0,\dots,K^4$ have dimensions $1\,761\,405\,952$, $148\,453\,376$,
$11\,713\,792$, $818\,688$ and $43\,264$, which gives ${\dim
C_1=1\,623\,889\,344}$
and
\[
  \begin{aligned}
    \dim_\kk\Lambda&=1600+159999^5\cdot1\,623\,889\,344\\
    &=170\,271\,818\,183\,326\,072\,615\,867\,045\,851\,312\,256.
  \end{aligned}
\]
The computation (\texttt{06-two-factor/}, over $\mathbb{F}_{2^{16}}$)
used minimal bimodule resolutions of the cones and determined, from the
dimensions of the terms and of the homology of the complex, cokernels of
dimensions $784\,704$ and $1\,377\,984$; it
stopped, before forming the bimodule matrices, at a linear system with
$3\,591\,168$ unknowns per right-hand side, beyond the bound of ${2\cdot10^6}$
set in advance. The same computation gives the dimensions of the Tate cohomology
groups in \Cref{prop:two-cone-profile} in degrees ${-4\leq a\leq7}$ for its sample
parameters.

\subsection*{Experiments for \texorpdfstring{\Cref{sec:obstructions}}{Section
10}}

\begin{enumerate}
  \item \emph{The algebra $\Lambda_0$} (\texttt{03-testbed-lambda0/}). Three
    runs over $\mathbb{F}_{2^{16}}$ in degrees up to six and one exact run over
    $\mathbb{F}_2(q,H_1,H_2)$ in degrees up to two compute minimal projective
    resolutions of $Z_0$ and the groups $\operatorname{Ext}^a(Z_0,Z_0)$ and
    $\operatorname{Ext}^a(Z_0,\Lambda_0)$. A separate script compares the
    surviving class with the cokernel of $\delta^0$; this map has rank one and
    kernel spanned by $(1,1)$.
  \item \emph{The algebra $\Lambda_1$ constructed over $T$}
  (\texttt{05-candidate1/}).
    Two runs over $\mathbb{F}_{2^{16}}$ construct minimal bimodule covers, the
    cocycle, the cone of dimension $312$, normalised lifts, the fibre of
    dimension $352$, the algebra $\Lambda_1$ and the module $Z_1$, and compute
    the dimensions of the Tate cohomology groups of the cone in degrees ${-4\leq
    a\leq4}$ (nonzero only in degrees
    $0$ and $1$, as in \Cref{prop:one-cone}), the matrix of $\delta^0$, and
    $\operatorname{Ext}^a$ for ${a=1,2,3}$.
  \item \emph{Further attempts using a single cone}
  (\texttt{07-one-factor-search/},
    over $\mathbb{F}_{2^{16}}$). For the six-dimensional trivial extension of
    ${\kk(1\to2)}$, its tensor product with $\kk[\epsilon]/(\epsilon^2)$, and the
    trivial extension of $T$, minimal resolutions of the simple modules give
    Ext dimensions up to degree twelve, directly for the first two and, for
    the third, directly in degrees $0$ and $1$ and by a tensor-product formula
    in degrees $2$ to $12$. None has polynomial Ext algebra on a generator of
    degree three. The same directory contains the six-dimensional example
    mentioned in \Cref{rem:weights}.
  \item \emph{Brackets for $T$} (\texttt{04-toda-bracket/}, over
    $\mathbb{F}_2(q)$ and over $\mathbb{F}_{2^n}$ for ${n}={8,12,16}$). An explicit
    defining system shows ${\langle\tau,\beta_0,\beta_0\rangle}={\set{0}}$, in
    agreement with \Cref{thm:tate-obstruction}.
\end{enumerate}

\subsection*{Checks of signs and finite cases}

The remaining scripts check signs and finite cases of general arguments; they
are not needed for the proofs. They cover the finite quotients of $G$ for
${m\leq10}$
(\texttt{D-C-selection-checks.py}); the signs of the bar decomposition, of the
simulation and two examples, including
\Cref{ex:derived-powers} (\texttt{D-A-checks.py}); the rectification for an
algebra with a nonzero null-homotopic relation
(\texttt{D-B-rectification-check.py});
and the signs of the cone in the proof of \Cref{thm:conversion}
(\texttt{D-D-cone-signs.py}).

\section{Verification and formalisation record}
\label{app:verification}

All checks
were carried out by AI models; none is a human verification, and the
formalised statements are the only results checked by a proof assistant. The
complete reports are in the repository described in \Cref{app:ai-declaration}.

\subsection*{Proof notes and fresh-context verification}

Before any section was written, every result received a written proof note,
either compiled from the first part of the work or drafted from scratch, and
the note was then checked by a fresh session of GPT-6 Astra that had not seen
the note being written, was given the statements and proofs with the earlier
verdicts and statuses removed as far as possible, and was instructed not to use
any remaining assessment as evidence. Each drafted section was checked
again, by another fresh session, against its note and its sources.
\Cref{tab:verification} lists, for each part of the article, the author of the
proof note, the verification of the note and the verification of the section;
the effort settings are those of the Codex command-line interface, where
\emph{ultra} is the highest.

\begin{table}[htbp]
  \centering
  \small
  \caption{Authors of the proof notes and verifications, by part of the
    article.}
  \label{tab:verification}
  \begin{tabular}{@{}>{\raggedright\arraybackslash}p{0.16\textwidth}%
    >{\raggedright\arraybackslash}p{0.21\textwidth}%
    >{\raggedright\arraybackslash}p{0.26\textwidth}%
    >{\raggedright\arraybackslash}p{0.26\textwidth}@{}}
  \hline
  Part&Proof note&Verification of the note&Verification of the section\\
  \hline
  \Cref{sec:criteria}&Claude Opus~5.5&GPT-6 Astra (max): errors and gaps in
    five blocks, repaired&GPT-6 Astra (ultra): two local slips and an overstated
    sentence repaired\\
  \Cref{sec:extinction,sec:simulation}&GPT-6 Astra (max)&GPT-6 Astra (ultra):
    no error; two wording repairs&GPT-6 Astra (ultra): wording repairs\\
  \Cref{sec:selection}&GPT-6 Astra (max); Claude Opus~5.5 for
    \Cref{subsec:selection-data}&GPT-6 Astra (max): no error&GPT-6 Astra
    (ultra): one range and wording repairs\\
  \Cref{sec:realisation}&GPT-6 Astra (ultra)&GPT-6 Astra (ultra): no error;
    one wording repair&GPT-6 Astra (ultra): wording and scope repairs\\
  \Cref{sec:conversion}&GPT-6 Astra (ultra)&GPT-6 Astra (ultra): no error&
    GPT-6 Astra (ultra): no error; one remark narrowed\\
  \Cref{sec:ar-route}&GPT-6 Astra (ultra)&GPT-6 Astra (ultra): no error; one
  missing input supplied, one statement narrowed&GPT-6 Astra (ultra): two
    steps restored from the note, one definition completed, wording repairs\\
  \Cref{sec:obstructions}&Claude Opus~5.5, from proofs by GPT-6 Astra
    (medium to max)&GPT-6 Astra (max): theorems confirmed, two readings
    corrected; the search computations not rechecked&GPT-6 Astra (ultra): a gap
    in the argument for a single cone closed by a direct proof
    (\Cref{coro:one-factor})\\
  \hline
  \end{tabular}
\end{table}

The errors found and repaired in this way include an invalid inference from
dimensions alone in an earlier version of \Cref{prop:infinitely-torsionfree},
an unstated convention on left and right modules in \Cref{prop:simple-witness},
a wrong lower
bound on the size of a construction of a counterexample to the Auslander--Reiten
conjecture using minimal bimodule resolutions, a misreading
of the scope of \Cref{thm:tate-obstruction}, and several statements whose
wording claimed more than their proofs. Two broader claims of the first part
of the work were withdrawn.

\subsection*{Formalisation}

The following results are formally verified in Lean~4 (version 4.33.1)
against Mathlib (commit \texttt{0df444a3}), without additional axioms beyond
\texttt{propext}, \texttt{Classical.choice} and \texttt{Quot.sound}. Items 1--5
are unconditional. Item~6 is conditional on explicit hypotheses, listed
below.
\begin{enumerate}
  \item \Cref{prop:coresolution}, for an arbitrary abelian category with enough
    projective objects, stated as the conjunction of the projective dimension
    bound and its sharpness.
  \item \Cref{thm:strong-nakayama} over an arbitrary ring, in two stages. For
    a resolution of a nonzero right module $E$ by finitely generated projective
    modules, the vanishing of the groups $\operatorname{Ext}^i_{A^{\op}}(E,A)$
    of Mathlib for all ${i\geq0}$ is equivalent to the exactness of the dual
    complex, and the cokernels $C_n$ of the dual complex are then finitely
    generated left modules of projective dimension exactly~$n$, for every
    ${n\geq1}$; moreover $C_n$ is the transpose of the presentation
    ${P_n\to P_{n-1}}$ of $\Omega^{n-1}E$. Minimality of this presentation, which
    the report includes in the definition of the transpose, is not formalised,
    since Mathlib has no projective covers.
  \item The group of \Cref{sec:selection}: \Cref{def:selection-group} as a
    presented group with exactly the stated relations, its finite presentation
    (\Cref{prop:group-presentation}, with Mathlib's notion of a finitely
    presented group), \Cref{prop:central-involutions},
    \Cref{prop:shift-automorphism}, \Cref{prop:finite-quotients} for every
    ${m\geq1}$, and \Cref{coro:z-independent}.
  \item The argument of \Cref{prop:rank-obstruction}, as a statement of linear
    algebra: if a linear form on a finite-dimensional vector space $\mathcal{V}$
    vanishes on $\mathsf{T}^tv$ for all large $t$, then it vanishes on
    $\mathsf{T}^tv$ for ${t\geq\dim\mathcal{V}}$; this holds over any field and
    without assuming that the vectors $\mathsf{T}^tv$ span $\mathcal{V}$. The
    passage from Grothendieck groups and rank functions to this statement is
    not formalised.
  \item \Cref{lemma:fractions}, for an arbitrary localisation with a calculus
    of right fractions; \Cref{lemma:cone-splitting}, in Mathlib's idempotent
    completion of a pretriangulated category; and \Cref{prop:odd-double}
    without the assertion about $K_0$, with the objects $C$ and $V$ and their
    decompositions constructed from the two cones of the proof.
  \item The results of \Cref{sec:obstructions}, in a $\kk$-linear
    pretriangulated category with finite-dimensional morphism spaces, under
    two further hypotheses: a perfect pairing between the spaces of morphisms
    ${X\to Y[a]}$ and ${Y\to X[-1-a]}$, natural and compatible with composition
    as in \Cref{thm:tate-duality} and the paragraph after it; and, for
    \Cref{prop:one-cone,prop:one-factor,coro:one-factor}, an object $s$ whose
    morphisms ${s\to s[n]}$, ${n\geq0}$, form the graded algebra $\kk[\tau]$ with
    $\tau$ of degree~$3$, as in \Cref{thm:T-ext}. Toda brackets are defined
    from triangles as in \Cref{subsec:toda}, and their properties used in the
    proofs, including \Cref{lemma:toda-juggling}, are proved. The formal
    statements cover \Cref{prop:one-cone}, \Cref{thm:tate-obstruction},
    \Cref{coro:tate-obstruction}, and the parts of \Cref{prop:one-factor} and
    \Cref{coro:one-factor} that concern the triangle: the vanishing of the
    composites, the dimension of $\widehat{\operatorname{Ext}}{}^0_A(s,Fs)$
    and the failure of surjectivity of ${(c,d)\mapsto(c-d)v}$. The assertions
    about $\operatorname{Ext}^1_\Lambda(Z,Z)$ and the identification of this
    map with $\delta^0$ require \Cref{thm:conversion} and are not formalised.
\end{enumerate}
The formalisation of \Cref{prop:coresolution}, statements and proof, was
written by Claude Opus~5.5. All other items were implemented by GPT-6 Astra
from specifications by Claude Opus~5.5, who then reviewed their formal
statements and, for item~6, the formal hypotheses. Acceptance required
five checks: a scan of the
sources for placeholders and new axioms, a build with warnings treated as
errors, an axiom report, a listing of the formal statements, and an
independent replay of the proof terms by the Lean kernel checker. The
formalisation is included in the repository described in
\Cref{app:ai-declaration} (directory \texttt{lean/formalisation}), with its check
logs (\texttt{lean/gates}).
A search of Mathlib shows that a complete formalisation of either construction
would require theories not currently available there, in particular derived
tensor products over
noncommutative algebras and dg resolutions for the main construction, and
stable categories of symmetric algebras with Tate duality for the construction
of the counterexample to the Auslander--Reiten conjecture. Of the smaller parts
identified by the feasibility review, only the finite computations of
\Cref{sec:ar-route} have not been formalised. The evaluation is recorded in
the repository.

\section{The preprint on Tachikawa's second conjecture}
\label{app:companion}

The two preprints discussed in this article were released together with a
third, \emph{A counterexample to Tachikawa's second
conjecture}~\cite{OAI26tachikawa}. Its main
theorem~\cite[Theorem~1.1]{OAI26tachikawa}
asserts that over
${\kk=\mathbb{F}_2(q,H_1,H_2)}$ there are a finite-dimensional symmetric algebra
$A$ and a nonprojective module $M$ with ${\operatorname{Ext}^i_A(M,M)=0}$ for all
${i>0}$. Since $A$ is self-injective, $M$ is then also a counterexample to the
Auslander--Reiten conjecture, and its consequences for the endomorphism
algebra ${\Gamma=\operatorname{End}_A(A\oplus M)^{\op}}$ include infinite
dominant dimension and infinite little left finitistic dimension. We have not
verified that preprint, and nothing in this article depends on it.

Given $A$ and $M$,
\Cref{thm:ar-to-findim} applies directly and gives infinite little finitistic
dimension of ${\operatorname{End}_A(A\oplus M)}$, with one more simple module than
$A$ if $A$ is basic and $M$ indecomposable. The preprint also states the
existence~\cite[Corollary~1.2(2)]{OAI26tachikawa} of a simple
$\Gamma$-module $S$ with ${\operatorname{Ext}^i_\Gamma(S,\Gamma)=0}$ for all
${i\geq0}$; \Cref{thm:ar-to-findim} establishes this property for a module that
need not be simple, and \Cref{thm:strong-nakayama}, applied
to $S$ regarded as a right module over $\Gamma^{\op}$, gives explicit modules
of every finite projective dimension over $\Gamma^{\op}$. The preprint's proof
of
infinite left finitistic dimension of $\Gamma$ proceeds through dominant
dimension instead. In \Cref{sec:ar-route}, a bimodule and a morphism in the
stable category of a symmetric algebra instead give a counterexample to the
Auslander--Reiten conjecture over a triangular algebra, without requiring a
module with vanishing self-extensions over the symmetric algebra itself.

\section{Declaration of the use of AI and outline of the process}
\label{app:ai-declaration}

This article was prepared autonomously by AI agents, as instructed by Gustavo
Jasso, between 7 and 8 October 2026. All its text, proofs, computations and
formal proofs were produced by AI models; no part of it was written by a human.
Gustavo Jasso designed the task and took the decisions recorded below. No human
has checked the mathematics of this article, and it is not peer reviewed.

\subsection*{Models}

The models and their roles were as follows; the names are those of the
project's records.
\begin{itemize}
  \item Claude Opus~5.5 (Anthropic), run in Claude Code: coordination of the
  work, the analysis of the preprints, the general criteria and obstructions
    of \Cref{sec:criteria} and \Cref{sec:obstructions} (part of them jointly
    with the other models, as recorded in the repository), the proof notes
    compiled from earlier results, the checking of verification reports, the
    first stage of the Lean formalisation, and the first draft of the prose of
    this article.
  \item GPT-6 Astra (OpenAI), run in Codex (command-line interface, version
    0.161), with reasoning effort between medium and its highest setting:
    independent verification of every proof in fresh sessions, all heavy
    computations, the proof notes for the constructions of the two preprints,
    the Tate-duality obstruction of \Cref{thm:tate-obstruction}, the later
    stages of the Lean formalisation, and the block-by-block review of the
    article.
  \item Claude Fable~5.1 (Anthropic): two consultations on possible simpler
    constructions and on the obstructions to them; the second was incomplete.
  \item Claude Sonnet (Anthropic), as subagents: retrieval of literature and
    bibliographic data, summaries of instructions, and one computation; their
    output was checked before use. In the second checks described in
    \Cref{app:verification}, Claude Sonnet~5.5 coordinated the work and
    drafted several edits, which were reviewed by Claude Opus~5.5.
\end{itemize}
A ChatGPT Pro~100 subscription and a premium seat in the Claude Team plan for
scientists (at a discounted rate) were available for the work.
% Usage figures: USAGE.md in the repository.

\subsection*{Process}

The first part (7--8 October 2026) analysed the main
preprint~\cite{OAI26findim} and searched for a simpler counterexample to the
little finitistic dimension conjecture; it produced the criteria of
\Cref{sec:criteria} and the obstructions of \Cref{sec:obstructions}, but found
no counterexample verifiable by hand. The second part
produced this article from that record.

Every mathematical claim carried a status (heuristic, plausible, supported by
computation, proved by an AI model, or verified by a second, independent check
of a different kind) and every change to a stronger status was accompanied by a
record of the supporting evidence. Proofs were checked step by step in fresh
sessions, sometimes of
the same model that wrote them. For the checks of the proof notes, earlier
verdicts were removed as far as possible, and the reviewers were instructed not
to use any remaining assessment as evidence; several errors and gaps were found
and
corrected in this way, and two broader claims were withdrawn. Computations used
exact arithmetic over finite fields and rational function fields, and their
scope is stated wherever they are used (\Cref{app:computations}). Citations were
checked against the sources, whose statements and locators were read before
being cited. Parts of \Cref{sec:criteria,sec:selection,sec:realisation} are
formally verified in Lean against Mathlib, and the arguments of
\Cref{sec:obstructions} are formally verified under explicitly stated
hypotheses (\Cref{app:verification}). The completed article was reviewed block by
block by fresh sessions of GPT-6 Astra without waiting for human approval of
each block, in three parts. The reviewers checked
$337$ blocks, corrected punctuation, typesetting and a few slips directly, and
proposed $15$ further changes, concerning notation and the precision of
statements about scope, attribution and evidence; each proposal was examined
by Claude Opus~5.5, and all were adopted. No error in a proof was found at this
stage.

\subsection*{Repository}

The repository
\url{https://github.com/gjassoah/findim-report} contains the complete record of
the work: the plan, the
notes, the verification reports, the scripts and their outputs, the Lean
sources and the evidence of their checks, the sources of this article, and a
log of all interactions between Gustavo Jasso and the agents, in which his
messages are reproduced verbatim and the agents' replies are summarised.


\printbibliography

\end{document}
